\Needspace{8\baselineskip}\originalchapter{2005 历史题：单粒子、波与输运}
本章取 Ron Parker 的 22.611J Fall 2005 历史题卷中九道与已选2003／2006题条件不同的公开题，均标“编者独立解答”；该批题卷未取得对应官方答案。来源PDF与完整去重判定见末章。相同原题在章末按原编号指回已选解答，不重复计题。下文$T$用于能量时已吸收$k_B$。

\originalsection{冷球与远隔双线圈镜}
\begin{exercise}\label{ex:hist05-1-3}
2005 PS1 原题3。半径$a$的均匀冷等离子体球中，电子、单价离子原先均以$n_0$占据同一球；将电子球作为整体相对固定的离子球微移$\boldsymbol\delta$，求小振动频率，并与无限长圆柱的同类位移比较。
\end{exercise}
\begin{solution}
编者独立解答。均匀带电球内场是$\mathbf E=\rho\mathbf r/(3\epsilon_0)$。两球重叠区的净场为$en_0\boldsymbol\delta/(3\epsilon_0)$，电子受力方向相反，因此$m_e\ddot{\boldsymbol\delta}=-e^2n_0\boldsymbol\delta/(3\epsilon_0)$，即
\begin{equation}\omega_{\rm ball}=\omega_{pe}/\sqrt3,\qquad \omega_{pe}^2=n_0e^2/(\epsilon_0m_e).\end{equation}
条件是$|\delta|\ll a$，边缘产生的高阶形变忽略。球的退极化系数$1/3$与作业\ref{ex:06-1-4}无限圆柱的$1/2$不同，所以两原题不能去重。
\end{solution}

\begin{exercise}\label{ex:hist05-2-1}
2005 PS2 原题1。两个半径$a$、同向电流$I$的圆线圈轴向相距$d\gg a$。一粒子在中间沿轴磁镜俘获，轴向振幅远小于$d$，Larmor半径$r_L$小于场变化尺度。(i)以各线圈远场偶极$B=\mu_0m/(2\pi r^3),m=I\pi a^2$求轴向振荡周期；(ii)用线圈中心场$B_c=\mu_0I/(2a)$估计中点可俘获的最小$v_\perp/v_\parallel$。
\end{exercise}
\begin{solution}
编者独立解答。令中心$z=0$，两线圈在$z=\pm d/2$。每个偶极系数$C=\mu_0Ia^2/2$，故
\begin{equation}
B(z)=C\bigl[(d/2-z)^{-3}+(d/2+z)^{-3}\bigr]
=B_m\left(1+24z^2/d^2+O(z^4/d^4)\right),\quad B_m=8\mu_0Ia^2/d^3.
\end{equation}
绝热磁矩$\mu=W_{\perp m}/B_m$守恒，沿轴势能$\mu B(z)$的二次项给$\omega_b^2=48W_{\perp m}/(m_pd^2)$，这里$m_p$表示被俘获粒子质量而非必为质子，周期$P_b=2\pi\sqrt{m_pd^2/(48W_{\perp m})}$。镜点须在到达线圈高场前出现：$W_{\parallel m}/W_{\perp m}\le B_c/B_m-1$，于是
\begin{equation}\left(\frac{v_\perp}{v_\parallel}\right)_{\min}\simeq\left[\frac{d^3}{16a^3}-1\right]^{-1/2}.\end{equation}
第二式把近线圈真实峰值仅以$B_c$估计；它与第一式的远场近似作用区域不同，不能在$r\sim a$继续使用偶极式。需$d^3>16a^3$且粒子磁矩保持。
\end{solution}

\begin{exercise}\label{ex:hist05-3-1}
2005 PS3 原题1。质量很大的中性原子近似静止，带电粒子$q,m$以速度$v_0$和碰撞参数$b$入射。题面印作原子诱导偶极$\mathbf p=-\alpha\mathbf E$，并给偶极场$\mathbf E_d=[3\hat{\mathbf r}(\mathbf p\cdot\hat{\mathbf r})-\mathbf p]/(4\pi\epsilon_0r^3)$。(i)写角动量和能量守恒；(ii)给出散射偏角积分；(iii)化为椭圆积分并取大$b$极限，与Coulomb散射比较。续页只另给积分表提示，不是新题。
\end{exercise}
\begin{solution}
编者独立解答，先照印刷负号计算。原子位于原点，正电荷在$\mathbf r$处时它在原子处的场是$\mathbf E_q=-q\hat{\mathbf r}/(4\pi\epsilon_0r^2)$，故印刷约定的$\mathbf p=-\alpha\mathbf E_q$与$\hat{\mathbf r}$同向。再由所给$\mathbf E_d$算出电荷受径向斥力$F_r=2\alpha q^2/[(4\pi\epsilon_0)^2r^5]$。记$C=\alpha q^2/[2(4\pi\epsilon_0)^2]$，则$U=+C/r^4$。角动量$L=mv_0b=mr^2\dot\phi$、能量
\begin{equation}\frac12mv_0^2=\frac12m\dot r^2+\frac{L^2}{2mr^2}+\frac{C}{r^4}.\end{equation}
设$u=b/r$、$\eta=2C/(mv_0^2b^4)$，正根$u_*^2=(\sqrt{1+4\eta}-1)/(2\eta)$，并设$b_*^2=(\sqrt{1+4\eta}+1)/(2\eta)$。散射偏角及完全椭圆积分$K$为
\begin{equation}
\chi=\pi-2\int_0^{u_*}\frac{\dd u}{\sqrt{1-u^2-\eta u^4}}
=\pi-2u_*K\!\left(-\frac{u_*^2}{b_*^2}\right),\quad
K(m)=\int_0^{\pi/2}\frac{\dd\varphi}{\sqrt{1-m\sin^2\varphi}}.
\end{equation}
这里$1-u^2-\eta u^4=(1-u^2/u_*^2)(1+u^2/b_*^2)$。大$b$即$\eta\ll1$，展开得$\chi=3\pi\eta/4+O(\eta^2)=3\pi C/(2mv_0^2b^4)+O(b^{-8})$；Rutherford小角偏转随$b^{-1}$而非$b^{-4}$。若$\alpha>0$被理解为通常物理极化率，惯例应是$\mathbf p=+\alpha\mathbf E_q$，相互作用改成$U=-C/r^4$，偏角一阶反号且高能存在俘获临界。不能悄悄把原题的负号当成吸引势。
\end{solution}

\originalsection{斜入射反射与碰撞 Alfv\'en 波}
\begin{exercise}\label{ex:hist05-7-2}
2005 PS7 原题2。真空平面波入射到介电常数$\epsilon=\epsilon_0(1-\omega_p^2/\omega^2)$的等离子体平面界面，电场垂直入射面。(a)由切向$E,H$连续求正常入射的复振幅反射系数$\Gamma=E_{\rm ref}/E_{\rm inc}$，指出完全反射频段，并以峰值电离层$n_e=10^{12}\,\mathrm m^{-3}$算频率；(b)对与法线夹角$\theta$的斜入射重算，数值取$45^\circ$。原题图仅标波矢、场方向和界面；两侧均视非磁性、无碰撞、半无限均匀介质。
\end{exercise}
\begin{solution}
编者独立解答。设$n_r=\sqrt{1-\omega_p^2/\omega^2}$并取使透射场向介质内衰减的平方根支。法向波数连续性与切向$E,H$匹配给TE系数
\begin{equation}
\Gamma_{\rm TE}(\theta)=\frac{\cos\theta-\sqrt{n_r^2-\sin^2\theta}}{\cos\theta+\sqrt{n_r^2-\sin^2\theta}}.
\end{equation}
正入射令$\theta=0$，得$(1-n_r)/(1+n_r)$。当$\omega<\omega_p/\cos\theta$时透射法向波数纯虚，无耗散半空间$|\Gamma|=1$；正入射阈值是$f_p=(2\pi)^{-1}\sqrt{n_ee^2/(\epsilon_0m_e)}\simeq8.98\,\mathrm{MHz}$，$45^\circ$阈值约$12.7\,\mathrm{MHz}$。等号是掠行阈值，等离子体有限厚度时另一界面可发生隧穿，不能沿用半空间的“全部反射”。本题斜入射条件未被2006 PS7的Airy密度坡题覆盖。
\end{solution}

\begin{exercise}\label{ex:hist05-8-1}
2005 PS8 原题1。均匀、静止、单价、$\mathbf B_0=B_0\hat z$的冷电子—离子流体扰动按$\exp(\ii\mathbf k\cdot\mathbf r-\ii\omega t)$，线性动量为
\begin{equation}
-\ii\omega m_s\mathbf v_s=q_s(\mathbf E+\mathbf v_s\times\mathbf B_0)-m_s\nu_{ss'}(\mathbf v_s-\mathbf v_{s'}),\qquad m_i\nu_{ie}=m_e\nu_{ei}.
\end{equation}
(a)将碰撞吸收到复质量$m_s^*=m_s(1+\ii\nu_{ss'}/\omega)$、复回旋频率$\Omega_s^*=m_s\Omega_s/m_s^*$，解$x,y$联立式；(b)在$|\omega|\ll|\Omega_{i,e}|$、一阶碰撞频率的原题排序下简化离子、电子响应并求$K_\perp,K_\times$；(c)由此求平行传播剪切Alfv\'en波最低阶阻尼。
\end{exercise}
\begin{solution}
编者独立解答。将另一种粒子速度视为此步的给定驱动，令$\mathbf F_s=(q_s/m_s)\mathbf E+\nu_{ss'}\mathbf v_{s'}$，横向线性系统的逆矩阵为
\begin{equation}
\begin{pmatrix}v_{sx}\\v_{sy}\end{pmatrix}
=\frac1{(\nu_{ss'}-\ii\omega)^2+\Omega_s^2}
\begin{pmatrix}\nu_{ss'}-\ii\omega&\Omega_s\\-\Omega_s&\nu_{ss'}-\ii\omega\end{pmatrix}
\begin{pmatrix}F_{sx}\\F_{sy}\end{pmatrix}.
\end{equation}
这同时给原题(a)的四个$x,y$式，代入两物种即可联立求解；复质量写法是同一逆矩阵。按原件保留到$O(\nu)$且先满足$|\omega|\ll|\Omega_s|$，两个物种共同的$\mathbf E\times\mathbf B$漂移在相对电流中抵消，离子极化漂移及摩擦留下
\begin{equation}
K_\perp\simeq\frac{\omega_{pi}^2}{\Omega_i^2}\left(1+\ii\gamma_c\right),\quad K_\times\simeq0,\quad
\gamma_c=\frac{2\omega\nu_{ie}}{\Omega_i^2}.
\end{equation}
这里取原题的低频、低碰撞排序，忽略电子惯性高阶项。剪切支$c^2k^2/\omega^2\simeq(c^2/v_A^2)(1+\ii\gamma_c)$，$v_A=B_0/\sqrt{\mu_0n_im_i}$，固定实$k$并令$\omega=\omega_0+\ii\gamma_t$得到
\begin{equation}\omega_0=kv_A,\qquad \gamma_t\simeq-\frac{\omega_0^2\nu_{ie}}{\Omega_i^2}.\end{equation}
$\gamma_c$是介电张量的小虚系数，$\gamma_t$才是有量纲阻尼率；不能混名。接近回旋共振或碰撞并非小扰动时此展开无效。
\end{solution}

\originalsection{旧参数冷波图与 SOL 温度}
\begin{exercise}\label{ex:hist05-8-2}
2005 PS8 原题2。氘等离子体Alcator C-Mod的两组参数分别为$(n_e,B_0)=(1\times10^{20}\,\mathrm m^{-3},5\,\mathrm T)$和$(4\times10^{20}\,\mathrm m^{-3},5\,\mathrm T)$。在$-2<\log_{10}(\omega/\Omega_{ci})<4$、$-2<\log_{10}n_{\rm ref}^2<3$窗口准确画传播区与截止区。原题未指定传播角；2006 PS8 题1改了两组密度和磁场，不能以其图冒充本题。
\end{exercise}
\begin{solution}
编者独立解答。取完全电离氘、$n_i=n_e$、$m_i\simeq2m_p$，在$\mathbf B_0\parallel\hat z$设带符号$\Omega_e=-eB_0/m_e,\Omega_i=eB_0/m_i$。冷张量的
\begin{equation}
R=1-\sum_s\frac{\omega_{ps}^2}{\omega(\omega+\Omega_s)},\quad
L=1-\sum_s\frac{\omega_{ps}^2}{\omega(\omega-\Omega_s)},\quad
P=1-\sum_s\frac{\omega_{ps}^2}{\omega^2},\quad S=(R+L)/2.
\end{equation}
平行传播画$R,L$，垂直传播画$O=P,X=RL/S$；$n_{\rm ref}^2>0$的区间为传播，零点是截止，$S=0$等为共振，不跨极点连线。两密度各一幅图由本卷\texttt{tools/historical\_2005\_cold.py}用真实旧参数重新采样，点集仅作图，零点和极点仍由上式判定。
\begin{center}\begin{tikzpicture}\begin{axis}[width=.77\textwidth,height=4.8cm,xlabel={$\log_{10}(\omega/\Omega_{ci})$},ylabel={$\log_{10}n_{\rm ref}^2$},xmin=-2,xmax=4,ymin=-2,ymax=3,unbounded coords=jump,title={$n_e=1\times10^{20}\,\mathrm m^{-3},\ B=5\,\mathrm T$},legend style={font=\small,at={(1.02,1)},anchor=north west}]
\addplot[thick,structurecolor] coordinates {(-2.00000,nan) (-1.99500,nan) (-1.99000,nan) (-1.98500,nan) (-1.98000,nan) (-1.97500,nan) (-1.97000,nan) (-1.96500,nan) (-1.96000,nan) (-1.95500,nan) (-1.95000,nan) (-1.94500,nan) (-1.94000,nan) (-1.93500,nan) (-1.93000,nan) (-1.92500,nan) (-1.92000,nan) (-1.91500,nan) (-1.91000,nan) (-1.90500,nan) (-1.90000,nan) (-1.89500,nan) (-1.89000,nan) (-1.88500,nan) (-1.88000,nan) (-1.87500,nan) (-1.87000,nan) (-1.86500,nan) (-1.86000,nan) (-1.85500,nan) (-1.85000,nan) (-1.84500,nan) (-1.84000,nan) (-1.83500,nan) (-1.83000,nan) (-1.82500,nan) (-1.82000,nan) (-1.81500,nan) (-1.81000,nan) (-1.80500,nan) (-1.80000,nan) (-1.79500,nan) (-1.79000,nan) (-1.78500,nan) (-1.78000,nan) (-1.77500,nan) (-1.77000,nan) (-1.76500,nan) (-1.76000,nan) (-1.75500,nan) (-1.75000,nan) (-1.74500,nan) (-1.74000,nan) (-1.73500,nan) (-1.73000,nan) (-1.72500,nan) (-1.72000,nan) (-1.71500,nan) (-1.71000,nan) (-1.70500,nan) (-1.70000,nan) (-1.69500,nan) (-1.69000,nan) (-1.68500,nan) (-1.68000,nan) (-1.67500,nan) (-1.67000,nan) (-1.66500,nan) (-1.66000,nan) (-1.65500,nan) (-1.65000,nan) (-1.64500,nan) (-1.64000,nan) (-1.63500,nan) (-1.63000,nan) (-1.62500,nan) (-1.62000,nan) (-1.61500,nan) (-1.61000,nan) (-1.60500,nan) (-1.60000,nan) (-1.59500,nan) (-1.59000,nan) (-1.58500,nan) (-1.58000,nan) (-1.57500,nan) (-1.57000,nan) (-1.56500,nan) (-1.56000,nan) (-1.55500,nan) (-1.55000,nan) (-1.54500,nan) (-1.54000,nan) (-1.53500,nan) (-1.53000,nan) (-1.52500,nan) (-1.52000,nan) (-1.51500,nan) (-1.51000,nan) (-1.50500,nan) (-1.50000,nan) (-1.49500,nan) (-1.49000,nan) (-1.48500,nan) (-1.48000,nan) (-1.47500,nan) (-1.47000,nan) (-1.46500,nan) (-1.46000,nan) (-1.45500,nan) (-1.45000,nan) (-1.44500,nan) (-1.44000,nan) (-1.43500,nan) (-1.43000,nan) (-1.42500,nan) (-1.42000,nan) (-1.41500,nan) (-1.41000,nan) (-1.40500,nan) (-1.40000,nan) (-1.39500,nan) (-1.39000,nan) (-1.38500,nan) (-1.38000,nan) (-1.37500,nan) (-1.37000,nan) (-1.36500,nan) (-1.36000,nan) (-1.35500,nan) (-1.35000,nan) (-1.34500,nan) (-1.34000,nan) (-1.33500,nan) (-1.33000,nan) (-1.32500,nan) (-1.32000,nan) (-1.31500,nan) (-1.31000,nan) (-1.30500,nan) (-1.30000,nan) (-1.29500,nan) (-1.29000,nan) (-1.28500,nan) (-1.28000,nan) (-1.27500,nan) (-1.27000,nan) (-1.26500,nan) (-1.26000,nan) (-1.25500,nan) (-1.25000,nan) (-1.24500,nan) (-1.24000,nan) (-1.23500,nan) (-1.23000,nan) (-1.22500,nan) (-1.22000,nan) (-1.21500,nan) (-1.21000,nan) (-1.20500,nan) (-1.20000,nan) (-1.19500,nan) (-1.19000,nan) (-1.18500,nan) (-1.18000,nan) (-1.17500,nan) (-1.17000,nan) (-1.16500,nan) (-1.16000,nan) (-1.15500,nan) (-1.15000,nan) (-1.14500,nan) (-1.14000,nan) (-1.13500,nan) (-1.13000,nan) (-1.12500,nan) (-1.12000,nan) (-1.11500,nan) (-1.11000,nan) (-1.10500,nan) (-1.10000,nan) (-1.09500,nan) (-1.09000,nan) (-1.08500,nan) (-1.08000,nan) (-1.07500,nan) (-1.07000,nan) (-1.06500,nan) (-1.06000,nan) (-1.05500,nan) (-1.05000,nan) (-1.04500,nan) (-1.04000,nan) (-1.03500,nan) (-1.03000,nan) (-1.02500,nan) (-1.02000,nan) (-1.01500,nan) (-1.01000,nan) (-1.00500,nan) (-1.00000,nan) (-0.99500,nan) (-0.99000,nan) (-0.98500,nan) (-0.98000,nan) (-0.97500,nan) (-0.97000,nan) (-0.96500,nan) (-0.96000,nan) (-0.95500,nan) (-0.95000,nan) (-0.94500,nan) (-0.94000,nan) (-0.93500,nan) (-0.93000,nan) (-0.92500,nan) (-0.92000,nan) (-0.91500,nan) (-0.91000,nan) (-0.90500,nan) (-0.90000,nan) (-0.89500,nan) (-0.89000,nan) (-0.88500,nan) (-0.88000,nan) (-0.87500,nan) (-0.87000,nan) (-0.86500,nan) (-0.86000,nan) (-0.85500,nan) (-0.85000,nan) (-0.84500,nan) (-0.84000,nan) (-0.83500,nan) (-0.83000,nan) (-0.82500,nan) (-0.82000,nan) (-0.81500,nan) (-0.81000,nan) (-0.80500,nan) (-0.80000,nan) (-0.79500,nan) (-0.79000,nan) (-0.78500,nan) (-0.78000,nan) (-0.77500,nan) (-0.77000,nan) (-0.76500,nan) (-0.76000,nan) (-0.75500,nan) (-0.75000,nan) (-0.74500,nan) (-0.74000,nan) (-0.73500,nan) (-0.73000,nan) (-0.72500,nan) (-0.72000,nan) (-0.71500,nan) (-0.71000,nan) (-0.70500,nan) (-0.70000,nan) (-0.69500,nan) (-0.69000,nan) (-0.68500,nan) (-0.68000,nan) (-0.67500,nan) (-0.67000,nan) (-0.66500,nan) (-0.66000,nan) (-0.65500,nan) (-0.65000,nan) (-0.64500,nan) (-0.64000,nan) (-0.63500,nan) (-0.63000,nan) (-0.62500,nan) (-0.62000,nan) (-0.61500,nan) (-0.61000,nan) (-0.60500,nan) (-0.60000,nan) (-0.59500,nan) (-0.59000,nan) (-0.58500,nan) (-0.58000,nan) (-0.57500,nan) (-0.57000,nan) (-0.56500,nan) (-0.56000,nan) (-0.55500,nan) (-0.55000,nan) (-0.54500,nan) (-0.54000,nan) (-0.53500,nan) (-0.53000,nan) (-0.52500,nan) (-0.52000,nan) (-0.51500,nan) (-0.51000,nan) (-0.50500,nan) (-0.50000,nan) (-0.49500,nan) (-0.49000,nan) (-0.48500,nan) (-0.48000,nan) (-0.47500,nan) (-0.47000,nan) (-0.46500,nan) (-0.46000,nan) (-0.45500,nan) (-0.45000,nan) (-0.44500,nan) (-0.44000,nan) (-0.43500,nan) (-0.43000,nan) (-0.42500,nan) (-0.42000,nan) (-0.41500,nan) (-0.41000,nan) (-0.40500,nan) (-0.40000,nan) (-0.39500,nan) (-0.39000,nan) (-0.38500,nan) (-0.38000,nan) (-0.37500,nan) (-0.37000,nan) (-0.36500,nan) (-0.36000,nan) (-0.35500,nan) (-0.35000,nan) (-0.34500,nan) (-0.34000,nan) (-0.33500,nan) (-0.33000,nan) (-0.32500,nan) (-0.32000,nan) (-0.31500,nan) (-0.31000,nan) (-0.30500,nan) (-0.30000,nan) (-0.29500,nan) (-0.29000,nan) (-0.28500,2.9984541) (-0.28000,2.9967421) (-0.27500,2.9950172) (-0.27000,2.9932792) (-0.26500,2.9915282) (-0.26000,2.9897641) (-0.25500,2.9879869) (-0.25000,2.9861964) (-0.24500,2.9843928) (-0.24000,2.9825758) (-0.23500,2.9807455) (-0.23000,2.9789019) (-0.22500,2.9770449) (-0.22000,2.9751745) (-0.21500,2.9732905) (-0.21000,2.9713931) (-0.20500,2.9694821) (-0.20000,2.9675575) (-0.19500,2.9656194) (-0.19000,2.9636675) (-0.18500,2.9617020) (-0.18000,2.9597227) (-0.17500,2.9577297) (-0.17000,2.9557229) (-0.16500,2.9537023) (-0.16000,2.9516678) (-0.15500,2.9496195) (-0.15000,2.9475573) (-0.14500,2.9454811) (-0.14000,2.9433909) (-0.13500,2.9412868) (-0.13000,2.9391687) (-0.12500,2.9370365) (-0.12000,2.9348903) (-0.11500,2.9327300) (-0.11000,2.9305556) (-0.10500,2.9283671) (-0.10000,2.9261644) (-0.09500,2.9239476) (-0.09000,2.9217166) (-0.08500,2.9194714) (-0.08000,2.9172120) (-0.07500,2.9149384) (-0.07000,2.9126505) (-0.06500,2.9103484) (-0.06000,2.9080320) (-0.05500,2.9057014) (-0.05000,2.9033564) (-0.04500,2.9009972) (-0.04000,2.8986236) (-0.03500,2.8962358) (-0.03000,2.8938336) (-0.02500,2.8914171) (-0.02000,2.8889863) (-0.01500,2.8865411) (-0.01000,2.8840816) (-0.00500,2.8816078) (0.00000,2.8791196) (0.00500,2.8766171) (0.01000,2.8741003) (0.01500,2.8715691) (0.02000,2.8690236) (0.02500,2.8664638) (0.03000,2.8638896) (0.03500,2.8613011) (0.04000,2.8586983) (0.04500,2.8560813) (0.05000,2.8534499) (0.05500,2.8508042) (0.06000,2.8481443) (0.06500,2.8454701) (0.07000,2.8427816) (0.07500,2.8400790) (0.08000,2.8373621) (0.08500,2.8346310) (0.09000,2.8318857) (0.09500,2.8291262) (0.10000,2.8263526) (0.10500,2.8235648) (0.11000,2.8207630) (0.11500,2.8179470) (0.12000,2.8151170) (0.12500,2.8122729) (0.13000,2.8094148) (0.13500,2.8065427) (0.14000,2.8036566) (0.14500,2.8007565) (0.15000,2.7978426) (0.15500,2.7949147) (0.16000,2.7919730) (0.16500,2.7890175) (0.17000,2.7860481) (0.17500,2.7830650) (0.18000,2.7800681) (0.18500,2.7770575) (0.19000,2.7740332) (0.19500,2.7709953) (0.20000,2.7679438) (0.20500,2.7648788) (0.21000,2.7618001) (0.21500,2.7587080) (0.22000,2.7556025) (0.22500,2.7524835) (0.23000,2.7493512) (0.23500,2.7462055) (0.24000,2.7430465) (0.24500,2.7398742) (0.25000,2.7366888) (0.25500,2.7334902) (0.26000,2.7302784) (0.26500,2.7270536) (0.27000,2.7238158) (0.27500,2.7205649) (0.28000,2.7173012) (0.28500,2.7140245) (0.29000,2.7107350) (0.29500,2.7074327) (0.30000,2.7041177) (0.30500,2.7007899) (0.31000,2.6974495) (0.31500,2.6940966) (0.32000,2.6907310) (0.32500,2.6873530) (0.33000,2.6839626) (0.33500,2.6805597) (0.34000,2.6771446) (0.34500,2.6737171) (0.35000,2.6702775) (0.35500,2.6668256) (0.36000,2.6633617) (0.36500,2.6598857) (0.37000,2.6563977) (0.37500,2.6528978) (0.38000,2.6493860) (0.38500,2.6458624) (0.39000,2.6423269) (0.39500,2.6387798) (0.40000,2.6352211) (0.40500,2.6316507) (0.41000,2.6280688) (0.41500,2.6244755) (0.42000,2.6208707) (0.42500,2.6172546) (0.43000,2.6136272) (0.43500,2.6099886) (0.44000,2.6063388) (0.44500,2.6026779) (0.45000,2.5990060) (0.45500,2.5953231) (0.46000,2.5916293) (0.46500,2.5879246) (0.47000,2.5842092) (0.47500,2.5804830) (0.48000,2.5767462) (0.48500,2.5729987) (0.49000,2.5692408) (0.49500,2.5654723) (0.50000,2.5616935) (0.50500,2.5579043) (0.51000,2.5541049) (0.51500,2.5502952) (0.52000,2.5464754) (0.52500,2.5426456) (0.53000,2.5388057) (0.53500,2.5349558) (0.54000,2.5310961) (0.54500,2.5272266) (0.55000,2.5233473) (0.55500,2.5194583) (0.56000,2.5155597) (0.56500,2.5116516) (0.57000,2.5077340) (0.57500,2.5038069) (0.58000,2.4998705) (0.58500,2.4959248) (0.59000,2.4919699) (0.59500,2.4880059) (0.60000,2.4840327) (0.60500,2.4800505) (0.61000,2.4760594) (0.61500,2.4720593) (0.62000,2.4680505) (0.62500,2.4640329) (0.63000,2.4600066) (0.63500,2.4559716) (0.64000,2.4519281) (0.64500,2.4478761) (0.65000,2.4438157) (0.65500,2.4397469) (0.66000,2.4356698) (0.66500,2.4315845) (0.67000,2.4274910) (0.67500,2.4233894) (0.68000,2.4192797) (0.68500,2.4151621) (0.69000,2.4110366) (0.69500,2.4069032) (0.70000,2.4027621) (0.70500,2.3986132) (0.71000,2.3944567) (0.71500,2.3902926) (0.72000,2.3861210) (0.72500,2.3819419) (0.73000,2.3777554) (0.73500,2.3735616) (0.74000,2.3693605) (0.74500,2.3651522) (0.75000,2.3609368) (0.75500,2.3567143) (0.76000,2.3524848) (0.76500,2.3482483) (0.77000,2.3440049) (0.77500,2.3397547) (0.78000,2.3354978) (0.78500,2.3312341) (0.79000,2.3269638) (0.79500,2.3226869) (0.80000,2.3184034) (0.80500,2.3141135) (0.81000,2.3098172) (0.81500,2.3055146) (0.82000,2.3012057) (0.82500,2.2968906) (0.83000,2.2925693) (0.83500,2.2882419) (0.84000,2.2839084) (0.84500,2.2795690) (0.85000,2.2752237) (0.85500,2.2708724) (0.86000,2.2665154) (0.86500,2.2621526) (0.87000,2.2577842) (0.87500,2.2534101) (0.88000,2.2490304) (0.88500,2.2446452) (0.89000,2.2402545) (0.89500,2.2358584) (0.90000,2.2314570) (0.90500,2.2270503) (0.91000,2.2226383) (0.91500,2.2182212) (0.92000,2.2137989) (0.92500,2.2093716) (0.93000,2.2049392) (0.93500,2.2005019) (0.94000,2.1960596) (0.94500,2.1916125) (0.95000,2.1871606) (0.95500,2.1827040) (0.96000,2.1782426) (0.96500,2.1737766) (0.97000,2.1693060) (0.97500,2.1648309) (0.98000,2.1603513) (0.98500,2.1558673) (0.99000,2.1513789) (0.99500,2.1468861) (1.00000,2.1423891) (1.00500,2.1378878) (1.01000,2.1333824) (1.01500,2.1288728) (1.02000,2.1243592) (1.02500,2.1198415) (1.03000,2.1153199) (1.03500,2.1107943) (1.04000,2.1062649) (1.04500,2.1017316) (1.05000,2.0971945) (1.05500,2.0926537) (1.06000,2.0881092) (1.06500,2.0835611) (1.07000,2.0790094) (1.07500,2.0744542) (1.08000,2.0698954) (1.08500,2.0653332) (1.09000,2.0607677) (1.09500,2.0561987) (1.10000,2.0516265) (1.10500,2.0470510) (1.11000,2.0424722) (1.11500,2.0378903) (1.12000,2.0333053) (1.12500,2.0287172) (1.13000,2.0241261) (1.13500,2.0195320) (1.14000,2.0149349) (1.14500,2.0103350) (1.15000,2.0057321) (1.15500,2.0011265) (1.16000,1.9965181) (1.16500,1.9919070) (1.17000,1.9872932) (1.17500,1.9826767) (1.18000,1.9780577) (1.18500,1.9734361) (1.19000,1.9688120) (1.19500,1.9641854) (1.20000,1.9595565) (1.20500,1.9549251) (1.21000,1.9502914) (1.21500,1.9456554) (1.22000,1.9410171) (1.22500,1.9363766) (1.23000,1.9317339) (1.23500,1.9270891) (1.24000,1.9224423) (1.24500,1.9177933) (1.25000,1.9131424) (1.25500,1.9084894) (1.26000,1.9038346) (1.26500,1.8991778) (1.27000,1.8945192) (1.27500,1.8898588) (1.28000,1.8851966) (1.28500,1.8805327) (1.29000,1.8758671) (1.29500,1.8711998) (1.30000,1.8665309) (1.30500,1.8618605) (1.31000,1.8571885) (1.31500,1.8525150) (1.32000,1.8478400) (1.32500,1.8431636) (1.33000,1.8384859) (1.33500,1.8338068) (1.34000,1.8291263) (1.34500,1.8244447) (1.35000,1.8197617) (1.35500,1.8150776) (1.36000,1.8103924) (1.36500,1.8057060) (1.37000,1.8010185) (1.37500,1.7963300) (1.38000,1.7916405) (1.38500,1.7869500) (1.39000,1.7822586) (1.39500,1.7775662) (1.40000,1.7728731) (1.40500,1.7681791) (1.41000,1.7634843) (1.41500,1.7587888) (1.42000,1.7540925) (1.42500,1.7493956) (1.43000,1.7446981) (1.43500,1.7399999) (1.44000,1.7353012) (1.44500,1.7306020) (1.45000,1.7259023) (1.45500,1.7212021) (1.46000,1.7165015) (1.46500,1.7118005) (1.47000,1.7070992) (1.47500,1.7023975) (1.48000,1.6976956) (1.48500,1.6929935) (1.49000,1.6882911) (1.49500,1.6835886) (1.50000,1.6788860) (1.50500,1.6741832) (1.51000,1.6694805) (1.51500,1.6647776) (1.52000,1.6600749) (1.52500,1.6553721) (1.53000,1.6506695) (1.53500,1.6459670) (1.54000,1.6412646) (1.54500,1.6365625) (1.55000,1.6318605) (1.55500,1.6271589) (1.56000,1.6224576) (1.56500,1.6177566) (1.57000,1.6130559) (1.57500,1.6083557) (1.58000,1.6036560) (1.58500,1.5989567) (1.59000,1.5942580) (1.59500,1.5895599) (1.60000,1.5848623) (1.60500,1.5801654) (1.61000,1.5754691) (1.61500,1.5707736) (1.62000,1.5660788) (1.62500,1.5613848) (1.63000,1.5566916) (1.63500,1.5519993) (1.64000,1.5473079) (1.64500,1.5426174) (1.65000,1.5379279) (1.65500,1.5332394) (1.66000,1.5285519) (1.66500,1.5238656) (1.67000,1.5191803) (1.67500,1.5144962) (1.68000,1.5098134) (1.68500,1.5051317) (1.69000,1.5004513) (1.69500,1.4957723) (1.70000,1.4910946) (1.70500,1.4864183) (1.71000,1.4817434) (1.71500,1.4770700) (1.72000,1.4723981) (1.72500,1.4677278) (1.73000,1.4630590) (1.73500,1.4583919) (1.74000,1.4537265) (1.74500,1.4490627) (1.75000,1.4444007) (1.75500,1.4397405) (1.76000,1.4350821) (1.76500,1.4304256) (1.77000,1.4257710) (1.77500,1.4211184) (1.78000,1.4164677) (1.78500,1.4118191) (1.79000,1.4071726) (1.79500,1.4025281) (1.80000,1.3978859) (1.80500,1.3932458) (1.81000,1.3886080) (1.81500,1.3839724) (1.82000,1.3793392) (1.82500,1.3747084) (1.83000,1.3700800) (1.83500,1.3654540) (1.84000,1.3608305) (1.84500,1.3562096) (1.85000,1.3515913) (1.85500,1.3469756) (1.86000,1.3423626) (1.86500,1.3377523) (1.87000,1.3331447) (1.87500,1.3285400) (1.88000,1.3239382) (1.88500,1.3193392) (1.89000,1.3147432) (1.89500,1.3101502) (1.90000,1.3055602) (1.90500,1.3009733) (1.91000,1.2963896) (1.91500,1.2918090) (1.92000,1.2872317) (1.92500,1.2826576) (1.93000,1.2780869) (1.93500,1.2735195) (1.94000,1.2689556) (1.94500,1.2643951) (1.95000,1.2598382) (1.95500,1.2552848) (1.96000,1.2507351) (1.96500,1.2461890) (1.97000,1.2416466) (1.97500,1.2371080) (1.98000,1.2325732) (1.98500,1.2280423) (1.99000,1.2235153) (1.99500,1.2189923) (2.00000,1.2144733) (2.00500,1.2099584) (2.01000,1.2054476) (2.01500,1.2009411) (2.02000,1.1964387) (2.02500,1.1919406) (2.03000,1.1874469) (2.03500,1.1829576) (2.04000,1.1784728) (2.04500,1.1739924) (2.05000,1.1695166) (2.05500,1.1650454) (2.06000,1.1605790) (2.06500,1.1561172) (2.07000,1.1516602) (2.07500,1.1472081) (2.08000,1.1427609) (2.08500,1.1383186) (2.09000,1.1338814) (2.09500,1.1294492) (2.10000,1.1250222) (2.10500,1.1206004) (2.11000,1.1161838) (2.11500,1.1117725) (2.12000,1.1073666) (2.12500,1.1029662) (2.13000,1.0985712) (2.13500,1.0941818) (2.14000,1.0897980) (2.14500,1.0854199) (2.15000,1.0810476) (2.15500,1.0766810) (2.16000,1.0723203) (2.16500,1.0679656) (2.17000,1.0636168) (2.17500,1.0592741) (2.18000,1.0549375) (2.18500,1.0506072) (2.19000,1.0462830) (2.19500,1.0419652) (2.20000,1.0376538) (2.20500,1.0333488) (2.21000,1.0290504) (2.21500,1.0247585) (2.22000,1.0204733) (2.22500,1.0161948) (2.23000,1.0119232) (2.23500,1.0076583) (2.24000,1.0034004) (2.24500,0.9991495) (2.25000,0.9949057) (2.25500,0.9906690) (2.26000,0.9864395) (2.26500,0.9822173) (2.27000,0.9780025) (2.27500,0.9737951) (2.28000,0.9695951) (2.28500,0.9654028) (2.29000,0.9612181) (2.29500,0.9570411) (2.30000,0.9528719) (2.30500,0.9487105) (2.31000,0.9445571) (2.31500,0.9404118) (2.32000,0.9362745) (2.32500,0.9321453) (2.33000,0.9280245) (2.33500,0.9239119) (2.34000,0.9198077) (2.34500,0.9157120) (2.35000,0.9116249) (2.35500,0.9075464) (2.36000,0.9034766) (2.36500,0.8994156) (2.37000,0.8953634) (2.37500,0.8913202) (2.38000,0.8872861) (2.38500,0.8832610) (2.39000,0.8792451) (2.39500,0.8752385) (2.40000,0.8712413) (2.40500,0.8672534) (2.41000,0.8632751) (2.41500,0.8593064) (2.42000,0.8553474) (2.42500,0.8513982) (2.43000,0.8474588) (2.43500,0.8435293) (2.44000,0.8396098) (2.44500,0.8357005) (2.45000,0.8318014) (2.45500,0.8279125) (2.46000,0.8240340) (2.46500,0.8201659) (2.47000,0.8163084) (2.47500,0.8124615) (2.48000,0.8086254) (2.48500,0.8048000) (2.49000,0.8009855) (2.49500,0.7971820) (2.50000,0.7933896) (2.50500,0.7896083) (2.51000,0.7858383) (2.51500,0.7820796) (2.52000,0.7783324) (2.52500,0.7745967) (2.53000,0.7708726) (2.53500,0.7671602) (2.54000,0.7634596) (2.54500,0.7597709) (2.55000,0.7560942) (2.55500,0.7524296) (2.56000,0.7487772) (2.56500,0.7451370) (2.57000,0.7415092) (2.57500,0.7378939) (2.58000,0.7342912) (2.58500,0.7307010) (2.59000,0.7271237) (2.59500,0.7235592) (2.60000,0.7200076) (2.60500,0.7164691) (2.61000,0.7129438) (2.61500,0.7094317) (2.62000,0.7059329) (2.62500,0.7024476) (2.63000,0.6989759) (2.63500,0.6955178) (2.64000,0.6920734) (2.64500,0.6886429) (2.65000,0.6852264) (2.65500,0.6818239) (2.66000,0.6784356) (2.66500,0.6750616) (2.67000,0.6717020) (2.67500,0.6683568) (2.68000,0.6650262) (2.68500,0.6617103) (2.69000,0.6584093) (2.69500,0.6551231) (2.70000,0.6518520) (2.70500,0.6485960) (2.71000,0.6453552) (2.71500,0.6421298) (2.72000,0.6389199) (2.72500,0.6357255) (2.73000,0.6325469) (2.73500,0.6293840) (2.74000,0.6262371) (2.74500,0.6231062) (2.75000,0.6199915) (2.75500,0.6168931) (2.76000,0.6138110) (2.76500,0.6107455) (2.77000,0.6076966) (2.77500,0.6046645) (2.78000,0.6016493) (2.78500,0.5986510) (2.79000,0.5956699) (2.79500,0.5927061) (2.80000,0.5897597) (2.80500,0.5868308) (2.81000,0.5839195) (2.81500,0.5810261) (2.82000,0.5781505) (2.82500,0.5752931) (2.83000,0.5724538) (2.83500,0.5696328) (2.84000,0.5668304) (2.84500,0.5640465) (2.85000,0.5612814) (2.85500,0.5585353) (2.86000,0.5558082) (2.86500,0.5531003) (2.87000,0.5504118) (2.87500,0.5477428) (2.88000,0.5450935) (2.88500,0.5424640) (2.89000,0.5398546) (2.89500,0.5372653) (2.90000,0.5346964) (2.90500,0.5321480) (2.91000,0.5296203) (2.91500,0.5271134) (2.92000,0.5246277) (2.92500,0.5221631) (2.93000,0.5197200) (2.93500,0.5172986) (2.94000,0.5148989) (2.94500,0.5125213) (2.95000,0.5101659) (2.95500,0.5078330) (2.96000,0.5055227) (2.96500,0.5032353) (2.97000,0.5009710) (2.97500,0.4987301) (2.98000,0.4965127) (2.98500,0.4943192) (2.99000,0.4921498) (2.99500,0.4900047) (3.00000,0.4878841) (3.00500,0.4857885) (3.01000,0.4837180) (3.01500,0.4816729) (3.02000,0.4796536) (3.02500,0.4776603) (3.03000,0.4756933) (3.03500,0.4737530) (3.04000,0.4718398) (3.04500,0.4699538) (3.05000,0.4680956) (3.05500,0.4662654) (3.06000,0.4644637) (3.06500,0.4626908) (3.07000,0.4609471) (3.07500,0.4592331) (3.08000,0.4575492) (3.08500,0.4558958) (3.09000,0.4542734) (3.09500,0.4526825) (3.10000,0.4511235) (3.10500,0.4495971) (3.11000,0.4481037) (3.11500,0.4466439) (3.12000,0.4452183) (3.12500,0.4438274) (3.13000,0.4424719) (3.13500,0.4411525) (3.14000,0.4398699) (3.14500,0.4386247) (3.15000,0.4374177) (3.15500,0.4362496) (3.16000,0.4351213) (3.16500,0.4340337) (3.17000,0.4329875) (3.17500,0.4319837) (3.18000,0.4310233) (3.18500,0.4301073) (3.19000,0.4292367) (3.19500,0.4284127) (3.20000,0.4276364) (3.20500,0.4269090) (3.21000,0.4262318) (3.21500,0.4256062) (3.22000,0.4250335) (3.22500,0.4245153) (3.23000,0.4240531) (3.23500,0.4236486) (3.24000,0.4233035) (3.24500,0.4230196) (3.25000,0.4227990) (3.25500,0.4226436) (3.26000,0.4225556) (3.26500,0.4225373) (3.27000,0.4225912) (3.27500,0.4227198) (3.28000,0.4229259) (3.28500,0.4232123) (3.29000,0.4235823) (3.29500,0.4240390) (3.30000,0.4245860) (3.30500,0.4252271) (3.31000,0.4259662) (3.31500,0.4268076) (3.32000,0.4277560) (3.32500,0.4288161) (3.33000,0.4299934) (3.33500,0.4312934) (3.34000,0.4327223) (3.34500,0.4342866) (3.35000,0.4359933) (3.35500,0.4378503) (3.36000,0.4398656) (3.36500,0.4420484) (3.37000,0.4444084) (3.37500,0.4469562) (3.38000,0.4497034) (3.38500,0.4526626) (3.39000,0.4558478) (3.39500,0.4592741) (3.40000,0.4629582) (3.40500,0.4669187) (3.41000,0.4711759) (3.41500,0.4757525) (3.42000,0.4806736) (3.42500,0.4859674) (3.43000,0.4916653) (3.43500,0.4978026) (3.44000,0.5044193) (3.44500,0.5115606) (3.45000,0.5192779) (3.45500,0.5276302) (3.46000,0.5366851) (3.46500,0.5465210) (3.47000,0.5572291) (3.47500,0.5689163) (3.48000,0.5817091) (3.48500,0.5957582) (3.49000,0.6112449) (3.49500,0.6283900) (3.50000,0.6474655) (3.50500,0.6688114) (3.51000,0.6928594) (3.51500,0.7201692) (3.52000,0.7514822) (3.52500,0.7878087) (3.53000,0.8305734) (3.53500,0.8818734) (3.54000,0.9449757) (3.54500,1.0253775) (3.55000,1.1334290) (3.55500,1.2924784) (3.56000,1.5776097) (3.56500,nan) (3.57000,nan) (3.57500,nan) (3.58000,nan) (3.58500,nan) (3.59000,nan) (3.59500,nan) (3.60000,nan) (3.60500,nan) (3.61000,nan) (3.61500,nan) (3.62000,nan) (3.62500,nan) (3.63000,nan) (3.63500,nan) (3.64000,nan) (3.64500,nan) (3.65000,nan) (3.65500,nan) (3.66000,nan) (3.66500,nan) (3.67000,nan) (3.67500,nan) (3.68000,nan) (3.68500,-1.7044149) (3.69000,-1.1231966) (3.69500,-0.8985377) (3.70000,-0.7611874) (3.70500,-0.6641370) (3.71000,-0.5902670) (3.71500,-0.5314026) (3.72000,-0.4830044) (3.72500,-0.4422929) (3.73000,-0.4074450) (3.73500,-0.3772022) (3.74000,-0.3506602) (3.74500,-0.3271486) (3.75000,-0.3061574) (3.75500,-0.2872899) (3.76000,-0.2702321) (3.76500,-0.2547315) (3.77000,-0.2405821) (3.77500,-0.2276141) (3.78000,-0.2156858) (3.78500,-0.2046782) (3.79000,-0.1944900) (3.79500,-0.1850349) (3.80000,-0.1762384) (3.80500,-0.1680361) (3.81000,-0.1603718) (3.81500,-0.1531964) (3.82000,-0.1464666) (3.82500,-0.1401441) (3.83000,-0.1341950) (3.83500,-0.1285891) (3.84000,-0.1232993) (3.84500,-0.1183014) (3.85000,-0.1135736) (3.85500,-0.1090962) (3.86000,-0.1048513) (3.86500,-0.1008228) (3.87000,-0.0969959) (3.87500,-0.0933572) (3.88000,-0.0898945) (3.88500,-0.0865966) (3.89000,-0.0834530) (3.89500,-0.0804544) (3.90000,-0.0775920) (3.90500,-0.0748578) (3.91000,-0.0722443) (3.91500,-0.0697446) (3.92000,-0.0673523) (3.92500,-0.0650615) (3.93000,-0.0628668) (3.93500,-0.0607629) (3.94000,-0.0587451) (3.94500,-0.0568089) (3.95000,-0.0549501) (3.95500,-0.0531649) (3.96000,-0.0514496) (3.96500,-0.0498008) (3.97000,-0.0482152) (3.97500,-0.0466898) (3.98000,-0.0452218) (3.98500,-0.0438086) (3.99000,-0.0424475) (3.99500,-0.0411363) (4.00000,-0.0398727)};\addlegendentry{R}
\addplot[thick,orange] coordinates {(-2.00000,nan) (-1.99500,nan) (-1.99000,nan) (-1.98500,nan) (-1.98000,nan) (-1.97500,nan) (-1.97000,nan) (-1.96500,nan) (-1.96000,nan) (-1.95500,nan) (-1.95000,nan) (-1.94500,nan) (-1.94000,nan) (-1.93500,nan) (-1.93000,nan) (-1.92500,nan) (-1.92000,nan) (-1.91500,nan) (-1.91000,nan) (-1.90500,nan) (-1.90000,nan) (-1.89500,nan) (-1.89000,nan) (-1.88500,nan) (-1.88000,nan) (-1.87500,nan) (-1.87000,nan) (-1.86500,nan) (-1.86000,nan) (-1.85500,nan) (-1.85000,nan) (-1.84500,nan) (-1.84000,nan) (-1.83500,nan) (-1.83000,nan) (-1.82500,nan) (-1.82000,nan) (-1.81500,nan) (-1.81000,nan) (-1.80500,nan) (-1.80000,nan) (-1.79500,nan) (-1.79000,nan) (-1.78500,nan) (-1.78000,nan) (-1.77500,nan) (-1.77000,nan) (-1.76500,nan) (-1.76000,nan) (-1.75500,nan) (-1.75000,nan) (-1.74500,nan) (-1.74000,nan) (-1.73500,nan) (-1.73000,nan) (-1.72500,nan) (-1.72000,nan) (-1.71500,nan) (-1.71000,nan) (-1.70500,nan) (-1.70000,nan) (-1.69500,nan) (-1.69000,nan) (-1.68500,nan) (-1.68000,nan) (-1.67500,nan) (-1.67000,nan) (-1.66500,nan) (-1.66000,nan) (-1.65500,nan) (-1.65000,nan) (-1.64500,nan) (-1.64000,nan) (-1.63500,nan) (-1.63000,nan) (-1.62500,nan) (-1.62000,nan) (-1.61500,nan) (-1.61000,nan) (-1.60500,nan) (-1.60000,nan) (-1.59500,nan) (-1.59000,nan) (-1.58500,nan) (-1.58000,nan) (-1.57500,nan) (-1.57000,nan) (-1.56500,nan) (-1.56000,nan) (-1.55500,nan) (-1.55000,nan) (-1.54500,nan) (-1.54000,nan) (-1.53500,nan) (-1.53000,nan) (-1.52500,nan) (-1.52000,nan) (-1.51500,nan) (-1.51000,nan) (-1.50500,nan) (-1.50000,nan) (-1.49500,nan) (-1.49000,nan) (-1.48500,nan) (-1.48000,nan) (-1.47500,nan) (-1.47000,nan) (-1.46500,nan) (-1.46000,nan) (-1.45500,nan) (-1.45000,nan) (-1.44500,nan) (-1.44000,nan) (-1.43500,nan) (-1.43000,nan) (-1.42500,nan) (-1.42000,nan) (-1.41500,nan) (-1.41000,nan) (-1.40500,nan) (-1.40000,nan) (-1.39500,nan) (-1.39000,nan) (-1.38500,nan) (-1.38000,nan) (-1.37500,nan) (-1.37000,nan) (-1.36500,nan) (-1.36000,nan) (-1.35500,nan) (-1.35000,nan) (-1.34500,nan) (-1.34000,nan) (-1.33500,nan) (-1.33000,nan) (-1.32500,nan) (-1.32000,nan) (-1.31500,nan) (-1.31000,nan) (-1.30500,nan) (-1.30000,nan) (-1.29500,nan) (-1.29000,nan) (-1.28500,nan) (-1.28000,nan) (-1.27500,nan) (-1.27000,nan) (-1.26500,nan) (-1.26000,nan) (-1.25500,nan) (-1.25000,nan) (-1.24500,nan) (-1.24000,nan) (-1.23500,nan) (-1.23000,nan) (-1.22500,nan) (-1.22000,nan) (-1.21500,nan) (-1.21000,nan) (-1.20500,nan) (-1.20000,nan) (-1.19500,nan) (-1.19000,nan) (-1.18500,nan) (-1.18000,nan) (-1.17500,nan) (-1.17000,nan) (-1.16500,nan) (-1.16000,nan) (-1.15500,nan) (-1.15000,nan) (-1.14500,nan) (-1.14000,nan) (-1.13500,nan) (-1.13000,nan) (-1.12500,nan) (-1.12000,nan) (-1.11500,nan) (-1.11000,nan) (-1.10500,nan) (-1.10000,nan) (-1.09500,nan) (-1.09000,nan) (-1.08500,nan) (-1.08000,nan) (-1.07500,nan) (-1.07000,nan) (-1.06500,nan) (-1.06000,nan) (-1.05500,nan) (-1.05000,nan) (-1.04500,nan) (-1.04000,nan) (-1.03500,nan) (-1.03000,nan) (-1.02500,nan) (-1.02000,nan) (-1.01500,nan) (-1.01000,nan) (-1.00500,nan) (-1.00000,nan) (-0.99500,nan) (-0.99000,nan) (-0.98500,nan) (-0.98000,nan) (-0.97500,nan) (-0.97000,nan) (-0.96500,nan) (-0.96000,nan) (-0.95500,nan) (-0.95000,nan) (-0.94500,nan) (-0.94000,nan) (-0.93500,nan) (-0.93000,nan) (-0.92500,nan) (-0.92000,nan) (-0.91500,nan) (-0.91000,nan) (-0.90500,nan) (-0.90000,nan) (-0.89500,nan) (-0.89000,nan) (-0.88500,nan) (-0.88000,nan) (-0.87500,nan) (-0.87000,nan) (-0.86500,nan) (-0.86000,nan) (-0.85500,nan) (-0.85000,nan) (-0.84500,nan) (-0.84000,nan) (-0.83500,nan) (-0.83000,nan) (-0.82500,nan) (-0.82000,nan) (-0.81500,nan) (-0.81000,nan) (-0.80500,nan) (-0.80000,nan) (-0.79500,nan) (-0.79000,nan) (-0.78500,nan) (-0.78000,nan) (-0.77500,nan) (-0.77000,nan) (-0.76500,nan) (-0.76000,nan) (-0.75500,nan) (-0.75000,nan) (-0.74500,nan) (-0.74000,nan) (-0.73500,nan) (-0.73000,nan) (-0.72500,nan) (-0.72000,nan) (-0.71500,nan) (-0.71000,nan) (-0.70500,nan) (-0.70000,nan) (-0.69500,nan) (-0.69000,nan) (-0.68500,nan) (-0.68000,nan) (-0.67500,nan) (-0.67000,nan) (-0.66500,nan) (-0.66000,nan) (-0.65500,nan) (-0.65000,nan) (-0.64500,nan) (-0.64000,nan) (-0.63500,nan) (-0.63000,nan) (-0.62500,nan) (-0.62000,nan) (-0.61500,nan) (-0.61000,nan) (-0.60500,nan) (-0.60000,nan) (-0.59500,nan) (-0.59000,nan) (-0.58500,nan) (-0.58000,nan) (-0.57500,nan) (-0.57000,nan) (-0.56500,nan) (-0.56000,nan) (-0.55500,nan) (-0.55000,nan) (-0.54500,nan) (-0.54000,nan) (-0.53500,nan) (-0.53000,nan) (-0.52500,nan) (-0.52000,nan) (-0.51500,nan) (-0.51000,nan) (-0.50500,nan) (-0.50000,nan) (-0.49500,nan) (-0.49000,nan) (-0.48500,nan) (-0.48000,nan) (-0.47500,nan) (-0.47000,nan) (-0.46500,nan) (-0.46000,nan) (-0.45500,nan) (-0.45000,nan) (-0.44500,nan) (-0.44000,nan) (-0.43500,nan) (-0.43000,nan) (-0.42500,nan) (-0.42000,nan) (-0.41500,nan) (-0.41000,nan) (-0.40500,nan) (-0.40000,nan) (-0.39500,nan) (-0.39000,nan) (-0.38500,nan) (-0.38000,nan) (-0.37500,nan) (-0.37000,nan) (-0.36500,nan) (-0.36000,nan) (-0.35500,nan) (-0.35000,nan) (-0.34500,nan) (-0.34000,nan) (-0.33500,nan) (-0.33000,nan) (-0.32500,nan) (-0.32000,nan) (-0.31500,nan) (-0.31000,nan) (-0.30500,nan) (-0.30000,nan) (-0.29500,nan) (-0.29000,nan) (-0.28500,nan) (-0.28000,nan) (-0.27500,nan) (-0.27000,nan) (-0.26500,nan) (-0.26000,nan) (-0.25500,nan) (-0.25000,nan) (-0.24500,nan) (-0.24000,nan) (-0.23500,nan) (-0.23000,nan) (-0.22500,nan) (-0.22000,nan) (-0.21500,nan) (-0.21000,nan) (-0.20500,nan) (-0.20000,nan) (-0.19500,nan) (-0.19000,nan) (-0.18500,nan) (-0.18000,nan) (-0.17500,nan) (-0.17000,nan) (-0.16500,nan) (-0.16000,nan) (-0.15500,nan) (-0.15000,nan) (-0.14500,nan) (-0.14000,nan) (-0.13500,nan) (-0.13000,nan) (-0.12500,nan) (-0.12000,nan) (-0.11500,nan) (-0.11000,nan) (-0.10500,nan) (-0.10000,nan) (-0.09500,nan) (-0.09000,nan) (-0.08500,nan) (-0.08000,nan) (-0.07500,nan) (-0.07000,nan) (-0.06500,nan) (-0.06000,nan) (-0.05500,nan) (-0.05000,nan) (-0.04500,nan) (-0.04000,nan) (-0.03500,nan) (-0.03000,nan) (-0.02500,nan) (-0.02000,nan) (-0.01500,nan) (-0.01000,nan) (-0.00500,nan) (0.00000,nan) (0.00500,nan) (0.01000,nan) (0.01500,nan) (0.02000,nan) (0.02500,nan) (0.03000,nan) (0.03500,nan) (0.04000,nan) (0.04500,nan) (0.05000,nan) (0.05500,nan) (0.06000,nan) (0.06500,nan) (0.07000,nan) (0.07500,nan) (0.08000,nan) (0.08500,nan) (0.09000,nan) (0.09500,nan) (0.10000,nan) (0.10500,nan) (0.11000,nan) (0.11500,nan) (0.12000,nan) (0.12500,nan) (0.13000,nan) (0.13500,nan) (0.14000,nan) (0.14500,nan) (0.15000,nan) (0.15500,nan) (0.16000,nan) (0.16500,nan) (0.17000,nan) (0.17500,nan) (0.18000,nan) (0.18500,nan) (0.19000,nan) (0.19500,nan) (0.20000,nan) (0.20500,nan) (0.21000,nan) (0.21500,nan) (0.22000,nan) (0.22500,nan) (0.23000,nan) (0.23500,nan) (0.24000,nan) (0.24500,nan) (0.25000,nan) (0.25500,nan) (0.26000,nan) (0.26500,nan) (0.27000,nan) (0.27500,nan) (0.28000,nan) (0.28500,nan) (0.29000,nan) (0.29500,nan) (0.30000,nan) (0.30500,nan) (0.31000,nan) (0.31500,nan) (0.32000,nan) (0.32500,nan) (0.33000,nan) (0.33500,nan) (0.34000,nan) (0.34500,nan) (0.35000,nan) (0.35500,nan) (0.36000,nan) (0.36500,nan) (0.37000,nan) (0.37500,nan) (0.38000,nan) (0.38500,nan) (0.39000,nan) (0.39500,nan) (0.40000,nan) (0.40500,nan) (0.41000,nan) (0.41500,nan) (0.42000,nan) (0.42500,nan) (0.43000,nan) (0.43500,nan) (0.44000,nan) (0.44500,nan) (0.45000,nan) (0.45500,nan) (0.46000,nan) (0.46500,nan) (0.47000,nan) (0.47500,nan) (0.48000,nan) (0.48500,nan) (0.49000,nan) (0.49500,nan) (0.50000,nan) (0.50500,nan) (0.51000,nan) (0.51500,nan) (0.52000,nan) (0.52500,nan) (0.53000,nan) (0.53500,nan) (0.54000,nan) (0.54500,nan) (0.55000,nan) (0.55500,nan) (0.56000,nan) (0.56500,nan) (0.57000,nan) (0.57500,nan) (0.58000,nan) (0.58500,nan) (0.59000,nan) (0.59500,nan) (0.60000,nan) (0.60500,nan) (0.61000,nan) (0.61500,nan) (0.62000,nan) (0.62500,nan) (0.63000,nan) (0.63500,nan) (0.64000,nan) (0.64500,nan) (0.65000,nan) (0.65500,nan) (0.66000,nan) (0.66500,nan) (0.67000,nan) (0.67500,nan) (0.68000,nan) (0.68500,nan) (0.69000,nan) (0.69500,nan) (0.70000,nan) (0.70500,nan) (0.71000,nan) (0.71500,nan) (0.72000,nan) (0.72500,nan) (0.73000,nan) (0.73500,nan) (0.74000,nan) (0.74500,nan) (0.75000,nan) (0.75500,nan) (0.76000,nan) (0.76500,nan) (0.77000,nan) (0.77500,nan) (0.78000,nan) (0.78500,nan) (0.79000,nan) (0.79500,nan) (0.80000,nan) (0.80500,nan) (0.81000,nan) (0.81500,nan) (0.82000,nan) (0.82500,nan) (0.83000,nan) (0.83500,nan) (0.84000,nan) (0.84500,nan) (0.85000,nan) (0.85500,nan) (0.86000,nan) (0.86500,nan) (0.87000,nan) (0.87500,nan) (0.88000,nan) (0.88500,nan) (0.89000,nan) (0.89500,nan) (0.90000,nan) (0.90500,nan) (0.91000,nan) (0.91500,nan) (0.92000,nan) (0.92500,nan) (0.93000,nan) (0.93500,nan) (0.94000,nan) (0.94500,nan) (0.95000,nan) (0.95500,nan) (0.96000,nan) (0.96500,nan) (0.97000,nan) (0.97500,nan) (0.98000,nan) (0.98500,nan) (0.99000,nan) (0.99500,nan) (1.00000,nan) (1.00500,nan) (1.01000,nan) (1.01500,nan) (1.02000,nan) (1.02500,nan) (1.03000,nan) (1.03500,nan) (1.04000,nan) (1.04500,nan) (1.05000,nan) (1.05500,nan) (1.06000,nan) (1.06500,nan) (1.07000,nan) (1.07500,nan) (1.08000,nan) (1.08500,nan) (1.09000,nan) (1.09500,nan) (1.10000,nan) (1.10500,nan) (1.11000,nan) (1.11500,nan) (1.12000,nan) (1.12500,nan) (1.13000,nan) (1.13500,nan) (1.14000,nan) (1.14500,nan) (1.15000,nan) (1.15500,nan) (1.16000,nan) (1.16500,nan) (1.17000,nan) (1.17500,nan) (1.18000,nan) (1.18500,nan) (1.19000,nan) (1.19500,nan) (1.20000,nan) (1.20500,nan) (1.21000,nan) (1.21500,nan) (1.22000,nan) (1.22500,nan) (1.23000,nan) (1.23500,nan) (1.24000,nan) (1.24500,nan) (1.25000,nan) (1.25500,nan) (1.26000,nan) (1.26500,nan) (1.27000,nan) (1.27500,nan) (1.28000,nan) (1.28500,nan) (1.29000,nan) (1.29500,nan) (1.30000,nan) (1.30500,nan) (1.31000,nan) (1.31500,nan) (1.32000,nan) (1.32500,nan) (1.33000,nan) (1.33500,nan) (1.34000,nan) (1.34500,nan) (1.35000,nan) (1.35500,nan) (1.36000,nan) (1.36500,nan) (1.37000,nan) (1.37500,nan) (1.38000,nan) (1.38500,nan) (1.39000,nan) (1.39500,nan) (1.40000,nan) (1.40500,nan) (1.41000,nan) (1.41500,nan) (1.42000,nan) (1.42500,nan) (1.43000,nan) (1.43500,nan) (1.44000,nan) (1.44500,nan) (1.45000,nan) (1.45500,nan) (1.46000,nan) (1.46500,nan) (1.47000,nan) (1.47500,nan) (1.48000,nan) (1.48500,nan) (1.49000,nan) (1.49500,nan) (1.50000,nan) (1.50500,nan) (1.51000,nan) (1.51500,nan) (1.52000,nan) (1.52500,nan) (1.53000,nan) (1.53500,nan) (1.54000,nan) (1.54500,nan) (1.55000,nan) (1.55500,nan) (1.56000,nan) (1.56500,nan) (1.57000,nan) (1.57500,nan) (1.58000,nan) (1.58500,nan) (1.59000,nan) (1.59500,nan) (1.60000,nan) (1.60500,nan) (1.61000,nan) (1.61500,nan) (1.62000,nan) (1.62500,nan) (1.63000,nan) (1.63500,nan) (1.64000,nan) (1.64500,nan) (1.65000,nan) (1.65500,nan) (1.66000,nan) (1.66500,nan) (1.67000,nan) (1.67500,nan) (1.68000,nan) (1.68500,nan) (1.69000,nan) (1.69500,nan) (1.70000,nan) (1.70500,nan) (1.71000,nan) (1.71500,nan) (1.72000,nan) (1.72500,nan) (1.73000,nan) (1.73500,nan) (1.74000,nan) (1.74500,nan) (1.75000,nan) (1.75500,nan) (1.76000,nan) (1.76500,nan) (1.77000,nan) (1.77500,nan) (1.78000,nan) (1.78500,nan) (1.79000,nan) (1.79500,nan) (1.80000,nan) (1.80500,nan) (1.81000,nan) (1.81500,nan) (1.82000,nan) (1.82500,nan) (1.83000,nan) (1.83500,nan) (1.84000,nan) (1.84500,nan) (1.85000,nan) (1.85500,nan) (1.86000,nan) (1.86500,nan) (1.87000,nan) (1.87500,nan) (1.88000,nan) (1.88500,nan) (1.89000,nan) (1.89500,nan) (1.90000,nan) (1.90500,nan) (1.91000,nan) (1.91500,nan) (1.92000,nan) (1.92500,nan) (1.93000,nan) (1.93500,nan) (1.94000,nan) (1.94500,nan) (1.95000,nan) (1.95500,nan) (1.96000,nan) (1.96500,nan) (1.97000,nan) (1.97500,nan) (1.98000,nan) (1.98500,nan) (1.99000,nan) (1.99500,nan) (2.00000,nan) (2.00500,nan) (2.01000,nan) (2.01500,nan) (2.02000,nan) (2.02500,nan) (2.03000,nan) (2.03500,nan) (2.04000,nan) (2.04500,nan) (2.05000,nan) (2.05500,nan) (2.06000,nan) (2.06500,nan) (2.07000,nan) (2.07500,nan) (2.08000,nan) (2.08500,nan) (2.09000,nan) (2.09500,nan) (2.10000,nan) (2.10500,nan) (2.11000,nan) (2.11500,nan) (2.12000,nan) (2.12500,nan) (2.13000,nan) (2.13500,nan) (2.14000,nan) (2.14500,nan) (2.15000,nan) (2.15500,nan) (2.16000,nan) (2.16500,nan) (2.17000,nan) (2.17500,nan) (2.18000,nan) (2.18500,nan) (2.19000,nan) (2.19500,nan) (2.20000,nan) (2.20500,nan) (2.21000,nan) (2.21500,nan) (2.22000,nan) (2.22500,nan) (2.23000,nan) (2.23500,nan) (2.24000,nan) (2.24500,nan) (2.25000,nan) (2.25500,nan) (2.26000,nan) (2.26500,nan) (2.27000,nan) (2.27500,nan) (2.28000,nan) (2.28500,nan) (2.29000,nan) (2.29500,nan) (2.30000,nan) (2.30500,nan) (2.31000,nan) (2.31500,nan) (2.32000,nan) (2.32500,nan) (2.33000,nan) (2.33500,nan) (2.34000,nan) (2.34500,nan) (2.35000,nan) (2.35500,nan) (2.36000,nan) (2.36500,nan) (2.37000,nan) (2.37500,nan) (2.38000,nan) (2.38500,nan) (2.39000,nan) (2.39500,nan) (2.40000,nan) (2.40500,nan) (2.41000,nan) (2.41500,nan) (2.42000,nan) (2.42500,nan) (2.43000,nan) (2.43500,nan) (2.44000,nan) (2.44500,nan) (2.45000,nan) (2.45500,nan) (2.46000,nan) (2.46500,nan) (2.47000,nan) (2.47500,nan) (2.48000,nan) (2.48500,nan) (2.49000,nan) (2.49500,nan) (2.50000,nan) (2.50500,nan) (2.51000,nan) (2.51500,nan) (2.52000,nan) (2.52500,nan) (2.53000,nan) (2.53500,nan) (2.54000,nan) (2.54500,nan) (2.55000,nan) (2.55500,nan) (2.56000,nan) (2.56500,nan) (2.57000,nan) (2.57500,nan) (2.58000,nan) (2.58500,nan) (2.59000,nan) (2.59500,nan) (2.60000,nan) (2.60500,nan) (2.61000,nan) (2.61500,nan) (2.62000,nan) (2.62500,nan) (2.63000,nan) (2.63500,nan) (2.64000,nan) (2.64500,nan) (2.65000,nan) (2.65500,nan) (2.66000,nan) (2.66500,nan) (2.67000,nan) (2.67500,nan) (2.68000,nan) (2.68500,nan) (2.69000,nan) (2.69500,nan) (2.70000,nan) (2.70500,nan) (2.71000,nan) (2.71500,nan) (2.72000,nan) (2.72500,nan) (2.73000,nan) (2.73500,nan) (2.74000,nan) (2.74500,nan) (2.75000,nan) (2.75500,nan) (2.76000,nan) (2.76500,nan) (2.77000,nan) (2.77500,nan) (2.78000,nan) (2.78500,nan) (2.79000,nan) (2.79500,nan) (2.80000,nan) (2.80500,nan) (2.81000,nan) (2.81500,nan) (2.82000,nan) (2.82500,nan) (2.83000,nan) (2.83500,nan) (2.84000,nan) (2.84500,nan) (2.85000,nan) (2.85500,nan) (2.86000,nan) (2.86500,nan) (2.87000,nan) (2.87500,nan) (2.88000,nan) (2.88500,nan) (2.89000,nan) (2.89500,nan) (2.90000,nan) (2.90500,nan) (2.91000,nan) (2.91500,nan) (2.92000,nan) (2.92500,nan) (2.93000,nan) (2.93500,nan) (2.94000,nan) (2.94500,nan) (2.95000,nan) (2.95500,nan) (2.96000,nan) (2.96500,nan) (2.97000,nan) (2.97500,nan) (2.98000,nan) (2.98500,nan) (2.99000,nan) (2.99500,nan) (3.00000,nan) (3.00500,nan) (3.01000,nan) (3.01500,nan) (3.02000,nan) (3.02500,nan) (3.03000,nan) (3.03500,nan) (3.04000,nan) (3.04500,nan) (3.05000,nan) (3.05500,nan) (3.06000,nan) (3.06500,-1.9854413) (3.07000,-1.6127255) (3.07500,-1.4171427) (3.08000,-1.2842305) (3.08500,-1.1837467) (3.09000,-1.1031364) (3.09500,-1.0359658) (3.10000,-0.9784939) (3.10500,-0.9283524) (3.11000,-0.8839465) (3.11500,-0.8441510) (3.12000,-0.8081422) (3.12500,-0.7752990) (3.13000,-0.7451417) (3.13500,-0.7172916) (3.14000,-0.6914450) (3.14500,-0.6673543) (3.15000,-0.6448150) (3.15500,-0.6236564) (3.16000,-0.6037344) (3.16500,-0.5849265) (3.17000,-0.5671271) (3.17500,-0.5502453) (3.18000,-0.5342016) (3.18500,-0.5189266) (3.19000,-0.5043591) (3.19500,-0.4904448) (3.20000,-0.4771354) (3.20500,-0.4643878) (3.21000,-0.4521634) (3.21500,-0.4404273) (3.22000,-0.4291480) (3.22500,-0.4182969) (3.23000,-0.4078480) (3.23500,-0.3977776) (3.24000,-0.3880641) (3.24500,-0.3786874) (3.25000,-0.3696294) (3.25500,-0.3608731) (3.26000,-0.3524030) (3.26500,-0.3442047) (3.27000,-0.3362648) (3.27500,-0.3285709) (3.28000,-0.3211114) (3.28500,-0.3138756) (3.29000,-0.3068535) (3.29500,-0.3000356) (3.30000,-0.2934130) (3.30500,-0.2869777) (3.31000,-0.2807217) (3.31500,-0.2746379) (3.32000,-0.2687192) (3.32500,-0.2629594) (3.33000,-0.2573522) (3.33500,-0.2518918) (3.34000,-0.2465729) (3.34500,-0.2413902) (3.35000,-0.2363389) (3.35500,-0.2314143) (3.36000,-0.2266120) (3.36500,-0.2219277) (3.37000,-0.2173576) (3.37500,-0.2128978) (3.38000,-0.2085447) (3.38500,-0.2042948) (3.39000,-0.2001449) (3.39500,-0.1960918) (3.40000,-0.1921324) (3.40500,-0.1882641) (3.41000,-0.1844839) (3.41500,-0.1807892) (3.42000,-0.1771776) (3.42500,-0.1736466) (3.43000,-0.1701939) (3.43500,-0.1668173) (3.44000,-0.1635146) (3.44500,-0.1602837) (3.45000,-0.1571228) (3.45500,-0.1540299) (3.46000,-0.1510031) (3.46500,-0.1480407) (3.47000,-0.1451411) (3.47500,-0.1423025) (3.48000,-0.1395235) (3.48500,-0.1368024) (3.49000,-0.1341378) (3.49500,-0.1315283) (3.50000,-0.1289726) (3.50500,-0.1264692) (3.51000,-0.1240169) (3.51500,-0.1216145) (3.52000,-0.1192608) (3.52500,-0.1169547) (3.53000,-0.1146949) (3.53500,-0.1124804) (3.54000,-0.1103102) (3.54500,-0.1081832) (3.55000,-0.1060984) (3.55500,-0.1040550) (3.56000,-0.1020518) (3.56500,-0.1000881) (3.57000,-0.0981630) (3.57500,-0.0962756) (3.58000,-0.0944251) (3.58500,-0.0926106) (3.59000,-0.0908315) (3.59500,-0.0890868) (3.60000,-0.0873760) (3.60500,-0.0856982) (3.61000,-0.0840528) (3.61500,-0.0824391) (3.62000,-0.0808564) (3.62500,-0.0793040) (3.63000,-0.0777815) (3.63500,-0.0762880) (3.64000,-0.0748230) (3.64500,-0.0733861) (3.65000,-0.0719764) (3.65500,-0.0705936) (3.66000,-0.0692371) (3.66500,-0.0679063) (3.67000,-0.0666008) (3.67500,-0.0653200) (3.68000,-0.0640635) (3.68500,-0.0628307) (3.69000,-0.0616212) (3.69500,-0.0604345) (3.70000,-0.0592702) (3.70500,-0.0581279) (3.71000,-0.0570071) (3.71500,-0.0559074) (3.72000,-0.0548284) (3.72500,-0.0537697) (3.73000,-0.0527309) (3.73500,-0.0517117) (3.74000,-0.0507116) (3.74500,-0.0497303) (3.75000,-0.0487675) (3.75500,-0.0478227) (3.76000,-0.0468957) (3.76500,-0.0459861) (3.77000,-0.0450935) (3.77500,-0.0442178) (3.78000,-0.0433584) (3.78500,-0.0425152) (3.79000,-0.0416879) (3.79500,-0.0408761) (3.80000,-0.0400795) (3.80500,-0.0392979) (3.81000,-0.0385310) (3.81500,-0.0377786) (3.82000,-0.0370402) (3.82500,-0.0363158) (3.83000,-0.0356050) (3.83500,-0.0349075) (3.84000,-0.0342232) (3.84500,-0.0335518) (3.85000,-0.0328931) (3.85500,-0.0322468) (3.86000,-0.0316126) (3.86500,-0.0309905) (3.87000,-0.0303801) (3.87500,-0.0297812) (3.88000,-0.0291936) (3.88500,-0.0286172) (3.89000,-0.0280517) (3.89500,-0.0274969) (3.90000,-0.0269526) (3.90500,-0.0264186) (3.91000,-0.0258948) (3.91500,-0.0253809) (3.92000,-0.0248768) (3.92500,-0.0243823) (3.93000,-0.0238972) (3.93500,-0.0234213) (3.94000,-0.0229545) (3.94500,-0.0224966) (3.95000,-0.0220475) (3.95500,-0.0216070) (3.96000,-0.0211748) (3.96500,-0.0207510) (3.97000,-0.0203353) (3.97500,-0.0199275) (3.98000,-0.0195276) (3.98500,-0.0191353) (3.99000,-0.0187506) (3.99500,-0.0183733) (4.00000,-0.0180033)};\addlegendentry{L}
\addplot[thick,dashed,gray] coordinates {(-2.00000,nan) (-1.99500,nan) (-1.99000,nan) (-1.98500,nan) (-1.98000,nan) (-1.97500,nan) (-1.97000,nan) (-1.96500,nan) (-1.96000,nan) (-1.95500,nan) (-1.95000,nan) (-1.94500,nan) (-1.94000,nan) (-1.93500,nan) (-1.93000,nan) (-1.92500,nan) (-1.92000,nan) (-1.91500,nan) (-1.91000,nan) (-1.90500,nan) (-1.90000,nan) (-1.89500,nan) (-1.89000,nan) (-1.88500,nan) (-1.88000,nan) (-1.87500,nan) (-1.87000,nan) (-1.86500,nan) (-1.86000,nan) (-1.85500,nan) (-1.85000,nan) (-1.84500,nan) (-1.84000,nan) (-1.83500,nan) (-1.83000,nan) (-1.82500,nan) (-1.82000,nan) (-1.81500,nan) (-1.81000,nan) (-1.80500,nan) (-1.80000,nan) (-1.79500,nan) (-1.79000,nan) (-1.78500,nan) (-1.78000,nan) (-1.77500,nan) (-1.77000,nan) (-1.76500,nan) (-1.76000,nan) (-1.75500,nan) (-1.75000,nan) (-1.74500,nan) (-1.74000,nan) (-1.73500,nan) (-1.73000,nan) (-1.72500,nan) (-1.72000,nan) (-1.71500,nan) (-1.71000,nan) (-1.70500,nan) (-1.70000,nan) (-1.69500,nan) (-1.69000,nan) (-1.68500,nan) (-1.68000,nan) (-1.67500,nan) (-1.67000,nan) (-1.66500,nan) (-1.66000,nan) (-1.65500,nan) (-1.65000,nan) (-1.64500,nan) (-1.64000,nan) (-1.63500,nan) (-1.63000,nan) (-1.62500,nan) (-1.62000,nan) (-1.61500,nan) (-1.61000,nan) (-1.60500,nan) (-1.60000,nan) (-1.59500,nan) (-1.59000,nan) (-1.58500,nan) (-1.58000,nan) (-1.57500,nan) (-1.57000,nan) (-1.56500,nan) (-1.56000,nan) (-1.55500,nan) (-1.55000,nan) (-1.54500,nan) (-1.54000,nan) (-1.53500,nan) (-1.53000,nan) (-1.52500,nan) (-1.52000,nan) (-1.51500,nan) (-1.51000,nan) (-1.50500,nan) (-1.50000,nan) (-1.49500,nan) (-1.49000,nan) (-1.48500,nan) (-1.48000,nan) (-1.47500,nan) (-1.47000,nan) (-1.46500,nan) (-1.46000,nan) (-1.45500,nan) (-1.45000,nan) (-1.44500,nan) (-1.44000,nan) (-1.43500,nan) (-1.43000,nan) (-1.42500,nan) (-1.42000,nan) (-1.41500,nan) (-1.41000,nan) (-1.40500,nan) (-1.40000,nan) (-1.39500,nan) (-1.39000,nan) (-1.38500,nan) (-1.38000,nan) (-1.37500,nan) (-1.37000,nan) (-1.36500,nan) (-1.36000,nan) (-1.35500,nan) (-1.35000,nan) (-1.34500,nan) (-1.34000,nan) (-1.33500,nan) (-1.33000,nan) (-1.32500,nan) (-1.32000,nan) (-1.31500,nan) (-1.31000,nan) (-1.30500,nan) (-1.30000,nan) (-1.29500,nan) (-1.29000,nan) (-1.28500,nan) (-1.28000,nan) (-1.27500,nan) (-1.27000,nan) (-1.26500,nan) (-1.26000,nan) (-1.25500,nan) (-1.25000,nan) (-1.24500,nan) (-1.24000,nan) (-1.23500,nan) (-1.23000,nan) (-1.22500,nan) (-1.22000,nan) (-1.21500,nan) (-1.21000,nan) (-1.20500,nan) (-1.20000,nan) (-1.19500,nan) (-1.19000,nan) (-1.18500,nan) (-1.18000,nan) (-1.17500,nan) (-1.17000,nan) (-1.16500,nan) (-1.16000,nan) (-1.15500,nan) (-1.15000,nan) (-1.14500,nan) (-1.14000,nan) (-1.13500,nan) (-1.13000,nan) (-1.12500,nan) (-1.12000,nan) (-1.11500,nan) (-1.11000,nan) (-1.10500,nan) (-1.10000,nan) (-1.09500,nan) (-1.09000,nan) (-1.08500,nan) (-1.08000,nan) (-1.07500,nan) (-1.07000,nan) (-1.06500,nan) (-1.06000,nan) (-1.05500,nan) (-1.05000,nan) (-1.04500,nan) (-1.04000,nan) (-1.03500,nan) (-1.03000,nan) (-1.02500,nan) (-1.02000,nan) (-1.01500,nan) (-1.01000,nan) (-1.00500,nan) (-1.00000,nan) (-0.99500,nan) (-0.99000,nan) (-0.98500,nan) (-0.98000,nan) (-0.97500,nan) (-0.97000,nan) (-0.96500,nan) (-0.96000,nan) (-0.95500,nan) (-0.95000,nan) (-0.94500,nan) (-0.94000,nan) (-0.93500,nan) (-0.93000,nan) (-0.92500,nan) (-0.92000,nan) (-0.91500,nan) (-0.91000,nan) (-0.90500,nan) (-0.90000,nan) (-0.89500,nan) (-0.89000,nan) (-0.88500,nan) (-0.88000,nan) (-0.87500,nan) (-0.87000,nan) (-0.86500,nan) (-0.86000,nan) (-0.85500,nan) (-0.85000,nan) (-0.84500,nan) (-0.84000,nan) (-0.83500,nan) (-0.83000,nan) (-0.82500,nan) (-0.82000,nan) (-0.81500,nan) (-0.81000,nan) (-0.80500,nan) (-0.80000,nan) (-0.79500,nan) (-0.79000,nan) (-0.78500,nan) (-0.78000,nan) (-0.77500,nan) (-0.77000,nan) (-0.76500,nan) (-0.76000,nan) (-0.75500,nan) (-0.75000,nan) (-0.74500,nan) (-0.74000,nan) (-0.73500,nan) (-0.73000,nan) (-0.72500,nan) (-0.72000,nan) (-0.71500,nan) (-0.71000,nan) (-0.70500,nan) (-0.70000,nan) (-0.69500,nan) (-0.69000,nan) (-0.68500,nan) (-0.68000,nan) (-0.67500,nan) (-0.67000,nan) (-0.66500,nan) (-0.66000,nan) (-0.65500,nan) (-0.65000,nan) (-0.64500,nan) (-0.64000,nan) (-0.63500,nan) (-0.63000,nan) (-0.62500,nan) (-0.62000,nan) (-0.61500,nan) (-0.61000,nan) (-0.60500,nan) (-0.60000,nan) (-0.59500,nan) (-0.59000,nan) (-0.58500,nan) (-0.58000,nan) (-0.57500,nan) (-0.57000,nan) (-0.56500,nan) (-0.56000,nan) (-0.55500,nan) (-0.55000,nan) (-0.54500,nan) (-0.54000,nan) (-0.53500,nan) (-0.53000,nan) (-0.52500,nan) (-0.52000,nan) (-0.51500,nan) (-0.51000,nan) (-0.50500,nan) (-0.50000,nan) (-0.49500,nan) (-0.49000,nan) (-0.48500,nan) (-0.48000,nan) (-0.47500,nan) (-0.47000,nan) (-0.46500,nan) (-0.46000,nan) (-0.45500,nan) (-0.45000,nan) (-0.44500,nan) (-0.44000,nan) (-0.43500,nan) (-0.43000,nan) (-0.42500,nan) (-0.42000,nan) (-0.41500,nan) (-0.41000,nan) (-0.40500,nan) (-0.40000,nan) (-0.39500,nan) (-0.39000,nan) (-0.38500,nan) (-0.38000,nan) (-0.37500,nan) (-0.37000,nan) (-0.36500,nan) (-0.36000,nan) (-0.35500,nan) (-0.35000,nan) (-0.34500,nan) (-0.34000,nan) (-0.33500,nan) (-0.33000,nan) (-0.32500,nan) (-0.32000,nan) (-0.31500,nan) (-0.31000,nan) (-0.30500,nan) (-0.30000,nan) (-0.29500,nan) (-0.29000,nan) (-0.28500,nan) (-0.28000,nan) (-0.27500,nan) (-0.27000,nan) (-0.26500,nan) (-0.26000,nan) (-0.25500,nan) (-0.25000,nan) (-0.24500,nan) (-0.24000,nan) (-0.23500,nan) (-0.23000,nan) (-0.22500,nan) (-0.22000,nan) (-0.21500,nan) (-0.21000,nan) (-0.20500,nan) (-0.20000,nan) (-0.19500,nan) (-0.19000,nan) (-0.18500,nan) (-0.18000,nan) (-0.17500,nan) (-0.17000,nan) (-0.16500,nan) (-0.16000,nan) (-0.15500,nan) (-0.15000,nan) (-0.14500,nan) (-0.14000,nan) (-0.13500,nan) (-0.13000,nan) (-0.12500,nan) (-0.12000,nan) (-0.11500,nan) (-0.11000,nan) (-0.10500,nan) (-0.10000,nan) (-0.09500,nan) (-0.09000,nan) (-0.08500,nan) (-0.08000,nan) (-0.07500,nan) (-0.07000,nan) (-0.06500,nan) (-0.06000,nan) (-0.05500,nan) (-0.05000,nan) (-0.04500,nan) (-0.04000,nan) (-0.03500,nan) (-0.03000,nan) (-0.02500,nan) (-0.02000,nan) (-0.01500,nan) (-0.01000,nan) (-0.00500,nan) (0.00000,nan) (0.00500,nan) (0.01000,nan) (0.01500,nan) (0.02000,nan) (0.02500,nan) (0.03000,nan) (0.03500,nan) (0.04000,nan) (0.04500,nan) (0.05000,nan) (0.05500,nan) (0.06000,nan) (0.06500,nan) (0.07000,nan) (0.07500,nan) (0.08000,nan) (0.08500,nan) (0.09000,nan) (0.09500,nan) (0.10000,nan) (0.10500,nan) (0.11000,nan) (0.11500,nan) (0.12000,nan) (0.12500,nan) (0.13000,nan) (0.13500,nan) (0.14000,nan) (0.14500,nan) (0.15000,nan) (0.15500,nan) (0.16000,nan) (0.16500,nan) (0.17000,nan) (0.17500,nan) (0.18000,nan) (0.18500,nan) (0.19000,nan) (0.19500,nan) (0.20000,nan) (0.20500,nan) (0.21000,nan) (0.21500,nan) (0.22000,nan) (0.22500,nan) (0.23000,nan) (0.23500,nan) (0.24000,nan) (0.24500,nan) (0.25000,nan) (0.25500,nan) (0.26000,nan) (0.26500,nan) (0.27000,nan) (0.27500,nan) (0.28000,nan) (0.28500,nan) (0.29000,nan) (0.29500,nan) (0.30000,nan) (0.30500,nan) (0.31000,nan) (0.31500,nan) (0.32000,nan) (0.32500,nan) (0.33000,nan) (0.33500,nan) (0.34000,nan) (0.34500,nan) (0.35000,nan) (0.35500,nan) (0.36000,nan) (0.36500,nan) (0.37000,nan) (0.37500,nan) (0.38000,nan) (0.38500,nan) (0.39000,nan) (0.39500,nan) (0.40000,nan) (0.40500,nan) (0.41000,nan) (0.41500,nan) (0.42000,nan) (0.42500,nan) (0.43000,nan) (0.43500,nan) (0.44000,nan) (0.44500,nan) (0.45000,nan) (0.45500,nan) (0.46000,nan) (0.46500,nan) (0.47000,nan) (0.47500,nan) (0.48000,nan) (0.48500,nan) (0.49000,nan) (0.49500,nan) (0.50000,nan) (0.50500,nan) (0.51000,nan) (0.51500,nan) (0.52000,nan) (0.52500,nan) (0.53000,nan) (0.53500,nan) (0.54000,nan) (0.54500,nan) (0.55000,nan) (0.55500,nan) (0.56000,nan) (0.56500,nan) (0.57000,nan) (0.57500,nan) (0.58000,nan) (0.58500,nan) (0.59000,nan) (0.59500,nan) (0.60000,nan) (0.60500,nan) (0.61000,nan) (0.61500,nan) (0.62000,nan) (0.62500,nan) (0.63000,nan) (0.63500,nan) (0.64000,nan) (0.64500,nan) (0.65000,nan) (0.65500,nan) (0.66000,nan) (0.66500,nan) (0.67000,nan) (0.67500,nan) (0.68000,nan) (0.68500,nan) (0.69000,nan) (0.69500,nan) (0.70000,nan) (0.70500,nan) (0.71000,nan) (0.71500,nan) (0.72000,nan) (0.72500,nan) (0.73000,nan) (0.73500,nan) (0.74000,nan) (0.74500,nan) (0.75000,nan) (0.75500,nan) (0.76000,nan) (0.76500,nan) (0.77000,nan) (0.77500,nan) (0.78000,nan) (0.78500,nan) (0.79000,nan) (0.79500,nan) (0.80000,nan) (0.80500,nan) (0.81000,nan) (0.81500,nan) (0.82000,nan) (0.82500,nan) (0.83000,nan) (0.83500,nan) (0.84000,nan) (0.84500,nan) (0.85000,nan) (0.85500,nan) (0.86000,nan) (0.86500,nan) (0.87000,nan) (0.87500,nan) (0.88000,nan) (0.88500,nan) (0.89000,nan) (0.89500,nan) (0.90000,nan) (0.90500,nan) (0.91000,nan) (0.91500,nan) (0.92000,nan) (0.92500,nan) (0.93000,nan) (0.93500,nan) (0.94000,nan) (0.94500,nan) (0.95000,nan) (0.95500,nan) (0.96000,nan) (0.96500,nan) (0.97000,nan) (0.97500,nan) (0.98000,nan) (0.98500,nan) (0.99000,nan) (0.99500,nan) (1.00000,nan) (1.00500,nan) (1.01000,nan) (1.01500,nan) (1.02000,nan) (1.02500,nan) (1.03000,nan) (1.03500,nan) (1.04000,nan) (1.04500,nan) (1.05000,nan) (1.05500,nan) (1.06000,nan) (1.06500,nan) (1.07000,nan) (1.07500,nan) (1.08000,nan) (1.08500,nan) (1.09000,nan) (1.09500,nan) (1.10000,nan) (1.10500,nan) (1.11000,nan) (1.11500,nan) (1.12000,nan) (1.12500,nan) (1.13000,nan) (1.13500,nan) (1.14000,nan) (1.14500,nan) (1.15000,nan) (1.15500,nan) (1.16000,nan) (1.16500,nan) (1.17000,nan) (1.17500,nan) (1.18000,nan) (1.18500,nan) (1.19000,nan) (1.19500,nan) (1.20000,nan) (1.20500,nan) (1.21000,nan) (1.21500,nan) (1.22000,nan) (1.22500,nan) (1.23000,nan) (1.23500,nan) (1.24000,nan) (1.24500,nan) (1.25000,nan) (1.25500,nan) (1.26000,nan) (1.26500,nan) (1.27000,nan) (1.27500,nan) (1.28000,nan) (1.28500,nan) (1.29000,nan) (1.29500,nan) (1.30000,nan) (1.30500,nan) (1.31000,nan) (1.31500,nan) (1.32000,nan) (1.32500,nan) (1.33000,nan) (1.33500,nan) (1.34000,nan) (1.34500,nan) (1.35000,nan) (1.35500,nan) (1.36000,nan) (1.36500,nan) (1.37000,nan) (1.37500,nan) (1.38000,nan) (1.38500,nan) (1.39000,nan) (1.39500,nan) (1.40000,nan) (1.40500,nan) (1.41000,nan) (1.41500,nan) (1.42000,nan) (1.42500,nan) (1.43000,nan) (1.43500,nan) (1.44000,nan) (1.44500,nan) (1.45000,nan) (1.45500,nan) (1.46000,nan) (1.46500,nan) (1.47000,nan) (1.47500,nan) (1.48000,nan) (1.48500,nan) (1.49000,nan) (1.49500,nan) (1.50000,nan) (1.50500,nan) (1.51000,nan) (1.51500,nan) (1.52000,nan) (1.52500,nan) (1.53000,nan) (1.53500,nan) (1.54000,nan) (1.54500,nan) (1.55000,nan) (1.55500,nan) (1.56000,nan) (1.56500,nan) (1.57000,nan) (1.57500,nan) (1.58000,nan) (1.58500,nan) (1.59000,nan) (1.59500,nan) (1.60000,nan) (1.60500,nan) (1.61000,nan) (1.61500,nan) (1.62000,nan) (1.62500,nan) (1.63000,nan) (1.63500,nan) (1.64000,nan) (1.64500,nan) (1.65000,nan) (1.65500,nan) (1.66000,nan) (1.66500,nan) (1.67000,nan) (1.67500,nan) (1.68000,nan) (1.68500,nan) (1.69000,nan) (1.69500,nan) (1.70000,nan) (1.70500,nan) (1.71000,nan) (1.71500,nan) (1.72000,nan) (1.72500,nan) (1.73000,nan) (1.73500,nan) (1.74000,nan) (1.74500,nan) (1.75000,nan) (1.75500,nan) (1.76000,nan) (1.76500,nan) (1.77000,nan) (1.77500,nan) (1.78000,nan) (1.78500,nan) (1.79000,nan) (1.79500,nan) (1.80000,nan) (1.80500,nan) (1.81000,nan) (1.81500,nan) (1.82000,nan) (1.82500,nan) (1.83000,nan) (1.83500,nan) (1.84000,nan) (1.84500,nan) (1.85000,nan) (1.85500,nan) (1.86000,nan) (1.86500,nan) (1.87000,nan) (1.87500,nan) (1.88000,nan) (1.88500,nan) (1.89000,nan) (1.89500,nan) (1.90000,nan) (1.90500,nan) (1.91000,nan) (1.91500,nan) (1.92000,nan) (1.92500,nan) (1.93000,nan) (1.93500,nan) (1.94000,nan) (1.94500,nan) (1.95000,nan) (1.95500,nan) (1.96000,nan) (1.96500,nan) (1.97000,nan) (1.97500,nan) (1.98000,nan) (1.98500,nan) (1.99000,nan) (1.99500,nan) (2.00000,nan) (2.00500,nan) (2.01000,nan) (2.01500,nan) (2.02000,nan) (2.02500,nan) (2.03000,nan) (2.03500,nan) (2.04000,nan) (2.04500,nan) (2.05000,nan) (2.05500,nan) (2.06000,nan) (2.06500,nan) (2.07000,nan) (2.07500,nan) (2.08000,nan) (2.08500,nan) (2.09000,nan) (2.09500,nan) (2.10000,nan) (2.10500,nan) (2.11000,nan) (2.11500,nan) (2.12000,nan) (2.12500,nan) (2.13000,nan) (2.13500,nan) (2.14000,nan) (2.14500,nan) (2.15000,nan) (2.15500,nan) (2.16000,nan) (2.16500,nan) (2.17000,nan) (2.17500,nan) (2.18000,nan) (2.18500,nan) (2.19000,nan) (2.19500,nan) (2.20000,nan) (2.20500,nan) (2.21000,nan) (2.21500,nan) (2.22000,nan) (2.22500,nan) (2.23000,nan) (2.23500,nan) (2.24000,nan) (2.24500,nan) (2.25000,nan) (2.25500,nan) (2.26000,nan) (2.26500,nan) (2.27000,nan) (2.27500,nan) (2.28000,nan) (2.28500,nan) (2.29000,nan) (2.29500,nan) (2.30000,nan) (2.30500,nan) (2.31000,nan) (2.31500,nan) (2.32000,nan) (2.32500,nan) (2.33000,nan) (2.33500,nan) (2.34000,nan) (2.34500,nan) (2.35000,nan) (2.35500,nan) (2.36000,nan) (2.36500,nan) (2.37000,nan) (2.37500,nan) (2.38000,nan) (2.38500,nan) (2.39000,nan) (2.39500,nan) (2.40000,nan) (2.40500,nan) (2.41000,nan) (2.41500,nan) (2.42000,nan) (2.42500,nan) (2.43000,nan) (2.43500,nan) (2.44000,nan) (2.44500,nan) (2.45000,nan) (2.45500,nan) (2.46000,nan) (2.46500,nan) (2.47000,nan) (2.47500,nan) (2.48000,nan) (2.48500,nan) (2.49000,nan) (2.49500,nan) (2.50000,nan) (2.50500,nan) (2.51000,nan) (2.51500,nan) (2.52000,nan) (2.52500,nan) (2.53000,nan) (2.53500,nan) (2.54000,nan) (2.54500,nan) (2.55000,nan) (2.55500,nan) (2.56000,nan) (2.56500,nan) (2.57000,nan) (2.57500,nan) (2.58000,nan) (2.58500,nan) (2.59000,nan) (2.59500,nan) (2.60000,nan) (2.60500,nan) (2.61000,nan) (2.61500,nan) (2.62000,nan) (2.62500,nan) (2.63000,nan) (2.63500,nan) (2.64000,nan) (2.64500,nan) (2.65000,nan) (2.65500,nan) (2.66000,nan) (2.66500,nan) (2.67000,nan) (2.67500,nan) (2.68000,nan) (2.68500,nan) (2.69000,nan) (2.69500,nan) (2.70000,nan) (2.70500,nan) (2.71000,nan) (2.71500,nan) (2.72000,nan) (2.72500,nan) (2.73000,nan) (2.73500,nan) (2.74000,nan) (2.74500,nan) (2.75000,nan) (2.75500,nan) (2.76000,nan) (2.76500,nan) (2.77000,nan) (2.77500,nan) (2.78000,nan) (2.78500,nan) (2.79000,nan) (2.79500,nan) (2.80000,nan) (2.80500,nan) (2.81000,nan) (2.81500,nan) (2.82000,nan) (2.82500,nan) (2.83000,nan) (2.83500,nan) (2.84000,nan) (2.84500,nan) (2.85000,nan) (2.85500,nan) (2.86000,nan) (2.86500,nan) (2.87000,nan) (2.87500,nan) (2.88000,nan) (2.88500,nan) (2.89000,nan) (2.89500,nan) (2.90000,nan) (2.90500,nan) (2.91000,nan) (2.91500,nan) (2.92000,nan) (2.92500,nan) (2.93000,nan) (2.93500,nan) (2.94000,nan) (2.94500,nan) (2.95000,nan) (2.95500,nan) (2.96000,nan) (2.96500,nan) (2.97000,nan) (2.97500,nan) (2.98000,nan) (2.98500,nan) (2.99000,nan) (2.99500,nan) (3.00000,nan) (3.00500,nan) (3.01000,nan) (3.01500,nan) (3.02000,nan) (3.02500,nan) (3.03000,nan) (3.03500,nan) (3.04000,nan) (3.04500,nan) (3.05000,nan) (3.05500,nan) (3.06000,nan) (3.06500,nan) (3.07000,nan) (3.07500,nan) (3.08000,nan) (3.08500,nan) (3.09000,nan) (3.09500,nan) (3.10000,nan) (3.10500,nan) (3.11000,nan) (3.11500,nan) (3.12000,nan) (3.12500,nan) (3.13000,nan) (3.13500,nan) (3.14000,nan) (3.14500,nan) (3.15000,nan) (3.15500,nan) (3.16000,nan) (3.16500,nan) (3.17000,nan) (3.17500,nan) (3.18000,nan) (3.18500,nan) (3.19000,nan) (3.19500,nan) (3.20000,nan) (3.20500,nan) (3.21000,nan) (3.21500,nan) (3.22000,nan) (3.22500,nan) (3.23000,nan) (3.23500,nan) (3.24000,nan) (3.24500,nan) (3.25000,nan) (3.25500,nan) (3.26000,nan) (3.26500,nan) (3.27000,nan) (3.27500,nan) (3.28000,nan) (3.28500,nan) (3.29000,nan) (3.29500,nan) (3.30000,nan) (3.30500,nan) (3.31000,nan) (3.31500,nan) (3.32000,nan) (3.32500,nan) (3.33000,nan) (3.33500,nan) (3.34000,nan) (3.34500,nan) (3.35000,nan) (3.35500,nan) (3.36000,nan) (3.36500,nan) (3.37000,nan) (3.37500,-1.8920908) (3.38000,-1.4523290) (3.38500,-1.2422168) (3.39000,-1.1039668) (3.39500,-1.0013845) (3.40000,-0.9201842) (3.40500,-0.8532436) (3.41000,-0.7964937) (3.41500,-0.7473898) (3.42000,-0.7042340) (3.42500,-0.6658368) (3.43000,-0.6313323) (3.43500,-0.6000705) (3.44000,-0.5715512) (3.44500,-0.5453811) (3.45000,-0.5212452) (3.45500,-0.4988877) (3.46000,-0.4780976) (3.46500,-0.4586991) (3.47000,-0.4405441) (3.47500,-0.4235067) (3.48000,-0.4074789) (3.48500,-0.3923674) (3.49000,-0.3780911) (3.49500,-0.3645790) (3.50000,-0.3517686) (3.50500,-0.3396046) (3.51000,-0.3280377) (3.51500,-0.3170242) (3.52000,-0.3065247) (3.52500,-0.2965037) (3.53000,-0.2869293) (3.53500,-0.2777724) (3.54000,-0.2690067) (3.54500,-0.2606081) (3.55000,-0.2525547) (3.55500,-0.2448262) (3.56000,-0.2374043) (3.56500,-0.2302719) (3.57000,-0.2234131) (3.57500,-0.2168136) (3.58000,-0.2104597) (3.58500,-0.2043390) (3.59000,-0.1984398) (3.59500,-0.1927513) (3.60000,-0.1872634) (3.60500,-0.1819665) (3.61000,-0.1768519) (3.61500,-0.1719113) (3.62000,-0.1671368) (3.62500,-0.1625211) (3.63000,-0.1580574) (3.63500,-0.1537392) (3.64000,-0.1495604) (3.64500,-0.1455152) (3.65000,-0.1415981) (3.65500,-0.1378041) (3.66000,-0.1341281) (3.66500,-0.1305656) (3.67000,-0.1271123) (3.67500,-0.1237638) (3.68000,-0.1205164) (3.68500,-0.1173661) (3.69000,-0.1143095) (3.69500,-0.1113431) (3.70000,-0.1084636) (3.70500,-0.1056680) (3.71000,-0.1029534) (3.71500,-0.1003168) (3.72000,-0.0977556) (3.72500,-0.0952672) (3.73000,-0.0928492) (3.73500,-0.0904991) (3.74000,-0.0882148) (3.74500,-0.0859940) (3.75000,-0.0838347) (3.75500,-0.0817348) (3.76000,-0.0796925) (3.76500,-0.0777059) (3.77000,-0.0757733) (3.77500,-0.0738930) (3.78000,-0.0720633) (3.78500,-0.0702826) (3.79000,-0.0685495) (3.79500,-0.0668625) (3.80000,-0.0652203) (3.80500,-0.0636213) (3.81000,-0.0620645) (3.81500,-0.0605484) (3.82000,-0.0590720) (3.82500,-0.0576340) (3.83000,-0.0562333) (3.83500,-0.0548689) (3.84000,-0.0535396) (3.84500,-0.0522445) (3.85000,-0.0509826) (3.85500,-0.0497530) (3.86000,-0.0485547) (3.86500,-0.0473869) (3.87000,-0.0462487) (3.87500,-0.0451393) (3.88000,-0.0440578) (3.88500,-0.0430036) (3.89000,-0.0419758) (3.89500,-0.0409738) (3.90000,-0.0399968) (3.90500,-0.0390442) (3.91000,-0.0381152) (3.91500,-0.0372094) (3.92000,-0.0363259) (3.92500,-0.0354643) (3.93000,-0.0346240) (3.93500,-0.0338044) (3.94000,-0.0330049) (3.94500,-0.0322250) (3.95000,-0.0314643) (3.95500,-0.0307221) (3.96000,-0.0299981) (3.96500,-0.0292917) (3.97000,-0.0286024) (3.97500,-0.0279300) (3.98000,-0.0272738) (3.98500,-0.0266336) (3.99000,-0.0260088) (3.99500,-0.0253991) (4.00000,-0.0248041)};\addlegendentry{O}
\addplot[thick,dotted,blue] coordinates {(-2.00000,nan) (-1.99500,nan) (-1.99000,nan) (-1.98500,nan) (-1.98000,nan) (-1.97500,nan) (-1.97000,nan) (-1.96500,nan) (-1.96000,nan) (-1.95500,nan) (-1.95000,nan) (-1.94500,nan) (-1.94000,nan) (-1.93500,nan) (-1.93000,nan) (-1.92500,nan) (-1.92000,nan) (-1.91500,nan) (-1.91000,nan) (-1.90500,nan) (-1.90000,nan) (-1.89500,nan) (-1.89000,nan) (-1.88500,nan) (-1.88000,nan) (-1.87500,nan) (-1.87000,nan) (-1.86500,nan) (-1.86000,nan) (-1.85500,nan) (-1.85000,nan) (-1.84500,nan) (-1.84000,nan) (-1.83500,nan) (-1.83000,nan) (-1.82500,nan) (-1.82000,nan) (-1.81500,nan) (-1.81000,nan) (-1.80500,nan) (-1.80000,nan) (-1.79500,nan) (-1.79000,nan) (-1.78500,nan) (-1.78000,nan) (-1.77500,nan) (-1.77000,nan) (-1.76500,nan) (-1.76000,nan) (-1.75500,nan) (-1.75000,nan) (-1.74500,nan) (-1.74000,nan) (-1.73500,nan) (-1.73000,nan) (-1.72500,nan) (-1.72000,nan) (-1.71500,nan) (-1.71000,nan) (-1.70500,nan) (-1.70000,nan) (-1.69500,nan) (-1.69000,nan) (-1.68500,nan) (-1.68000,nan) (-1.67500,nan) (-1.67000,nan) (-1.66500,nan) (-1.66000,nan) (-1.65500,nan) (-1.65000,nan) (-1.64500,nan) (-1.64000,nan) (-1.63500,nan) (-1.63000,nan) (-1.62500,nan) (-1.62000,nan) (-1.61500,nan) (-1.61000,nan) (-1.60500,nan) (-1.60000,nan) (-1.59500,nan) (-1.59000,nan) (-1.58500,nan) (-1.58000,nan) (-1.57500,nan) (-1.57000,nan) (-1.56500,nan) (-1.56000,nan) (-1.55500,nan) (-1.55000,nan) (-1.54500,nan) (-1.54000,nan) (-1.53500,nan) (-1.53000,nan) (-1.52500,nan) (-1.52000,nan) (-1.51500,nan) (-1.51000,nan) (-1.50500,nan) (-1.50000,nan) (-1.49500,nan) (-1.49000,nan) (-1.48500,nan) (-1.48000,nan) (-1.47500,nan) (-1.47000,nan) (-1.46500,nan) (-1.46000,nan) (-1.45500,nan) (-1.45000,nan) (-1.44500,nan) (-1.44000,nan) (-1.43500,nan) (-1.43000,nan) (-1.42500,nan) (-1.42000,nan) (-1.41500,nan) (-1.41000,nan) (-1.40500,nan) (-1.40000,nan) (-1.39500,nan) (-1.39000,nan) (-1.38500,nan) (-1.38000,nan) (-1.37500,nan) (-1.37000,nan) (-1.36500,nan) (-1.36000,nan) (-1.35500,nan) (-1.35000,nan) (-1.34500,nan) (-1.34000,nan) (-1.33500,nan) (-1.33000,nan) (-1.32500,nan) (-1.32000,nan) (-1.31500,nan) (-1.31000,nan) (-1.30500,nan) (-1.30000,nan) (-1.29500,nan) (-1.29000,nan) (-1.28500,nan) (-1.28000,nan) (-1.27500,nan) (-1.27000,nan) (-1.26500,nan) (-1.26000,nan) (-1.25500,nan) (-1.25000,nan) (-1.24500,nan) (-1.24000,nan) (-1.23500,nan) (-1.23000,nan) (-1.22500,nan) (-1.22000,nan) (-1.21500,nan) (-1.21000,nan) (-1.20500,nan) (-1.20000,nan) (-1.19500,nan) (-1.19000,nan) (-1.18500,nan) (-1.18000,nan) (-1.17500,nan) (-1.17000,nan) (-1.16500,nan) (-1.16000,nan) (-1.15500,nan) (-1.15000,nan) (-1.14500,nan) (-1.14000,nan) (-1.13500,nan) (-1.13000,nan) (-1.12500,nan) (-1.12000,nan) (-1.11500,nan) (-1.11000,nan) (-1.10500,nan) (-1.10000,nan) (-1.09500,nan) (-1.09000,nan) (-1.08500,nan) (-1.08000,nan) (-1.07500,nan) (-1.07000,nan) (-1.06500,nan) (-1.06000,nan) (-1.05500,nan) (-1.05000,nan) (-1.04500,nan) (-1.04000,nan) (-1.03500,nan) (-1.03000,nan) (-1.02500,nan) (-1.02000,nan) (-1.01500,nan) (-1.01000,nan) (-1.00500,nan) (-1.00000,nan) (-0.99500,nan) (-0.99000,nan) (-0.98500,nan) (-0.98000,nan) (-0.97500,nan) (-0.97000,nan) (-0.96500,nan) (-0.96000,nan) (-0.95500,nan) (-0.95000,nan) (-0.94500,nan) (-0.94000,nan) (-0.93500,nan) (-0.93000,nan) (-0.92500,nan) (-0.92000,nan) (-0.91500,nan) (-0.91000,nan) (-0.90500,nan) (-0.90000,nan) (-0.89500,nan) (-0.89000,nan) (-0.88500,nan) (-0.88000,nan) (-0.87500,nan) (-0.87000,nan) (-0.86500,nan) (-0.86000,nan) (-0.85500,nan) (-0.85000,nan) (-0.84500,nan) (-0.84000,nan) (-0.83500,nan) (-0.83000,nan) (-0.82500,nan) (-0.82000,nan) (-0.81500,nan) (-0.81000,nan) (-0.80500,nan) (-0.80000,nan) (-0.79500,nan) (-0.79000,nan) (-0.78500,nan) (-0.78000,nan) (-0.77500,nan) (-0.77000,nan) (-0.76500,nan) (-0.76000,nan) (-0.75500,nan) (-0.75000,nan) (-0.74500,nan) (-0.74000,nan) (-0.73500,nan) (-0.73000,nan) (-0.72500,nan) (-0.72000,nan) (-0.71500,nan) (-0.71000,nan) (-0.70500,nan) (-0.70000,nan) (-0.69500,nan) (-0.69000,nan) (-0.68500,nan) (-0.68000,nan) (-0.67500,nan) (-0.67000,nan) (-0.66500,nan) (-0.66000,nan) (-0.65500,nan) (-0.65000,nan) (-0.64500,nan) (-0.64000,nan) (-0.63500,nan) (-0.63000,nan) (-0.62500,nan) (-0.62000,nan) (-0.61500,nan) (-0.61000,nan) (-0.60500,nan) (-0.60000,nan) (-0.59500,nan) (-0.59000,nan) (-0.58500,nan) (-0.58000,nan) (-0.57500,nan) (-0.57000,nan) (-0.56500,nan) (-0.56000,nan) (-0.55500,nan) (-0.55000,nan) (-0.54500,nan) (-0.54000,nan) (-0.53500,nan) (-0.53000,nan) (-0.52500,nan) (-0.52000,nan) (-0.51500,nan) (-0.51000,nan) (-0.50500,nan) (-0.50000,nan) (-0.49500,nan) (-0.49000,nan) (-0.48500,nan) (-0.48000,nan) (-0.47500,nan) (-0.47000,nan) (-0.46500,nan) (-0.46000,nan) (-0.45500,nan) (-0.45000,nan) (-0.44500,nan) (-0.44000,nan) (-0.43500,nan) (-0.43000,nan) (-0.42500,nan) (-0.42000,nan) (-0.41500,nan) (-0.41000,nan) (-0.40500,nan) (-0.40000,nan) (-0.39500,nan) (-0.39000,nan) (-0.38500,nan) (-0.38000,nan) (-0.37500,nan) (-0.37000,nan) (-0.36500,nan) (-0.36000,nan) (-0.35500,nan) (-0.35000,nan) (-0.34500,nan) (-0.34000,nan) (-0.33500,nan) (-0.33000,nan) (-0.32500,nan) (-0.32000,nan) (-0.31500,nan) (-0.31000,nan) (-0.30500,nan) (-0.30000,nan) (-0.29500,nan) (-0.29000,nan) (-0.28500,nan) (-0.28000,nan) (-0.27500,nan) (-0.27000,nan) (-0.26500,nan) (-0.26000,nan) (-0.25500,nan) (-0.25000,nan) (-0.24500,nan) (-0.24000,nan) (-0.23500,nan) (-0.23000,nan) (-0.22500,nan) (-0.22000,nan) (-0.21500,nan) (-0.21000,nan) (-0.20500,nan) (-0.20000,nan) (-0.19500,nan) (-0.19000,nan) (-0.18500,nan) (-0.18000,nan) (-0.17500,nan) (-0.17000,nan) (-0.16500,nan) (-0.16000,nan) (-0.15500,nan) (-0.15000,nan) (-0.14500,nan) (-0.14000,nan) (-0.13500,nan) (-0.13000,nan) (-0.12500,nan) (-0.12000,nan) (-0.11500,nan) (-0.11000,nan) (-0.10500,nan) (-0.10000,nan) (-0.09500,nan) (-0.09000,nan) (-0.08500,nan) (-0.08000,nan) (-0.07500,nan) (-0.07000,nan) (-0.06500,nan) (-0.06000,nan) (-0.05500,nan) (-0.05000,nan) (-0.04500,nan) (-0.04000,nan) (-0.03500,nan) (-0.03000,nan) (-0.02500,nan) (-0.02000,nan) (-0.01500,nan) (-0.01000,nan) (-0.00500,nan) (0.00000,nan) (0.00500,nan) (0.01000,nan) (0.01500,nan) (0.02000,nan) (0.02500,nan) (0.03000,nan) (0.03500,nan) (0.04000,nan) (0.04500,nan) (0.05000,nan) (0.05500,nan) (0.06000,nan) (0.06500,nan) (0.07000,nan) (0.07500,nan) (0.08000,nan) (0.08500,nan) (0.09000,nan) (0.09500,nan) (0.10000,nan) (0.10500,nan) (0.11000,nan) (0.11500,nan) (0.12000,nan) (0.12500,nan) (0.13000,nan) (0.13500,nan) (0.14000,nan) (0.14500,nan) (0.15000,nan) (0.15500,nan) (0.16000,nan) (0.16500,nan) (0.17000,nan) (0.17500,nan) (0.18000,nan) (0.18500,nan) (0.19000,nan) (0.19500,nan) (0.20000,nan) (0.20500,nan) (0.21000,nan) (0.21500,nan) (0.22000,nan) (0.22500,nan) (0.23000,nan) (0.23500,nan) (0.24000,nan) (0.24500,nan) (0.25000,nan) (0.25500,nan) (0.26000,nan) (0.26500,nan) (0.27000,nan) (0.27500,nan) (0.28000,nan) (0.28500,nan) (0.29000,nan) (0.29500,nan) (0.30000,nan) (0.30500,nan) (0.31000,nan) (0.31500,nan) (0.32000,nan) (0.32500,nan) (0.33000,nan) (0.33500,nan) (0.34000,nan) (0.34500,nan) (0.35000,nan) (0.35500,nan) (0.36000,nan) (0.36500,nan) (0.37000,nan) (0.37500,nan) (0.38000,nan) (0.38500,nan) (0.39000,nan) (0.39500,nan) (0.40000,nan) (0.40500,nan) (0.41000,nan) (0.41500,nan) (0.42000,nan) (0.42500,nan) (0.43000,nan) (0.43500,nan) (0.44000,nan) (0.44500,nan) (0.45000,nan) (0.45500,nan) (0.46000,nan) (0.46500,nan) (0.47000,nan) (0.47500,nan) (0.48000,nan) (0.48500,nan) (0.49000,nan) (0.49500,nan) (0.50000,nan) (0.50500,nan) (0.51000,nan) (0.51500,nan) (0.52000,nan) (0.52500,nan) (0.53000,nan) (0.53500,nan) (0.54000,nan) (0.54500,nan) (0.55000,nan) (0.55500,nan) (0.56000,nan) (0.56500,nan) (0.57000,nan) (0.57500,nan) (0.58000,nan) (0.58500,nan) (0.59000,nan) (0.59500,nan) (0.60000,nan) (0.60500,nan) (0.61000,nan) (0.61500,nan) (0.62000,nan) (0.62500,nan) (0.63000,nan) (0.63500,nan) (0.64000,nan) (0.64500,nan) (0.65000,nan) (0.65500,nan) (0.66000,nan) (0.66500,nan) (0.67000,nan) (0.67500,nan) (0.68000,nan) (0.68500,nan) (0.69000,nan) (0.69500,nan) (0.70000,nan) (0.70500,nan) (0.71000,nan) (0.71500,nan) (0.72000,nan) (0.72500,nan) (0.73000,nan) (0.73500,nan) (0.74000,nan) (0.74500,nan) (0.75000,nan) (0.75500,nan) (0.76000,nan) (0.76500,nan) (0.77000,nan) (0.77500,nan) (0.78000,nan) (0.78500,nan) (0.79000,nan) (0.79500,nan) (0.80000,nan) (0.80500,nan) (0.81000,nan) (0.81500,nan) (0.82000,nan) (0.82500,nan) (0.83000,nan) (0.83500,nan) (0.84000,nan) (0.84500,nan) (0.85000,nan) (0.85500,nan) (0.86000,nan) (0.86500,nan) (0.87000,nan) (0.87500,nan) (0.88000,nan) (0.88500,nan) (0.89000,nan) (0.89500,nan) (0.90000,nan) (0.90500,nan) (0.91000,nan) (0.91500,nan) (0.92000,nan) (0.92500,nan) (0.93000,nan) (0.93500,nan) (0.94000,nan) (0.94500,nan) (0.95000,nan) (0.95500,nan) (0.96000,nan) (0.96500,nan) (0.97000,nan) (0.97500,nan) (0.98000,nan) (0.98500,nan) (0.99000,nan) (0.99500,nan) (1.00000,nan) (1.00500,nan) (1.01000,nan) (1.01500,nan) (1.02000,nan) (1.02500,nan) (1.03000,nan) (1.03500,nan) (1.04000,nan) (1.04500,nan) (1.05000,nan) (1.05500,nan) (1.06000,nan) (1.06500,nan) (1.07000,nan) (1.07500,nan) (1.08000,nan) (1.08500,nan) (1.09000,nan) (1.09500,nan) (1.10000,nan) (1.10500,nan) (1.11000,nan) (1.11500,nan) (1.12000,nan) (1.12500,nan) (1.13000,nan) (1.13500,nan) (1.14000,nan) (1.14500,nan) (1.15000,nan) (1.15500,nan) (1.16000,nan) (1.16500,nan) (1.17000,nan) (1.17500,nan) (1.18000,nan) (1.18500,nan) (1.19000,nan) (1.19500,nan) (1.20000,nan) (1.20500,nan) (1.21000,nan) (1.21500,nan) (1.22000,nan) (1.22500,nan) (1.23000,nan) (1.23500,nan) (1.24000,nan) (1.24500,nan) (1.25000,nan) (1.25500,nan) (1.26000,nan) (1.26500,nan) (1.27000,nan) (1.27500,nan) (1.28000,nan) (1.28500,nan) (1.29000,nan) (1.29500,nan) (1.30000,nan) (1.30500,nan) (1.31000,nan) (1.31500,nan) (1.32000,nan) (1.32500,nan) (1.33000,nan) (1.33500,nan) (1.34000,nan) (1.34500,nan) (1.35000,nan) (1.35500,nan) (1.36000,nan) (1.36500,nan) (1.37000,nan) (1.37500,nan) (1.38000,nan) (1.38500,nan) (1.39000,nan) (1.39500,nan) (1.40000,nan) (1.40500,nan) (1.41000,nan) (1.41500,nan) (1.42000,nan) (1.42500,nan) (1.43000,nan) (1.43500,nan) (1.44000,nan) (1.44500,nan) (1.45000,nan) (1.45500,nan) (1.46000,nan) (1.46500,nan) (1.47000,nan) (1.47500,nan) (1.48000,nan) (1.48500,nan) (1.49000,nan) (1.49500,nan) (1.50000,nan) (1.50500,nan) (1.51000,nan) (1.51500,nan) (1.52000,nan) (1.52500,nan) (1.53000,nan) (1.53500,nan) (1.54000,nan) (1.54500,nan) (1.55000,nan) (1.55500,nan) (1.56000,nan) (1.56500,nan) (1.57000,nan) (1.57500,nan) (1.58000,nan) (1.58500,nan) (1.59000,nan) (1.59500,nan) (1.60000,nan) (1.60500,nan) (1.61000,nan) (1.61500,nan) (1.62000,nan) (1.62500,nan) (1.63000,nan) (1.63500,nan) (1.64000,nan) (1.64500,nan) (1.65000,nan) (1.65500,nan) (1.66000,nan) (1.66500,nan) (1.67000,nan) (1.67500,nan) (1.68000,nan) (1.68500,nan) (1.69000,nan) (1.69500,nan) (1.70000,nan) (1.70500,nan) (1.71000,nan) (1.71500,nan) (1.72000,nan) (1.72500,nan) (1.73000,nan) (1.73500,nan) (1.74000,nan) (1.74500,nan) (1.75000,nan) (1.75500,nan) (1.76000,nan) (1.76500,nan) (1.77000,nan) (1.77500,nan) (1.78000,nan) (1.78500,nan) (1.79000,nan) (1.79500,nan) (1.80000,nan) (1.80500,nan) (1.81000,nan) (1.81500,nan) (1.82000,nan) (1.82500,nan) (1.83000,nan) (1.83500,nan) (1.84000,nan) (1.84500,nan) (1.85000,nan) (1.85500,nan) (1.86000,nan) (1.86500,nan) (1.87000,nan) (1.87500,nan) (1.88000,nan) (1.88500,nan) (1.89000,nan) (1.89500,nan) (1.90000,nan) (1.90500,nan) (1.91000,nan) (1.91500,nan) (1.92000,nan) (1.92500,nan) (1.93000,nan) (1.93500,nan) (1.94000,nan) (1.94500,nan) (1.95000,nan) (1.95500,nan) (1.96000,nan) (1.96500,nan) (1.97000,nan) (1.97500,nan) (1.98000,nan) (1.98500,nan) (1.99000,nan) (1.99500,nan) (2.00000,nan) (2.00500,nan) (2.01000,nan) (2.01500,nan) (2.02000,nan) (2.02500,nan) (2.03000,nan) (2.03500,nan) (2.04000,nan) (2.04500,nan) (2.05000,nan) (2.05500,nan) (2.06000,nan) (2.06500,nan) (2.07000,nan) (2.07500,nan) (2.08000,nan) (2.08500,nan) (2.09000,nan) (2.09500,nan) (2.10000,nan) (2.10500,nan) (2.11000,nan) (2.11500,nan) (2.12000,nan) (2.12500,nan) (2.13000,nan) (2.13500,nan) (2.14000,nan) (2.14500,nan) (2.15000,nan) (2.15500,nan) (2.16000,nan) (2.16500,nan) (2.17000,nan) (2.17500,nan) (2.18000,nan) (2.18500,nan) (2.19000,nan) (2.19500,nan) (2.20000,nan) (2.20500,nan) (2.21000,nan) (2.21500,nan) (2.22000,nan) (2.22500,nan) (2.23000,nan) (2.23500,nan) (2.24000,nan) (2.24500,nan) (2.25000,nan) (2.25500,nan) (2.26000,nan) (2.26500,nan) (2.27000,nan) (2.27500,nan) (2.28000,nan) (2.28500,nan) (2.29000,nan) (2.29500,nan) (2.30000,nan) (2.30500,nan) (2.31000,nan) (2.31500,nan) (2.32000,nan) (2.32500,nan) (2.33000,nan) (2.33500,nan) (2.34000,nan) (2.34500,nan) (2.35000,nan) (2.35500,nan) (2.36000,nan) (2.36500,nan) (2.37000,nan) (2.37500,nan) (2.38000,nan) (2.38500,nan) (2.39000,nan) (2.39500,nan) (2.40000,nan) (2.40500,nan) (2.41000,nan) (2.41500,nan) (2.42000,nan) (2.42500,nan) (2.43000,nan) (2.43500,nan) (2.44000,nan) (2.44500,nan) (2.45000,nan) (2.45500,nan) (2.46000,nan) (2.46500,nan) (2.47000,nan) (2.47500,nan) (2.48000,nan) (2.48500,nan) (2.49000,nan) (2.49500,nan) (2.50000,nan) (2.50500,nan) (2.51000,nan) (2.51500,nan) (2.52000,nan) (2.52500,nan) (2.53000,nan) (2.53500,nan) (2.54000,nan) (2.54500,nan) (2.55000,nan) (2.55500,nan) (2.56000,nan) (2.56500,nan) (2.57000,nan) (2.57500,nan) (2.58000,nan) (2.58500,nan) (2.59000,nan) (2.59500,nan) (2.60000,nan) (2.60500,nan) (2.61000,nan) (2.61500,nan) (2.62000,nan) (2.62500,nan) (2.63000,nan) (2.63500,nan) (2.64000,nan) (2.64500,nan) (2.65000,nan) (2.65500,nan) (2.66000,nan) (2.66500,nan) (2.67000,nan) (2.67500,nan) (2.68000,nan) (2.68500,nan) (2.69000,nan) (2.69500,nan) (2.70000,nan) (2.70500,nan) (2.71000,nan) (2.71500,nan) (2.72000,nan) (2.72500,nan) (2.73000,nan) (2.73500,nan) (2.74000,nan) (2.74500,nan) (2.75000,nan) (2.75500,nan) (2.76000,nan) (2.76500,nan) (2.77000,nan) (2.77500,nan) (2.78000,nan) (2.78500,nan) (2.79000,nan) (2.79500,nan) (2.80000,nan) (2.80500,nan) (2.81000,nan) (2.81500,nan) (2.82000,nan) (2.82500,nan) (2.83000,nan) (2.83500,nan) (2.84000,nan) (2.84500,nan) (2.85000,nan) (2.85500,nan) (2.86000,nan) (2.86500,nan) (2.87000,nan) (2.87500,nan) (2.88000,nan) (2.88500,nan) (2.89000,nan) (2.89500,nan) (2.90000,nan) (2.90500,nan) (2.91000,nan) (2.91500,nan) (2.92000,nan) (2.92500,nan) (2.93000,nan) (2.93500,nan) (2.94000,nan) (2.94500,nan) (2.95000,nan) (2.95500,nan) (2.96000,nan) (2.96500,nan) (2.97000,nan) (2.97500,nan) (2.98000,nan) (2.98500,nan) (2.99000,nan) (2.99500,nan) (3.00000,nan) (3.00500,nan) (3.01000,nan) (3.01500,nan) (3.02000,nan) (3.02500,nan) (3.03000,nan) (3.03500,nan) (3.04000,nan) (3.04500,nan) (3.05000,nan) (3.05500,nan) (3.06000,nan) (3.06500,-1.6859561) (3.07000,-1.3153455) (3.07500,-1.1218478) (3.08000,-0.9910006) (3.08500,-0.8925616) (3.09000,-0.8139757) (3.09500,-0.7488091) (3.10000,-0.6933207) (3.10500,-0.6451419) (3.11000,-0.6026778) (3.11500,-0.5648032) (3.12000,-0.5306941) (3.12500,-0.4997294) (3.13000,-0.4714291) (3.13500,-0.4454145) (3.14000,-0.4213816) (3.14500,-0.3990827) (3.15000,-0.3783131) (3.15500,-0.3589020) (3.16000,-0.3407050) (3.16500,-0.3235993) (3.17000,-0.3074794) (3.17500,-0.2922539) (3.18000,-0.2778433) (3.18500,-0.2641778) (3.19000,-0.2511960) (3.19500,-0.2388433) (3.20000,-0.2270714) (3.20500,-0.2158366) (3.21000,-0.2051002) (3.21500,-0.1948269) (3.22000,-0.1849849) (3.22500,-0.1755452) (3.23000,-0.1664817) (3.23500,-0.1577701) (3.24000,-0.1493884) (3.24500,-0.1413163) (3.25000,-0.1335351) (3.25500,-0.1260274) (3.26000,-0.1187773) (3.26500,-0.1117698) (3.27000,-0.1049911) (3.27500,-0.0984281) (3.28000,-0.0920688) (3.28500,-0.0859017) (3.29000,-0.0799162) (3.29500,-0.0741022) (3.30000,-0.0684501) (3.30500,-0.0629510) (3.31000,-0.0575961) (3.31500,-0.0523775) (3.32000,-0.0472873) (3.32500,-0.0423180) (3.33000,-0.0374626) (3.33500,-0.0327140) (3.34000,-0.0280658) (3.34500,-0.0235115) (3.35000,-0.0190448) (3.35500,-0.0146597) (3.36000,-0.0103503) (3.36500,-0.0061106) (3.37000,-0.0019351) (3.37500,0.0021820) (3.38000,0.0062464) (3.38500,0.0102635) (3.39000,0.0142391) (3.39500,0.0181788) (3.40000,0.0220883) (3.40500,0.0259733) (3.41000,0.0298399) (3.41500,0.0336940) (3.42000,0.0375419) (3.42500,0.0413901) (3.43000,0.0452453) (3.43500,0.0491145) (3.44000,0.0530050) (3.44500,0.0569248) (3.45000,0.0608820) (3.45500,0.0648856) (3.46000,0.0689449) (3.46500,0.0730702) (3.47000,0.0772725) (3.47500,0.0815638) (3.48000,0.0859574) (3.48500,0.0904675) (3.49000,0.0951101) (3.49500,0.0999029) (3.50000,0.1048655) (3.50500,0.1100201) (3.51000,0.1153915) (3.51500,0.1210078) (3.52000,0.1269012) (3.52500,0.1331084) (3.53000,0.1396718) (3.53500,0.1466406) (3.54000,0.1540722) (3.54500,0.1620344) (3.55000,0.1706074) (3.55500,0.1798875) (3.56000,0.1899910) (3.56500,0.2010601) (3.57000,0.2132704) (3.57500,0.2268421) (3.58000,0.2420554) (3.58500,0.2592727) (3.59000,0.2789726) (3.59500,0.3018017) (3.60000,0.3286595) (3.60500,0.3608403) (3.61000,0.4002904) (3.61500,0.4501056) (3.62000,0.5156061) (3.62500,0.6070205) (3.63000,0.7479290) (3.63500,1.0162853) (3.64000,nan) (3.64500,nan) (3.65000,nan) (3.65500,nan) (3.66000,nan) (3.66500,nan) (3.67000,nan) (3.67500,nan) (3.68000,nan) (3.68500,-1.4131863) (3.69000,-0.8583087) (3.69500,-0.6563802) (3.70000,-0.5388488) (3.70500,-0.4592300) (3.71000,-0.4008066) (3.71500,-0.3557185) (3.72000,-0.3196762) (3.72500,-0.2901024) (3.73000,-0.2653385) (3.73500,-0.2442620) (3.74000,-0.2260822) (3.74500,-0.2102239) (3.75000,-0.1962579) (3.75500,-0.1838564) (3.76000,-0.1727648) (3.76500,-0.1627817) (3.77000,-0.1537456) (3.77500,-0.1455255) (3.78000,-0.1380140) (3.78500,-0.1311220) (3.79000,-0.1247748) (3.79500,-0.1189097) (3.80000,-0.1134733) (3.80500,-0.1084199) (3.81000,-0.1037103) (3.81500,-0.0993104) (3.82000,-0.0951907) (3.82500,-0.0913254) (3.83000,-0.0876916) (3.83500,-0.0842694) (3.84000,-0.0810410) (3.84500,-0.0779906) (3.85000,-0.0751042) (3.85500,-0.0723691) (3.86000,-0.0697740) (3.86500,-0.0673087) (3.87000,-0.0649640) (3.87500,-0.0627315) (3.88000,-0.0606037) (3.88500,-0.0585737) (3.89000,-0.0566352) (3.89500,-0.0547823) (3.90000,-0.0530100) (3.90500,-0.0513132) (3.91000,-0.0496876) (3.91500,-0.0481290) (3.92000,-0.0466336) (3.92500,-0.0451980) (3.93000,-0.0438189) (3.93500,-0.0424933) (3.94000,-0.0412184) (3.94500,-0.0399915) (3.95000,-0.0388103) (3.95500,-0.0376725) (3.96000,-0.0365760) (3.96500,-0.0355187) (3.97000,-0.0344989) (3.97500,-0.0335148) (3.98000,-0.0325647) (3.98500,-0.0316471) (3.99000,-0.0307607) (3.99500,-0.0299039) (4.00000,-0.0290756)};\addlegendentry{X}
\end{axis}\end{tikzpicture}\end{center}
\begin{center}\begin{tikzpicture}\begin{axis}[width=.77\textwidth,height=4.8cm,xlabel={$\log_{10}(\omega/\Omega_{ci})$},ylabel={$\log_{10}n_{\rm ref}^2$},xmin=-2,xmax=4,ymin=-2,ymax=3,unbounded coords=jump,title={$n_e=4\times10^{20}\,\mathrm m^{-3},\ B=5\,\mathrm T$},legend style={font=\small,at={(1.02,1)},anchor=north west}]
\addplot[thick,structurecolor] coordinates {(-2.00000,nan) (-1.99500,nan) (-1.99000,nan) (-1.98500,nan) (-1.98000,nan) (-1.97500,nan) (-1.97000,nan) (-1.96500,nan) (-1.96000,nan) (-1.95500,nan) (-1.95000,nan) (-1.94500,nan) (-1.94000,nan) (-1.93500,nan) (-1.93000,nan) (-1.92500,nan) (-1.92000,nan) (-1.91500,nan) (-1.91000,nan) (-1.90500,nan) (-1.90000,nan) (-1.89500,nan) (-1.89000,nan) (-1.88500,nan) (-1.88000,nan) (-1.87500,nan) (-1.87000,nan) (-1.86500,nan) (-1.86000,nan) (-1.85500,nan) (-1.85000,nan) (-1.84500,nan) (-1.84000,nan) (-1.83500,nan) (-1.83000,nan) (-1.82500,nan) (-1.82000,nan) (-1.81500,nan) (-1.81000,nan) (-1.80500,nan) (-1.80000,nan) (-1.79500,nan) (-1.79000,nan) (-1.78500,nan) (-1.78000,nan) (-1.77500,nan) (-1.77000,nan) (-1.76500,nan) (-1.76000,nan) (-1.75500,nan) (-1.75000,nan) (-1.74500,nan) (-1.74000,nan) (-1.73500,nan) (-1.73000,nan) (-1.72500,nan) (-1.72000,nan) (-1.71500,nan) (-1.71000,nan) (-1.70500,nan) (-1.70000,nan) (-1.69500,nan) (-1.69000,nan) (-1.68500,nan) (-1.68000,nan) (-1.67500,nan) (-1.67000,nan) (-1.66500,nan) (-1.66000,nan) (-1.65500,nan) (-1.65000,nan) (-1.64500,nan) (-1.64000,nan) (-1.63500,nan) (-1.63000,nan) (-1.62500,nan) (-1.62000,nan) (-1.61500,nan) (-1.61000,nan) (-1.60500,nan) (-1.60000,nan) (-1.59500,nan) (-1.59000,nan) (-1.58500,nan) (-1.58000,nan) (-1.57500,nan) (-1.57000,nan) (-1.56500,nan) (-1.56000,nan) (-1.55500,nan) (-1.55000,nan) (-1.54500,nan) (-1.54000,nan) (-1.53500,nan) (-1.53000,nan) (-1.52500,nan) (-1.52000,nan) (-1.51500,nan) (-1.51000,nan) (-1.50500,nan) (-1.50000,nan) (-1.49500,nan) (-1.49000,nan) (-1.48500,nan) (-1.48000,nan) (-1.47500,nan) (-1.47000,nan) (-1.46500,nan) (-1.46000,nan) (-1.45500,nan) (-1.45000,nan) (-1.44500,nan) (-1.44000,nan) (-1.43500,nan) (-1.43000,nan) (-1.42500,nan) (-1.42000,nan) (-1.41500,nan) (-1.41000,nan) (-1.40500,nan) (-1.40000,nan) (-1.39500,nan) (-1.39000,nan) (-1.38500,nan) (-1.38000,nan) (-1.37500,nan) (-1.37000,nan) (-1.36500,nan) (-1.36000,nan) (-1.35500,nan) (-1.35000,nan) (-1.34500,nan) (-1.34000,nan) (-1.33500,nan) (-1.33000,nan) (-1.32500,nan) (-1.32000,nan) (-1.31500,nan) (-1.31000,nan) (-1.30500,nan) (-1.30000,nan) (-1.29500,nan) (-1.29000,nan) (-1.28500,nan) (-1.28000,nan) (-1.27500,nan) (-1.27000,nan) (-1.26500,nan) (-1.26000,nan) (-1.25500,nan) (-1.25000,nan) (-1.24500,nan) (-1.24000,nan) (-1.23500,nan) (-1.23000,nan) (-1.22500,nan) (-1.22000,nan) (-1.21500,nan) (-1.21000,nan) (-1.20500,nan) (-1.20000,nan) (-1.19500,nan) (-1.19000,nan) (-1.18500,nan) (-1.18000,nan) (-1.17500,nan) (-1.17000,nan) (-1.16500,nan) (-1.16000,nan) (-1.15500,nan) (-1.15000,nan) (-1.14500,nan) (-1.14000,nan) (-1.13500,nan) (-1.13000,nan) (-1.12500,nan) (-1.12000,nan) (-1.11500,nan) (-1.11000,nan) (-1.10500,nan) (-1.10000,nan) (-1.09500,nan) (-1.09000,nan) (-1.08500,nan) (-1.08000,nan) (-1.07500,nan) (-1.07000,nan) (-1.06500,nan) (-1.06000,nan) (-1.05500,nan) (-1.05000,nan) (-1.04500,nan) (-1.04000,nan) (-1.03500,nan) (-1.03000,nan) (-1.02500,nan) (-1.02000,nan) (-1.01500,nan) (-1.01000,nan) (-1.00500,nan) (-1.00000,nan) (-0.99500,nan) (-0.99000,nan) (-0.98500,nan) (-0.98000,nan) (-0.97500,nan) (-0.97000,nan) (-0.96500,nan) (-0.96000,nan) (-0.95500,nan) (-0.95000,nan) (-0.94500,nan) (-0.94000,nan) (-0.93500,nan) (-0.93000,nan) (-0.92500,nan) (-0.92000,nan) (-0.91500,nan) (-0.91000,nan) (-0.90500,nan) (-0.90000,nan) (-0.89500,nan) (-0.89000,nan) (-0.88500,nan) (-0.88000,nan) (-0.87500,nan) (-0.87000,nan) (-0.86500,nan) (-0.86000,nan) (-0.85500,nan) (-0.85000,nan) (-0.84500,nan) (-0.84000,nan) (-0.83500,nan) (-0.83000,nan) (-0.82500,nan) (-0.82000,nan) (-0.81500,nan) (-0.81000,nan) (-0.80500,nan) (-0.80000,nan) (-0.79500,nan) (-0.79000,nan) (-0.78500,nan) (-0.78000,nan) (-0.77500,nan) (-0.77000,nan) (-0.76500,nan) (-0.76000,nan) (-0.75500,nan) (-0.75000,nan) (-0.74500,nan) (-0.74000,nan) (-0.73500,nan) (-0.73000,nan) (-0.72500,nan) (-0.72000,nan) (-0.71500,nan) (-0.71000,nan) (-0.70500,nan) (-0.70000,nan) (-0.69500,nan) (-0.69000,nan) (-0.68500,nan) (-0.68000,nan) (-0.67500,nan) (-0.67000,nan) (-0.66500,nan) (-0.66000,nan) (-0.65500,nan) (-0.65000,nan) (-0.64500,nan) (-0.64000,nan) (-0.63500,nan) (-0.63000,nan) (-0.62500,nan) (-0.62000,nan) (-0.61500,nan) (-0.61000,nan) (-0.60500,nan) (-0.60000,nan) (-0.59500,nan) (-0.59000,nan) (-0.58500,nan) (-0.58000,nan) (-0.57500,nan) (-0.57000,nan) (-0.56500,nan) (-0.56000,nan) (-0.55500,nan) (-0.55000,nan) (-0.54500,nan) (-0.54000,nan) (-0.53500,nan) (-0.53000,nan) (-0.52500,nan) (-0.52000,nan) (-0.51500,nan) (-0.51000,nan) (-0.50500,nan) (-0.50000,nan) (-0.49500,nan) (-0.49000,nan) (-0.48500,nan) (-0.48000,nan) (-0.47500,nan) (-0.47000,nan) (-0.46500,nan) (-0.46000,nan) (-0.45500,nan) (-0.45000,nan) (-0.44500,nan) (-0.44000,nan) (-0.43500,nan) (-0.43000,nan) (-0.42500,nan) (-0.42000,nan) (-0.41500,nan) (-0.41000,nan) (-0.40500,nan) (-0.40000,nan) (-0.39500,nan) (-0.39000,nan) (-0.38500,nan) (-0.38000,nan) (-0.37500,nan) (-0.37000,nan) (-0.36500,nan) (-0.36000,nan) (-0.35500,nan) (-0.35000,nan) (-0.34500,nan) (-0.34000,nan) (-0.33500,nan) (-0.33000,nan) (-0.32500,nan) (-0.32000,nan) (-0.31500,nan) (-0.31000,nan) (-0.30500,nan) (-0.30000,nan) (-0.29500,nan) (-0.29000,nan) (-0.28500,nan) (-0.28000,nan) (-0.27500,nan) (-0.27000,nan) (-0.26500,nan) (-0.26000,nan) (-0.25500,nan) (-0.25000,nan) (-0.24500,nan) (-0.24000,nan) (-0.23500,nan) (-0.23000,nan) (-0.22500,nan) (-0.22000,nan) (-0.21500,nan) (-0.21000,nan) (-0.20500,nan) (-0.20000,nan) (-0.19500,nan) (-0.19000,nan) (-0.18500,nan) (-0.18000,nan) (-0.17500,nan) (-0.17000,nan) (-0.16500,nan) (-0.16000,nan) (-0.15500,nan) (-0.15000,nan) (-0.14500,nan) (-0.14000,nan) (-0.13500,nan) (-0.13000,nan) (-0.12500,nan) (-0.12000,nan) (-0.11500,nan) (-0.11000,nan) (-0.10500,nan) (-0.10000,nan) (-0.09500,nan) (-0.09000,nan) (-0.08500,nan) (-0.08000,nan) (-0.07500,nan) (-0.07000,nan) (-0.06500,nan) (-0.06000,nan) (-0.05500,nan) (-0.05000,nan) (-0.04500,nan) (-0.04000,nan) (-0.03500,nan) (-0.03000,nan) (-0.02500,nan) (-0.02000,nan) (-0.01500,nan) (-0.01000,nan) (-0.00500,nan) (0.00000,nan) (0.00500,nan) (0.01000,nan) (0.01500,nan) (0.02000,nan) (0.02500,nan) (0.03000,nan) (0.03500,nan) (0.04000,nan) (0.04500,nan) (0.05000,nan) (0.05500,nan) (0.06000,nan) (0.06500,nan) (0.07000,nan) (0.07500,nan) (0.08000,nan) (0.08500,nan) (0.09000,nan) (0.09500,nan) (0.10000,nan) (0.10500,nan) (0.11000,nan) (0.11500,nan) (0.12000,nan) (0.12500,nan) (0.13000,nan) (0.13500,nan) (0.14000,nan) (0.14500,nan) (0.15000,nan) (0.15500,nan) (0.16000,nan) (0.16500,nan) (0.17000,nan) (0.17500,nan) (0.18000,nan) (0.18500,nan) (0.19000,nan) (0.19500,nan) (0.20000,nan) (0.20500,nan) (0.21000,nan) (0.21500,nan) (0.22000,nan) (0.22500,nan) (0.23000,nan) (0.23500,nan) (0.24000,nan) (0.24500,nan) (0.25000,nan) (0.25500,nan) (0.26000,nan) (0.26500,nan) (0.27000,nan) (0.27500,nan) (0.28000,nan) (0.28500,nan) (0.29000,nan) (0.29500,nan) (0.30000,nan) (0.30500,nan) (0.31000,nan) (0.31500,nan) (0.32000,nan) (0.32500,nan) (0.33000,nan) (0.33500,nan) (0.34000,nan) (0.34500,nan) (0.35000,nan) (0.35500,nan) (0.36000,nan) (0.36500,nan) (0.37000,nan) (0.37500,nan) (0.38000,nan) (0.38500,nan) (0.39000,nan) (0.39500,nan) (0.40000,nan) (0.40500,nan) (0.41000,nan) (0.41500,nan) (0.42000,nan) (0.42500,nan) (0.43000,nan) (0.43500,nan) (0.44000,nan) (0.44500,nan) (0.45000,nan) (0.45500,nan) (0.46000,nan) (0.46500,nan) (0.47000,nan) (0.47500,nan) (0.48000,nan) (0.48500,nan) (0.49000,nan) (0.49500,nan) (0.50000,nan) (0.50500,nan) (0.51000,nan) (0.51500,nan) (0.52000,nan) (0.52500,nan) (0.53000,nan) (0.53500,nan) (0.54000,nan) (0.54500,nan) (0.55000,nan) (0.55500,nan) (0.56000,nan) (0.56500,nan) (0.57000,nan) (0.57500,nan) (0.58000,nan) (0.58500,nan) (0.59000,nan) (0.59500,nan) (0.60000,nan) (0.60500,nan) (0.61000,nan) (0.61500,nan) (0.62000,nan) (0.62500,nan) (0.63000,nan) (0.63500,nan) (0.64000,nan) (0.64500,nan) (0.65000,nan) (0.65500,nan) (0.66000,nan) (0.66500,nan) (0.67000,nan) (0.67500,nan) (0.68000,nan) (0.68500,nan) (0.69000,nan) (0.69500,nan) (0.70000,nan) (0.70500,2.9993704) (0.71000,2.9952013) (0.71500,2.9910245) (0.72000,2.9868401) (0.72500,2.9826480) (0.73000,2.9784484) (0.73500,2.9742413) (0.74000,2.9700268) (0.74500,2.9658049) (0.75000,2.9615757) (0.75500,2.9573393) (0.76000,2.9530957) (0.76500,2.9488450) (0.77000,2.9445872) (0.77500,2.9403225) (0.78000,2.9360508) (0.78500,2.9317722) (0.79000,2.9274868) (0.79500,2.9231947) (0.80000,2.9188959) (0.80500,2.9145904) (0.81000,2.9102783) (0.81500,2.9059597) (0.82000,2.9016347) (0.82500,2.8973032) (0.83000,2.8929654) (0.83500,2.8886213) (0.84000,2.8842710) (0.84500,2.8799145) (0.85000,2.8755519) (0.85500,2.8711832) (0.86000,2.8668085) (0.86500,2.8624279) (0.87000,2.8580413) (0.87500,2.8536489) (0.88000,2.8492507) (0.88500,2.8448468) (0.89000,2.8404372) (0.89500,2.8360220) (0.90000,2.8316012) (0.90500,2.8271749) (0.91000,2.8227431) (0.91500,2.8183060) (0.92000,2.8138634) (0.92500,2.8094156) (0.93000,2.8049625) (0.93500,2.8005042) (0.94000,2.7960407) (0.94500,2.7915722) (0.95000,2.7870986) (0.95500,2.7826200) (0.96000,2.7781365) (0.96500,2.7736480) (0.97000,2.7691548) (0.97500,2.7646567) (0.98000,2.7601538) (0.98500,2.7556463) (0.99000,2.7511341) (0.99500,2.7466174) (1.00000,2.7420960) (1.00500,2.7375702) (1.01000,2.7330399) (1.01500,2.7285052) (1.02000,2.7239661) (1.02500,2.7194227) (1.03000,2.7148750) (1.03500,2.7103232) (1.04000,2.7057671) (1.04500,2.7012069) (1.05000,2.6966426) (1.05500,2.6920743) (1.06000,2.6875019) (1.06500,2.6829257) (1.07000,2.6783455) (1.07500,2.6737614) (1.08000,2.6691735) (1.08500,2.6645819) (1.09000,2.6599865) (1.09500,2.6553874) (1.10000,2.6507846) (1.10500,2.6461783) (1.11000,2.6415684) (1.11500,2.6369549) (1.12000,2.6323380) (1.12500,2.6277177) (1.13000,2.6230939) (1.13500,2.6184668) (1.14000,2.6138363) (1.14500,2.6092026) (1.15000,2.6045657) (1.15500,2.5999255) (1.16000,2.5952822) (1.16500,2.5906358) (1.17000,2.5859862) (1.17500,2.5813337) (1.18000,2.5766781) (1.18500,2.5720196) (1.19000,2.5673581) (1.19500,2.5626938) (1.20000,2.5580265) (1.20500,2.5533565) (1.21000,2.5486837) (1.21500,2.5440082) (1.22000,2.5393299) (1.22500,2.5346490) (1.23000,2.5299655) (1.23500,2.5252793) (1.24000,2.5205906) (1.24500,2.5158994) (1.25000,2.5112057) (1.25500,2.5065095) (1.26000,2.5018109) (1.26500,2.4971099) (1.27000,2.4924066) (1.27500,2.4877009) (1.28000,2.4829930) (1.28500,2.4782828) (1.29000,2.4735704) (1.29500,2.4688558) (1.30000,2.4641391) (1.30500,2.4594202) (1.31000,2.4546993) (1.31500,2.4499763) (1.32000,2.4452513) (1.32500,2.4405244) (1.33000,2.4357954) (1.33500,2.4310646) (1.34000,2.4263318) (1.34500,2.4215972) (1.35000,2.4168608) (1.35500,2.4121226) (1.36000,2.4073826) (1.36500,2.4026409) (1.37000,2.3978975) (1.37500,2.3931524) (1.38000,2.3884057) (1.38500,2.3836573) (1.39000,2.3789074) (1.39500,2.3741559) (1.40000,2.3694029) (1.40500,2.3646485) (1.41000,2.3598925) (1.41500,2.3551352) (1.42000,2.3503764) (1.42500,2.3456162) (1.43000,2.3408548) (1.43500,2.3360920) (1.44000,2.3313279) (1.44500,2.3265625) (1.45000,2.3217960) (1.45500,2.3170282) (1.46000,2.3122593) (1.46500,2.3074892) (1.47000,2.3027180) (1.47500,2.2979457) (1.48000,2.2931724) (1.48500,2.2883980) (1.49000,2.2836227) (1.49500,2.2788463) (1.50000,2.2740690) (1.50500,2.2692908) (1.51000,2.2645117) (1.51500,2.2597318) (1.52000,2.2549510) (1.52500,2.2501694) (1.53000,2.2453870) (1.53500,2.2406039) (1.54000,2.2358200) (1.54500,2.2310354) (1.55000,2.2262502) (1.55500,2.2214643) (1.56000,2.2166778) (1.56500,2.2118907) (1.57000,2.2071030) (1.57500,2.2023148) (1.58000,2.1975261) (1.58500,2.1927369) (1.59000,2.1879472) (1.59500,2.1831571) (1.60000,2.1783666) (1.60500,2.1735757) (1.61000,2.1687845) (1.61500,2.1639929) (1.62000,2.1592010) (1.62500,2.1544089) (1.63000,2.1496165) (1.63500,2.1448239) (1.64000,2.1400311) (1.64500,2.1352381) (1.65000,2.1304450) (1.65500,2.1256517) (1.66000,2.1208584) (1.66500,2.1160650) (1.67000,2.1112716) (1.67500,2.1064781) (1.68000,2.1016847) (1.68500,2.0968913) (1.69000,2.0920980) (1.69500,2.0873048) (1.70000,2.0825116) (1.70500,2.0777187) (1.71000,2.0729259) (1.71500,2.0681333) (1.72000,2.0633410) (1.72500,2.0585488) (1.73000,2.0537570) (1.73500,2.0489655) (1.74000,2.0441743) (1.74500,2.0393835) (1.75000,2.0345931) (1.75500,2.0298030) (1.76000,2.0250134) (1.76500,2.0202243) (1.77000,2.0154357) (1.77500,2.0106476) (1.78000,2.0058601) (1.78500,2.0010731) (1.79000,1.9962867) (1.79500,1.9915010) (1.80000,1.9867159) (1.80500,1.9819315) (1.81000,1.9771479) (1.81500,1.9723649) (1.82000,1.9675828) (1.82500,1.9628014) (1.83000,1.9580209) (1.83500,1.9532412) (1.84000,1.9484624) (1.84500,1.9436845) (1.85000,1.9389076) (1.85500,1.9341317) (1.86000,1.9293567) (1.86500,1.9245828) (1.87000,1.9198099) (1.87500,1.9150381) (1.88000,1.9102674) (1.88500,1.9054979) (1.89000,1.9007295) (1.89500,1.8959624) (1.90000,1.8911965) (1.90500,1.8864319) (1.91000,1.8816685) (1.91500,1.8769065) (1.92000,1.8721458) (1.92500,1.8673866) (1.93000,1.8626287) (1.93500,1.8578723) (1.94000,1.8531174) (1.94500,1.8483640) (1.95000,1.8436122) (1.95500,1.8388619) (1.96000,1.8341133) (1.96500,1.8293663) (1.97000,1.8246209) (1.97500,1.8198773) (1.98000,1.8151354) (1.98500,1.8103953) (1.99000,1.8056571) (1.99500,1.8009206) (2.00000,1.7961860) (2.00500,1.7914534) (2.01000,1.7867227) (2.01500,1.7819940) (2.02000,1.7772673) (2.02500,1.7725426) (2.03000,1.7678200) (2.03500,1.7630996) (2.04000,1.7583813) (2.04500,1.7536652) (2.05000,1.7489514) (2.05500,1.7442398) (2.06000,1.7395306) (2.06500,1.7348236) (2.07000,1.7301191) (2.07500,1.7254170) (2.08000,1.7207173) (2.08500,1.7160202) (2.09000,1.7113256) (2.09500,1.7066335) (2.10000,1.7019441) (2.10500,1.6972574) (2.11000,1.6925733) (2.11500,1.6878920) (2.12000,1.6832134) (2.12500,1.6785377) (2.13000,1.6738649) (2.13500,1.6691949) (2.14000,1.6645279) (2.14500,1.6598639) (2.15000,1.6552030) (2.15500,1.6505451) (2.16000,1.6458903) (2.16500,1.6412387) (2.17000,1.6365903) (2.17500,1.6319452) (2.18000,1.6273034) (2.18500,1.6226649) (2.19000,1.6180299) (2.19500,1.6133983) (2.20000,1.6087701) (2.20500,1.6041455) (2.21000,1.5995245) (2.21500,1.5949072) (2.22000,1.5902935) (2.22500,1.5856836) (2.23000,1.5810774) (2.23500,1.5764751) (2.24000,1.5718767) (2.24500,1.5672822) (2.25000,1.5626917) (2.25500,1.5581053) (2.26000,1.5535230) (2.26500,1.5489448) (2.27000,1.5443708) (2.27500,1.5398011) (2.28000,1.5352357) (2.28500,1.5306747) (2.29000,1.5261181) (2.29500,1.5215660) (2.30000,1.5170185) (2.30500,1.5124756) (2.31000,1.5079373) (2.31500,1.5034037) (2.32000,1.4988750) (2.32500,1.4943511) (2.33000,1.4898320) (2.33500,1.4853180) (2.34000,1.4808090) (2.34500,1.4763051) (2.35000,1.4718063) (2.35500,1.4673128) (2.36000,1.4628246) (2.36500,1.4583417) (2.37000,1.4538642) (2.37500,1.4493923) (2.38000,1.4449259) (2.38500,1.4404651) (2.39000,1.4360100) (2.39500,1.4315607) (2.40000,1.4271172) (2.40500,1.4226797) (2.41000,1.4182481) (2.41500,1.4138226) (2.42000,1.4094032) (2.42500,1.4049900) (2.43000,1.4005832) (2.43500,1.3961826) (2.44000,1.3917885) (2.44500,1.3874010) (2.45000,1.3830200) (2.45500,1.3786457) (2.46000,1.3742782) (2.46500,1.3699175) (2.47000,1.3655638) (2.47500,1.3612170) (2.48000,1.3568774) (2.48500,1.3525449) (2.49000,1.3482198) (2.49500,1.3439019) (2.50000,1.3395915) (2.50500,1.3352887) (2.51000,1.3309935) (2.51500,1.3267060) (2.52000,1.3224263) (2.52500,1.3181546) (2.53000,1.3138908) (2.53500,1.3096351) (2.54000,1.3053877) (2.54500,1.3011486) (2.55000,1.2969178) (2.55500,1.2926956) (2.56000,1.2884819) (2.56500,1.2842770) (2.57000,1.2800809) (2.57500,1.2758937) (2.58000,1.2717156) (2.58500,1.2675466) (2.59000,1.2633868) (2.59500,1.2592365) (2.60000,1.2550956) (2.60500,1.2509643) (2.61000,1.2468427) (2.61500,1.2427309) (2.62000,1.2386291) (2.62500,1.2345374) (2.63000,1.2304558) (2.63500,1.2263846) (2.64000,1.2223239) (2.64500,1.2182737) (2.65000,1.2142342) (2.65500,1.2102055) (2.66000,1.2061879) (2.66500,1.2021813) (2.67000,1.1981860) (2.67500,1.1942020) (2.68000,1.1902296) (2.68500,1.1862689) (2.69000,1.1823199) (2.69500,1.1783830) (2.70000,1.1744581) (2.70500,1.1705455) (2.71000,1.1666453) (2.71500,1.1627577) (2.72000,1.1588828) (2.72500,1.1550208) (2.73000,1.1511718) (2.73500,1.1473361) (2.74000,1.1435137) (2.74500,1.1397049) (2.75000,1.1359098) (2.75500,1.1321287) (2.76000,1.1283616) (2.76500,1.1246088) (2.77000,1.1208704) (2.77500,1.1171467) (2.78000,1.1134378) (2.78500,1.1097439) (2.79000,1.1060652) (2.79500,1.1024020) (2.80000,1.0987543) (2.80500,1.0951226) (2.81000,1.0915068) (2.81500,1.0879073) (2.82000,1.0843242) (2.82500,1.0807579) (2.83000,1.0772085) (2.83500,1.0736762) (2.84000,1.0701613) (2.84500,1.0666640) (2.85000,1.0631846) (2.85500,1.0597233) (2.86000,1.0562804) (2.86500,1.0528560) (2.87000,1.0494506) (2.87500,1.0460643) (2.88000,1.0426974) (2.88500,1.0393502) (2.89000,1.0360230) (2.89500,1.0327160) (2.90000,1.0294296) (2.90500,1.0261641) (2.91000,1.0229198) (2.91500,1.0196969) (2.92000,1.0164958) (2.92500,1.0133169) (2.93000,1.0101604) (2.93500,1.0070267) (2.94000,1.0039162) (2.94500,1.0008291) (2.95000,0.9977660) (2.95500,0.9947271) (2.96000,0.9917128) (2.96500,0.9887236) (2.97000,0.9857598) (2.97500,0.9828218) (2.98000,0.9799101) (2.98500,0.9770250) (2.99000,0.9741671) (2.99500,0.9713368) (3.00000,0.9685345) (3.00500,0.9657608) (3.01000,0.9630161) (3.01500,0.9603009) (3.02000,0.9576157) (3.02500,0.9549611) (3.03000,0.9523376) (3.03500,0.9497457) (3.04000,0.9471861) (3.04500,0.9446594) (3.05000,0.9421661) (3.05500,0.9397069) (3.06000,0.9372824) (3.06500,0.9348933) (3.07000,0.9325403) (3.07500,0.9302241) (3.08000,0.9279454) (3.08500,0.9257049) (3.09000,0.9235035) (3.09500,0.9213420) (3.10000,0.9192212) (3.10500,0.9171419) (3.11000,0.9151050) (3.11500,0.9131115) (3.12000,0.9111622) (3.12500,0.9092583) (3.13000,0.9074006) (3.13500,0.9055903) (3.14000,0.9038285) (3.14500,0.9021162) (3.15000,0.9004546) (3.15500,0.8988451) (3.16000,0.8972888) (3.16500,0.8957870) (3.17000,0.8943412) (3.17500,0.8929527) (3.18000,0.8916231) (3.18500,0.8903539) (3.19000,0.8891467) (3.19500,0.8880032) (3.20000,0.8869251) (3.20500,0.8859142) (3.21000,0.8849726) (3.21500,0.8841021) (3.22000,0.8833049) (3.22500,0.8825831) (3.23000,0.8819391) (3.23500,0.8813752) (3.24000,0.8808940) (3.24500,0.8804980) (3.25000,0.8801902) (3.25500,0.8799733) (3.26000,0.8798505) (3.26500,0.8798250) (3.27000,0.8799002) (3.27500,0.8800796) (3.28000,0.8803672) (3.28500,0.8807668) (3.29000,0.8812827) (3.29500,0.8819194) (3.30000,0.8826816) (3.30500,0.8835744) (3.31000,0.8846030) (3.31500,0.8857733) (3.32000,0.8870912) (3.32500,0.8885631) (3.33000,0.8901960) (3.33500,0.8919972) (3.34000,0.8939745) (3.34500,0.8961364) (3.35000,0.8984917) (3.35500,0.9010504) (3.36000,0.9038226) (3.36500,0.9068198) (3.37000,0.9100539) (3.37500,0.9135381) (3.38000,0.9172867) (3.38500,0.9213150) (3.39000,0.9256398) (3.39500,0.9302795) (3.40000,0.9352540) (3.40500,0.9405852) (3.41000,0.9462972) (3.41500,0.9524165) (3.42000,0.9589726) (3.42500,0.9659978) (3.43000,0.9735283) (3.43500,0.9816044) (3.44000,0.9902715) (3.44500,0.9995804) (3.45000,1.0095886) (3.45500,1.0203617) (3.46000,1.0319742) (3.46500,1.0445122) (3.47000,1.0580749) (3.47500,1.0727781) (3.48000,1.0887576) (3.48500,1.1061742) (3.49000,1.1252206) (3.49500,1.1461294) (3.50000,1.1691857) (3.50500,1.1947437) (3.51000,1.2232510) (3.51500,1.2552840) (3.52000,1.2916031) (3.52500,1.3332391) (3.53000,1.3816393) (3.53500,1.4389263) (3.54000,1.5083948) (3.54500,1.5955733) (3.55000,1.7108468) (3.55500,1.8776025) (3.56000,2.1709685) (3.56500,nan) (3.57000,nan) (3.57500,nan) (3.58000,nan) (3.58500,nan) (3.59000,nan) (3.59500,nan) (3.60000,nan) (3.60500,nan) (3.61000,nan) (3.61500,nan) (3.62000,nan) (3.62500,nan) (3.63000,nan) (3.63500,nan) (3.64000,nan) (3.64500,nan) (3.65000,nan) (3.65500,nan) (3.66000,nan) (3.66500,nan) (3.67000,nan) (3.67500,nan) (3.68000,nan) (3.68500,nan) (3.69000,nan) (3.69500,nan) (3.70000,nan) (3.70500,nan) (3.71000,nan) (3.71500,nan) (3.72000,nan) (3.72500,nan) (3.73000,nan) (3.73500,nan) (3.74000,nan) (3.74500,nan) (3.75000,nan) (3.75500,nan) (3.76000,nan) (3.76500,nan) (3.77000,nan) (3.77500,nan) (3.78000,nan) (3.78500,nan) (3.79000,nan) (3.79500,nan) (3.80000,nan) (3.80500,nan) (3.81000,nan) (3.81500,nan) (3.82000,nan) (3.82500,nan) (3.83000,nan) (3.83500,nan) (3.84000,-1.9451457) (3.84500,-1.3353474) (3.85000,-1.0993939) (3.85500,-0.9528925) (3.86000,-0.8476560) (3.86500,-0.7662446) (3.87000,-0.7003461) (3.87500,-0.6453439) (3.88000,-0.5984071) (3.88500,-0.5576743) (3.89000,-0.5218565) (3.89500,-0.4900241) (3.90000,-0.4614847) (3.90500,-0.4357093) (3.91000,-0.4122843) (3.91500,-0.3908804) (3.92000,-0.3712315) (3.92500,-0.3531191) (3.93000,-0.3363621) (3.93500,-0.3208085) (3.94000,-0.3063296) (3.94500,-0.2928155) (3.95000,-0.2801717) (3.95500,-0.2683161) (3.96000,-0.2571775) (3.96500,-0.2466932) (3.97000,-0.2368079) (3.97500,-0.2274729) (3.98000,-0.2186446) (3.98500,-0.2102843) (3.99000,-0.2023570) (3.99500,-0.1948316) (4.00000,-0.1876796)};\addlegendentry{R}
\addplot[thick,orange] coordinates {(-2.00000,nan) (-1.99500,nan) (-1.99000,nan) (-1.98500,nan) (-1.98000,nan) (-1.97500,nan) (-1.97000,nan) (-1.96500,nan) (-1.96000,nan) (-1.95500,nan) (-1.95000,nan) (-1.94500,nan) (-1.94000,nan) (-1.93500,nan) (-1.93000,nan) (-1.92500,nan) (-1.92000,nan) (-1.91500,nan) (-1.91000,nan) (-1.90500,nan) (-1.90000,nan) (-1.89500,nan) (-1.89000,nan) (-1.88500,nan) (-1.88000,nan) (-1.87500,nan) (-1.87000,nan) (-1.86500,nan) (-1.86000,nan) (-1.85500,nan) (-1.85000,nan) (-1.84500,nan) (-1.84000,nan) (-1.83500,nan) (-1.83000,nan) (-1.82500,nan) (-1.82000,nan) (-1.81500,nan) (-1.81000,nan) (-1.80500,nan) (-1.80000,nan) (-1.79500,nan) (-1.79000,nan) (-1.78500,nan) (-1.78000,nan) (-1.77500,nan) (-1.77000,nan) (-1.76500,nan) (-1.76000,nan) (-1.75500,nan) (-1.75000,nan) (-1.74500,nan) (-1.74000,nan) (-1.73500,nan) (-1.73000,nan) (-1.72500,nan) (-1.72000,nan) (-1.71500,nan) (-1.71000,nan) (-1.70500,nan) (-1.70000,nan) (-1.69500,nan) (-1.69000,nan) (-1.68500,nan) (-1.68000,nan) (-1.67500,nan) (-1.67000,nan) (-1.66500,nan) (-1.66000,nan) (-1.65500,nan) (-1.65000,nan) (-1.64500,nan) (-1.64000,nan) (-1.63500,nan) (-1.63000,nan) (-1.62500,nan) (-1.62000,nan) (-1.61500,nan) (-1.61000,nan) (-1.60500,nan) (-1.60000,nan) (-1.59500,nan) (-1.59000,nan) (-1.58500,nan) (-1.58000,nan) (-1.57500,nan) (-1.57000,nan) (-1.56500,nan) (-1.56000,nan) (-1.55500,nan) (-1.55000,nan) (-1.54500,nan) (-1.54000,nan) (-1.53500,nan) (-1.53000,nan) (-1.52500,nan) (-1.52000,nan) (-1.51500,nan) (-1.51000,nan) (-1.50500,nan) (-1.50000,nan) (-1.49500,nan) (-1.49000,nan) (-1.48500,nan) (-1.48000,nan) (-1.47500,nan) (-1.47000,nan) (-1.46500,nan) (-1.46000,nan) (-1.45500,nan) (-1.45000,nan) (-1.44500,nan) (-1.44000,nan) (-1.43500,nan) (-1.43000,nan) (-1.42500,nan) (-1.42000,nan) (-1.41500,nan) (-1.41000,nan) (-1.40500,nan) (-1.40000,nan) (-1.39500,nan) (-1.39000,nan) (-1.38500,nan) (-1.38000,nan) (-1.37500,nan) (-1.37000,nan) (-1.36500,nan) (-1.36000,nan) (-1.35500,nan) (-1.35000,nan) (-1.34500,nan) (-1.34000,nan) (-1.33500,nan) (-1.33000,nan) (-1.32500,nan) (-1.32000,nan) (-1.31500,nan) (-1.31000,nan) (-1.30500,nan) (-1.30000,nan) (-1.29500,nan) (-1.29000,nan) (-1.28500,nan) (-1.28000,nan) (-1.27500,nan) (-1.27000,nan) (-1.26500,nan) (-1.26000,nan) (-1.25500,nan) (-1.25000,nan) (-1.24500,nan) (-1.24000,nan) (-1.23500,nan) (-1.23000,nan) (-1.22500,nan) (-1.22000,nan) (-1.21500,nan) (-1.21000,nan) (-1.20500,nan) (-1.20000,nan) (-1.19500,nan) (-1.19000,nan) (-1.18500,nan) (-1.18000,nan) (-1.17500,nan) (-1.17000,nan) (-1.16500,nan) (-1.16000,nan) (-1.15500,nan) (-1.15000,nan) (-1.14500,nan) (-1.14000,nan) (-1.13500,nan) (-1.13000,nan) (-1.12500,nan) (-1.12000,nan) (-1.11500,nan) (-1.11000,nan) (-1.10500,nan) (-1.10000,nan) (-1.09500,nan) (-1.09000,nan) (-1.08500,nan) (-1.08000,nan) (-1.07500,nan) (-1.07000,nan) (-1.06500,nan) (-1.06000,nan) (-1.05500,nan) (-1.05000,nan) (-1.04500,nan) (-1.04000,nan) (-1.03500,nan) (-1.03000,nan) (-1.02500,nan) (-1.02000,nan) (-1.01500,nan) (-1.01000,nan) (-1.00500,nan) (-1.00000,nan) (-0.99500,nan) (-0.99000,nan) (-0.98500,nan) (-0.98000,nan) (-0.97500,nan) (-0.97000,nan) (-0.96500,nan) (-0.96000,nan) (-0.95500,nan) (-0.95000,nan) (-0.94500,nan) (-0.94000,nan) (-0.93500,nan) (-0.93000,nan) (-0.92500,nan) (-0.92000,nan) (-0.91500,nan) (-0.91000,nan) (-0.90500,nan) (-0.90000,nan) (-0.89500,nan) (-0.89000,nan) (-0.88500,nan) (-0.88000,nan) (-0.87500,nan) (-0.87000,nan) (-0.86500,nan) (-0.86000,nan) (-0.85500,nan) (-0.85000,nan) (-0.84500,nan) (-0.84000,nan) (-0.83500,nan) (-0.83000,nan) (-0.82500,nan) (-0.82000,nan) (-0.81500,nan) (-0.81000,nan) (-0.80500,nan) (-0.80000,nan) (-0.79500,nan) (-0.79000,nan) (-0.78500,nan) (-0.78000,nan) (-0.77500,nan) (-0.77000,nan) (-0.76500,nan) (-0.76000,nan) (-0.75500,nan) (-0.75000,nan) (-0.74500,nan) (-0.74000,nan) (-0.73500,nan) (-0.73000,nan) (-0.72500,nan) (-0.72000,nan) (-0.71500,nan) (-0.71000,nan) (-0.70500,nan) (-0.70000,nan) (-0.69500,nan) (-0.69000,nan) (-0.68500,nan) (-0.68000,nan) (-0.67500,nan) (-0.67000,nan) (-0.66500,nan) (-0.66000,nan) (-0.65500,nan) (-0.65000,nan) (-0.64500,nan) (-0.64000,nan) (-0.63500,nan) (-0.63000,nan) (-0.62500,nan) (-0.62000,nan) (-0.61500,nan) (-0.61000,nan) (-0.60500,nan) (-0.60000,nan) (-0.59500,nan) (-0.59000,nan) (-0.58500,nan) (-0.58000,nan) (-0.57500,nan) (-0.57000,nan) (-0.56500,nan) (-0.56000,nan) (-0.55500,nan) (-0.55000,nan) (-0.54500,nan) (-0.54000,nan) (-0.53500,nan) (-0.53000,nan) (-0.52500,nan) (-0.52000,nan) (-0.51500,nan) (-0.51000,nan) (-0.50500,nan) (-0.50000,nan) (-0.49500,nan) (-0.49000,nan) (-0.48500,nan) (-0.48000,nan) (-0.47500,nan) (-0.47000,nan) (-0.46500,nan) (-0.46000,nan) (-0.45500,nan) (-0.45000,nan) (-0.44500,nan) (-0.44000,nan) (-0.43500,nan) (-0.43000,nan) (-0.42500,nan) (-0.42000,nan) (-0.41500,nan) (-0.41000,nan) (-0.40500,nan) (-0.40000,nan) (-0.39500,nan) (-0.39000,nan) (-0.38500,nan) (-0.38000,nan) (-0.37500,nan) (-0.37000,nan) (-0.36500,nan) (-0.36000,nan) (-0.35500,nan) (-0.35000,nan) (-0.34500,nan) (-0.34000,nan) (-0.33500,nan) (-0.33000,nan) (-0.32500,nan) (-0.32000,nan) (-0.31500,nan) (-0.31000,nan) (-0.30500,nan) (-0.30000,nan) (-0.29500,nan) (-0.29000,nan) (-0.28500,nan) (-0.28000,nan) (-0.27500,nan) (-0.27000,nan) (-0.26500,nan) (-0.26000,nan) (-0.25500,nan) (-0.25000,nan) (-0.24500,nan) (-0.24000,nan) (-0.23500,nan) (-0.23000,nan) (-0.22500,nan) (-0.22000,nan) (-0.21500,nan) (-0.21000,nan) (-0.20500,nan) (-0.20000,nan) (-0.19500,nan) (-0.19000,nan) (-0.18500,nan) (-0.18000,nan) (-0.17500,nan) (-0.17000,nan) (-0.16500,nan) (-0.16000,nan) (-0.15500,nan) (-0.15000,nan) (-0.14500,nan) (-0.14000,nan) (-0.13500,nan) (-0.13000,nan) (-0.12500,nan) (-0.12000,nan) (-0.11500,nan) (-0.11000,nan) (-0.10500,nan) (-0.10000,nan) (-0.09500,nan) (-0.09000,nan) (-0.08500,nan) (-0.08000,nan) (-0.07500,nan) (-0.07000,nan) (-0.06500,nan) (-0.06000,nan) (-0.05500,nan) (-0.05000,nan) (-0.04500,nan) (-0.04000,nan) (-0.03500,nan) (-0.03000,nan) (-0.02500,nan) (-0.02000,nan) (-0.01500,nan) (-0.01000,nan) (-0.00500,nan) (0.00000,nan) (0.00500,nan) (0.01000,nan) (0.01500,nan) (0.02000,nan) (0.02500,nan) (0.03000,nan) (0.03500,nan) (0.04000,nan) (0.04500,nan) (0.05000,nan) (0.05500,nan) (0.06000,nan) (0.06500,nan) (0.07000,nan) (0.07500,nan) (0.08000,nan) (0.08500,nan) (0.09000,nan) (0.09500,nan) (0.10000,nan) (0.10500,nan) (0.11000,nan) (0.11500,nan) (0.12000,nan) (0.12500,nan) (0.13000,nan) (0.13500,nan) (0.14000,nan) (0.14500,nan) (0.15000,nan) (0.15500,nan) (0.16000,nan) (0.16500,nan) (0.17000,nan) (0.17500,nan) (0.18000,nan) (0.18500,nan) (0.19000,nan) (0.19500,nan) (0.20000,nan) (0.20500,nan) (0.21000,nan) (0.21500,nan) (0.22000,nan) (0.22500,nan) (0.23000,nan) (0.23500,nan) (0.24000,nan) (0.24500,nan) (0.25000,nan) (0.25500,nan) (0.26000,nan) (0.26500,nan) (0.27000,nan) (0.27500,nan) (0.28000,nan) (0.28500,nan) (0.29000,nan) (0.29500,nan) (0.30000,nan) (0.30500,nan) (0.31000,nan) (0.31500,nan) (0.32000,nan) (0.32500,nan) (0.33000,nan) (0.33500,nan) (0.34000,nan) (0.34500,nan) (0.35000,nan) (0.35500,nan) (0.36000,nan) (0.36500,nan) (0.37000,nan) (0.37500,nan) (0.38000,nan) (0.38500,nan) (0.39000,nan) (0.39500,nan) (0.40000,nan) (0.40500,nan) (0.41000,nan) (0.41500,nan) (0.42000,nan) (0.42500,nan) (0.43000,nan) (0.43500,nan) (0.44000,nan) (0.44500,nan) (0.45000,nan) (0.45500,nan) (0.46000,nan) (0.46500,nan) (0.47000,nan) (0.47500,nan) (0.48000,nan) (0.48500,nan) (0.49000,nan) (0.49500,nan) (0.50000,nan) (0.50500,nan) (0.51000,nan) (0.51500,nan) (0.52000,nan) (0.52500,nan) (0.53000,nan) (0.53500,nan) (0.54000,nan) (0.54500,nan) (0.55000,nan) (0.55500,nan) (0.56000,nan) (0.56500,nan) (0.57000,nan) (0.57500,nan) (0.58000,nan) (0.58500,nan) (0.59000,nan) (0.59500,nan) (0.60000,nan) (0.60500,nan) (0.61000,nan) (0.61500,nan) (0.62000,nan) (0.62500,nan) (0.63000,nan) (0.63500,nan) (0.64000,nan) (0.64500,nan) (0.65000,nan) (0.65500,nan) (0.66000,nan) (0.66500,nan) (0.67000,nan) (0.67500,nan) (0.68000,nan) (0.68500,nan) (0.69000,nan) (0.69500,nan) (0.70000,nan) (0.70500,nan) (0.71000,nan) (0.71500,nan) (0.72000,nan) (0.72500,nan) (0.73000,nan) (0.73500,nan) (0.74000,nan) (0.74500,nan) (0.75000,nan) (0.75500,nan) (0.76000,nan) (0.76500,nan) (0.77000,nan) (0.77500,nan) (0.78000,nan) (0.78500,nan) (0.79000,nan) (0.79500,nan) (0.80000,nan) (0.80500,nan) (0.81000,nan) (0.81500,nan) (0.82000,nan) (0.82500,nan) (0.83000,nan) (0.83500,nan) (0.84000,nan) (0.84500,nan) (0.85000,nan) (0.85500,nan) (0.86000,nan) (0.86500,nan) (0.87000,nan) (0.87500,nan) (0.88000,nan) (0.88500,nan) (0.89000,nan) (0.89500,nan) (0.90000,nan) (0.90500,nan) (0.91000,nan) (0.91500,nan) (0.92000,nan) (0.92500,nan) (0.93000,nan) (0.93500,nan) (0.94000,nan) (0.94500,nan) (0.95000,nan) (0.95500,nan) (0.96000,nan) (0.96500,nan) (0.97000,nan) (0.97500,nan) (0.98000,nan) (0.98500,nan) (0.99000,nan) (0.99500,nan) (1.00000,nan) (1.00500,nan) (1.01000,nan) (1.01500,nan) (1.02000,nan) (1.02500,nan) (1.03000,nan) (1.03500,nan) (1.04000,nan) (1.04500,nan) (1.05000,nan) (1.05500,nan) (1.06000,nan) (1.06500,nan) (1.07000,nan) (1.07500,nan) (1.08000,nan) (1.08500,nan) (1.09000,nan) (1.09500,nan) (1.10000,nan) (1.10500,nan) (1.11000,nan) (1.11500,nan) (1.12000,nan) (1.12500,nan) (1.13000,nan) (1.13500,nan) (1.14000,nan) (1.14500,nan) (1.15000,nan) (1.15500,nan) (1.16000,nan) (1.16500,nan) (1.17000,nan) (1.17500,nan) (1.18000,nan) (1.18500,nan) (1.19000,nan) (1.19500,nan) (1.20000,nan) (1.20500,nan) (1.21000,nan) (1.21500,nan) (1.22000,nan) (1.22500,nan) (1.23000,nan) (1.23500,nan) (1.24000,nan) (1.24500,nan) (1.25000,nan) (1.25500,nan) (1.26000,nan) (1.26500,nan) (1.27000,nan) (1.27500,nan) (1.28000,nan) (1.28500,nan) (1.29000,nan) (1.29500,nan) (1.30000,nan) (1.30500,nan) (1.31000,nan) (1.31500,nan) (1.32000,nan) (1.32500,nan) (1.33000,nan) (1.33500,nan) (1.34000,nan) (1.34500,nan) (1.35000,nan) (1.35500,nan) (1.36000,nan) (1.36500,nan) (1.37000,nan) (1.37500,nan) (1.38000,nan) (1.38500,nan) (1.39000,nan) (1.39500,nan) (1.40000,nan) (1.40500,nan) (1.41000,nan) (1.41500,nan) (1.42000,nan) (1.42500,nan) (1.43000,nan) (1.43500,nan) (1.44000,nan) (1.44500,nan) (1.45000,nan) (1.45500,nan) (1.46000,nan) (1.46500,nan) (1.47000,nan) (1.47500,nan) (1.48000,nan) (1.48500,nan) (1.49000,nan) (1.49500,nan) (1.50000,nan) (1.50500,nan) (1.51000,nan) (1.51500,nan) (1.52000,nan) (1.52500,nan) (1.53000,nan) (1.53500,nan) (1.54000,nan) (1.54500,nan) (1.55000,nan) (1.55500,nan) (1.56000,nan) (1.56500,nan) (1.57000,nan) (1.57500,nan) (1.58000,nan) (1.58500,nan) (1.59000,nan) (1.59500,nan) (1.60000,nan) (1.60500,nan) (1.61000,nan) (1.61500,nan) (1.62000,nan) (1.62500,nan) (1.63000,nan) (1.63500,nan) (1.64000,nan) (1.64500,nan) (1.65000,nan) (1.65500,nan) (1.66000,nan) (1.66500,nan) (1.67000,nan) (1.67500,nan) (1.68000,nan) (1.68500,nan) (1.69000,nan) (1.69500,nan) (1.70000,nan) (1.70500,nan) (1.71000,nan) (1.71500,nan) (1.72000,nan) (1.72500,nan) (1.73000,nan) (1.73500,nan) (1.74000,nan) (1.74500,nan) (1.75000,nan) (1.75500,nan) (1.76000,nan) (1.76500,nan) (1.77000,nan) (1.77500,nan) (1.78000,nan) (1.78500,nan) (1.79000,nan) (1.79500,nan) (1.80000,nan) (1.80500,nan) (1.81000,nan) (1.81500,nan) (1.82000,nan) (1.82500,nan) (1.83000,nan) (1.83500,nan) (1.84000,nan) (1.84500,nan) (1.85000,nan) (1.85500,nan) (1.86000,nan) (1.86500,nan) (1.87000,nan) (1.87500,nan) (1.88000,nan) (1.88500,nan) (1.89000,nan) (1.89500,nan) (1.90000,nan) (1.90500,nan) (1.91000,nan) (1.91500,nan) (1.92000,nan) (1.92500,nan) (1.93000,nan) (1.93500,nan) (1.94000,nan) (1.94500,nan) (1.95000,nan) (1.95500,nan) (1.96000,nan) (1.96500,nan) (1.97000,nan) (1.97500,nan) (1.98000,nan) (1.98500,nan) (1.99000,nan) (1.99500,nan) (2.00000,nan) (2.00500,nan) (2.01000,nan) (2.01500,nan) (2.02000,nan) (2.02500,nan) (2.03000,nan) (2.03500,nan) (2.04000,nan) (2.04500,nan) (2.05000,nan) (2.05500,nan) (2.06000,nan) (2.06500,nan) (2.07000,nan) (2.07500,nan) (2.08000,nan) (2.08500,nan) (2.09000,nan) (2.09500,nan) (2.10000,nan) (2.10500,nan) (2.11000,nan) (2.11500,nan) (2.12000,nan) (2.12500,nan) (2.13000,nan) (2.13500,nan) (2.14000,nan) (2.14500,nan) (2.15000,nan) (2.15500,nan) (2.16000,nan) (2.16500,nan) (2.17000,nan) (2.17500,nan) (2.18000,nan) (2.18500,nan) (2.19000,nan) (2.19500,nan) (2.20000,nan) (2.20500,nan) (2.21000,nan) (2.21500,nan) (2.22000,nan) (2.22500,nan) (2.23000,nan) (2.23500,nan) (2.24000,nan) (2.24500,nan) (2.25000,nan) (2.25500,nan) (2.26000,nan) (2.26500,nan) (2.27000,nan) (2.27500,nan) (2.28000,nan) (2.28500,nan) (2.29000,nan) (2.29500,nan) (2.30000,nan) (2.30500,nan) (2.31000,nan) (2.31500,nan) (2.32000,nan) (2.32500,nan) (2.33000,nan) (2.33500,nan) (2.34000,nan) (2.34500,nan) (2.35000,nan) (2.35500,nan) (2.36000,nan) (2.36500,nan) (2.37000,nan) (2.37500,nan) (2.38000,nan) (2.38500,nan) (2.39000,nan) (2.39500,nan) (2.40000,nan) (2.40500,nan) (2.41000,nan) (2.41500,nan) (2.42000,nan) (2.42500,nan) (2.43000,nan) (2.43500,nan) (2.44000,nan) (2.44500,nan) (2.45000,nan) (2.45500,nan) (2.46000,nan) (2.46500,nan) (2.47000,nan) (2.47500,nan) (2.48000,nan) (2.48500,nan) (2.49000,nan) (2.49500,nan) (2.50000,nan) (2.50500,nan) (2.51000,nan) (2.51500,nan) (2.52000,nan) (2.52500,nan) (2.53000,nan) (2.53500,nan) (2.54000,nan) (2.54500,nan) (2.55000,nan) (2.55500,nan) (2.56000,nan) (2.56500,nan) (2.57000,nan) (2.57500,nan) (2.58000,nan) (2.58500,nan) (2.59000,nan) (2.59500,nan) (2.60000,nan) (2.60500,nan) (2.61000,nan) (2.61500,nan) (2.62000,nan) (2.62500,nan) (2.63000,nan) (2.63500,nan) (2.64000,nan) (2.64500,nan) (2.65000,nan) (2.65500,nan) (2.66000,nan) (2.66500,nan) (2.67000,nan) (2.67500,nan) (2.68000,nan) (2.68500,nan) (2.69000,nan) (2.69500,nan) (2.70000,nan) (2.70500,nan) (2.71000,nan) (2.71500,nan) (2.72000,nan) (2.72500,nan) (2.73000,nan) (2.73500,nan) (2.74000,nan) (2.74500,nan) (2.75000,nan) (2.75500,nan) (2.76000,nan) (2.76500,nan) (2.77000,nan) (2.77500,nan) (2.78000,nan) (2.78500,nan) (2.79000,nan) (2.79500,nan) (2.80000,nan) (2.80500,nan) (2.81000,nan) (2.81500,nan) (2.82000,nan) (2.82500,nan) (2.83000,nan) (2.83500,nan) (2.84000,nan) (2.84500,nan) (2.85000,nan) (2.85500,nan) (2.86000,nan) (2.86500,nan) (2.87000,nan) (2.87500,nan) (2.88000,nan) (2.88500,nan) (2.89000,nan) (2.89500,nan) (2.90000,nan) (2.90500,nan) (2.91000,nan) (2.91500,nan) (2.92000,nan) (2.92500,nan) (2.93000,nan) (2.93500,nan) (2.94000,nan) (2.94500,nan) (2.95000,nan) (2.95500,nan) (2.96000,nan) (2.96500,nan) (2.97000,nan) (2.97500,nan) (2.98000,nan) (2.98500,nan) (2.99000,nan) (2.99500,nan) (3.00000,nan) (3.00500,nan) (3.01000,nan) (3.01500,nan) (3.02000,nan) (3.02500,nan) (3.03000,nan) (3.03500,nan) (3.04000,nan) (3.04500,nan) (3.05000,nan) (3.05500,nan) (3.06000,nan) (3.06500,nan) (3.07000,nan) (3.07500,nan) (3.08000,nan) (3.08500,nan) (3.09000,nan) (3.09500,nan) (3.10000,nan) (3.10500,nan) (3.11000,nan) (3.11500,nan) (3.12000,nan) (3.12500,nan) (3.13000,nan) (3.13500,nan) (3.14000,nan) (3.14500,nan) (3.15000,nan) (3.15500,nan) (3.16000,nan) (3.16500,nan) (3.17000,nan) (3.17500,nan) (3.18000,nan) (3.18500,nan) (3.19000,nan) (3.19500,nan) (3.20000,nan) (3.20500,nan) (3.21000,nan) (3.21500,nan) (3.22000,nan) (3.22500,nan) (3.23000,nan) (3.23500,nan) (3.24000,nan) (3.24500,nan) (3.25000,nan) (3.25500,nan) (3.26000,nan) (3.26500,nan) (3.27000,nan) (3.27500,nan) (3.28000,nan) (3.28500,nan) (3.29000,nan) (3.29500,nan) (3.30000,nan) (3.30500,nan) (3.31000,nan) (3.31500,nan) (3.32000,nan) (3.32500,nan) (3.33000,nan) (3.33500,nan) (3.34000,nan) (3.34500,nan) (3.35000,nan) (3.35500,nan) (3.36000,nan) (3.36500,nan) (3.37000,nan) (3.37500,nan) (3.38000,nan) (3.38500,nan) (3.39000,nan) (3.39500,nan) (3.40000,nan) (3.40500,nan) (3.41000,nan) (3.41500,nan) (3.42000,nan) (3.42500,nan) (3.43000,nan) (3.43500,nan) (3.44000,nan) (3.44500,nan) (3.45000,nan) (3.45500,nan) (3.46000,nan) (3.46500,nan) (3.47000,nan) (3.47500,nan) (3.48000,nan) (3.48500,nan) (3.49000,nan) (3.49500,nan) (3.50000,nan) (3.50500,nan) (3.51000,nan) (3.51500,-1.6373135) (3.52000,-1.4036336) (3.52500,-1.2544401) (3.53000,-1.1450672) (3.53500,-1.0589587) (3.54000,-0.9881264) (3.54500,-0.9280994) (3.55000,-0.8761208) (3.55500,-0.8303702) (3.56000,-0.7895822) (3.56500,-0.7528418) (3.57000,-0.7194655) (3.57500,-0.6889293) (3.58000,-0.6608228) (3.58500,-0.6348186) (3.59000,-0.6106510) (3.59500,-0.5881017) (3.60000,-0.5669891) (3.60500,-0.5471604) (3.61000,-0.5284859) (3.61500,-0.5108546) (3.62000,-0.4941707) (3.62500,-0.4783508) (3.63000,-0.4633223) (3.63500,-0.4490211) (3.64000,-0.4353906) (3.64500,-0.4223805) (3.65000,-0.4099460) (3.65500,-0.3980468) (3.66000,-0.3866468) (3.66500,-0.3757133) (3.67000,-0.3652167) (3.67500,-0.3551301) (3.68000,-0.3454290) (3.68500,-0.3360909) (3.69000,-0.3270954) (3.69500,-0.3184235) (3.70000,-0.3100579) (3.70500,-0.3019825) (3.71000,-0.2941824) (3.71500,-0.2866440) (3.72000,-0.2793544) (3.72500,-0.2723017) (3.73000,-0.2654750) (3.73500,-0.2588638) (3.74000,-0.2524584) (3.74500,-0.2462499) (3.75000,-0.2402297) (3.75500,-0.2343900) (3.76000,-0.2287231) (3.76500,-0.2232220) (3.77000,-0.2178801) (3.77500,-0.2126912) (3.78000,-0.2076493) (3.78500,-0.2027487) (3.79000,-0.1979843) (3.79500,-0.1933509) (3.80000,-0.1888438) (3.80500,-0.1844585) (3.81000,-0.1801907) (3.81500,-0.1760362) (3.82000,-0.1719913) (3.82500,-0.1680521) (3.83000,-0.1642152) (3.83500,-0.1604771) (3.84000,-0.1568346) (3.84500,-0.1532848) (3.85000,-0.1498245) (3.85500,-0.1464510) (3.86000,-0.1431616) (3.86500,-0.1399537) (3.87000,-0.1368249) (3.87500,-0.1337727) (3.88000,-0.1307948) (3.88500,-0.1278892) (3.89000,-0.1250536) (3.89500,-0.1222860) (3.90000,-0.1195845) (3.90500,-0.1169473) (3.91000,-0.1143724) (3.91500,-0.1118582) (3.92000,-0.1094030) (3.92500,-0.1070052) (3.93000,-0.1046632) (3.93500,-0.1023754) (3.94000,-0.1001406) (3.94500,-0.0979572) (3.95000,-0.0958238) (3.95500,-0.0937393) (3.96000,-0.0917023) (3.96500,-0.0897115) (3.97000,-0.0877659) (3.97500,-0.0858642) (3.98000,-0.0840054) (3.98500,-0.0821883) (3.99000,-0.0804119) (3.99500,-0.0786752) (4.00000,-0.0769772)};\addlegendentry{L}
\addplot[thick,dashed,gray] coordinates {(-2.00000,nan) (-1.99500,nan) (-1.99000,nan) (-1.98500,nan) (-1.98000,nan) (-1.97500,nan) (-1.97000,nan) (-1.96500,nan) (-1.96000,nan) (-1.95500,nan) (-1.95000,nan) (-1.94500,nan) (-1.94000,nan) (-1.93500,nan) (-1.93000,nan) (-1.92500,nan) (-1.92000,nan) (-1.91500,nan) (-1.91000,nan) (-1.90500,nan) (-1.90000,nan) (-1.89500,nan) (-1.89000,nan) (-1.88500,nan) (-1.88000,nan) (-1.87500,nan) (-1.87000,nan) (-1.86500,nan) (-1.86000,nan) (-1.85500,nan) (-1.85000,nan) (-1.84500,nan) (-1.84000,nan) (-1.83500,nan) (-1.83000,nan) (-1.82500,nan) (-1.82000,nan) (-1.81500,nan) (-1.81000,nan) (-1.80500,nan) (-1.80000,nan) (-1.79500,nan) (-1.79000,nan) (-1.78500,nan) (-1.78000,nan) (-1.77500,nan) (-1.77000,nan) (-1.76500,nan) (-1.76000,nan) (-1.75500,nan) (-1.75000,nan) (-1.74500,nan) (-1.74000,nan) (-1.73500,nan) (-1.73000,nan) (-1.72500,nan) (-1.72000,nan) (-1.71500,nan) (-1.71000,nan) (-1.70500,nan) (-1.70000,nan) (-1.69500,nan) (-1.69000,nan) (-1.68500,nan) (-1.68000,nan) (-1.67500,nan) (-1.67000,nan) (-1.66500,nan) (-1.66000,nan) (-1.65500,nan) (-1.65000,nan) (-1.64500,nan) (-1.64000,nan) (-1.63500,nan) (-1.63000,nan) (-1.62500,nan) (-1.62000,nan) (-1.61500,nan) (-1.61000,nan) (-1.60500,nan) (-1.60000,nan) (-1.59500,nan) (-1.59000,nan) (-1.58500,nan) (-1.58000,nan) (-1.57500,nan) (-1.57000,nan) (-1.56500,nan) (-1.56000,nan) (-1.55500,nan) (-1.55000,nan) (-1.54500,nan) (-1.54000,nan) (-1.53500,nan) (-1.53000,nan) (-1.52500,nan) (-1.52000,nan) (-1.51500,nan) (-1.51000,nan) (-1.50500,nan) (-1.50000,nan) (-1.49500,nan) (-1.49000,nan) (-1.48500,nan) (-1.48000,nan) (-1.47500,nan) (-1.47000,nan) (-1.46500,nan) (-1.46000,nan) (-1.45500,nan) (-1.45000,nan) (-1.44500,nan) (-1.44000,nan) (-1.43500,nan) (-1.43000,nan) (-1.42500,nan) (-1.42000,nan) (-1.41500,nan) (-1.41000,nan) (-1.40500,nan) (-1.40000,nan) (-1.39500,nan) (-1.39000,nan) (-1.38500,nan) (-1.38000,nan) (-1.37500,nan) (-1.37000,nan) (-1.36500,nan) (-1.36000,nan) (-1.35500,nan) (-1.35000,nan) (-1.34500,nan) (-1.34000,nan) (-1.33500,nan) (-1.33000,nan) (-1.32500,nan) (-1.32000,nan) (-1.31500,nan) (-1.31000,nan) (-1.30500,nan) (-1.30000,nan) (-1.29500,nan) (-1.29000,nan) (-1.28500,nan) (-1.28000,nan) (-1.27500,nan) (-1.27000,nan) (-1.26500,nan) (-1.26000,nan) (-1.25500,nan) (-1.25000,nan) (-1.24500,nan) (-1.24000,nan) (-1.23500,nan) (-1.23000,nan) (-1.22500,nan) (-1.22000,nan) (-1.21500,nan) (-1.21000,nan) (-1.20500,nan) (-1.20000,nan) (-1.19500,nan) (-1.19000,nan) (-1.18500,nan) (-1.18000,nan) (-1.17500,nan) (-1.17000,nan) (-1.16500,nan) (-1.16000,nan) (-1.15500,nan) (-1.15000,nan) (-1.14500,nan) (-1.14000,nan) (-1.13500,nan) (-1.13000,nan) (-1.12500,nan) (-1.12000,nan) (-1.11500,nan) (-1.11000,nan) (-1.10500,nan) (-1.10000,nan) (-1.09500,nan) (-1.09000,nan) (-1.08500,nan) (-1.08000,nan) (-1.07500,nan) (-1.07000,nan) (-1.06500,nan) (-1.06000,nan) (-1.05500,nan) (-1.05000,nan) (-1.04500,nan) (-1.04000,nan) (-1.03500,nan) (-1.03000,nan) (-1.02500,nan) (-1.02000,nan) (-1.01500,nan) (-1.01000,nan) (-1.00500,nan) (-1.00000,nan) (-0.99500,nan) (-0.99000,nan) (-0.98500,nan) (-0.98000,nan) (-0.97500,nan) (-0.97000,nan) (-0.96500,nan) (-0.96000,nan) (-0.95500,nan) (-0.95000,nan) (-0.94500,nan) (-0.94000,nan) (-0.93500,nan) (-0.93000,nan) (-0.92500,nan) (-0.92000,nan) (-0.91500,nan) (-0.91000,nan) (-0.90500,nan) (-0.90000,nan) (-0.89500,nan) (-0.89000,nan) (-0.88500,nan) (-0.88000,nan) (-0.87500,nan) (-0.87000,nan) (-0.86500,nan) (-0.86000,nan) (-0.85500,nan) (-0.85000,nan) (-0.84500,nan) (-0.84000,nan) (-0.83500,nan) (-0.83000,nan) (-0.82500,nan) (-0.82000,nan) (-0.81500,nan) (-0.81000,nan) (-0.80500,nan) (-0.80000,nan) (-0.79500,nan) (-0.79000,nan) (-0.78500,nan) (-0.78000,nan) (-0.77500,nan) (-0.77000,nan) (-0.76500,nan) (-0.76000,nan) (-0.75500,nan) (-0.75000,nan) (-0.74500,nan) (-0.74000,nan) (-0.73500,nan) (-0.73000,nan) (-0.72500,nan) (-0.72000,nan) (-0.71500,nan) (-0.71000,nan) (-0.70500,nan) (-0.70000,nan) (-0.69500,nan) (-0.69000,nan) (-0.68500,nan) (-0.68000,nan) (-0.67500,nan) (-0.67000,nan) (-0.66500,nan) (-0.66000,nan) (-0.65500,nan) (-0.65000,nan) (-0.64500,nan) (-0.64000,nan) (-0.63500,nan) (-0.63000,nan) (-0.62500,nan) (-0.62000,nan) (-0.61500,nan) (-0.61000,nan) (-0.60500,nan) (-0.60000,nan) (-0.59500,nan) (-0.59000,nan) (-0.58500,nan) (-0.58000,nan) (-0.57500,nan) (-0.57000,nan) (-0.56500,nan) (-0.56000,nan) (-0.55500,nan) (-0.55000,nan) (-0.54500,nan) (-0.54000,nan) (-0.53500,nan) (-0.53000,nan) (-0.52500,nan) (-0.52000,nan) (-0.51500,nan) (-0.51000,nan) (-0.50500,nan) (-0.50000,nan) (-0.49500,nan) (-0.49000,nan) (-0.48500,nan) (-0.48000,nan) (-0.47500,nan) (-0.47000,nan) (-0.46500,nan) (-0.46000,nan) (-0.45500,nan) (-0.45000,nan) (-0.44500,nan) (-0.44000,nan) (-0.43500,nan) (-0.43000,nan) (-0.42500,nan) (-0.42000,nan) (-0.41500,nan) (-0.41000,nan) (-0.40500,nan) (-0.40000,nan) (-0.39500,nan) (-0.39000,nan) (-0.38500,nan) (-0.38000,nan) (-0.37500,nan) (-0.37000,nan) (-0.36500,nan) (-0.36000,nan) (-0.35500,nan) (-0.35000,nan) (-0.34500,nan) (-0.34000,nan) (-0.33500,nan) (-0.33000,nan) (-0.32500,nan) (-0.32000,nan) (-0.31500,nan) (-0.31000,nan) (-0.30500,nan) (-0.30000,nan) (-0.29500,nan) (-0.29000,nan) (-0.28500,nan) (-0.28000,nan) (-0.27500,nan) (-0.27000,nan) (-0.26500,nan) (-0.26000,nan) (-0.25500,nan) (-0.25000,nan) (-0.24500,nan) (-0.24000,nan) (-0.23500,nan) (-0.23000,nan) (-0.22500,nan) (-0.22000,nan) (-0.21500,nan) (-0.21000,nan) (-0.20500,nan) (-0.20000,nan) (-0.19500,nan) (-0.19000,nan) (-0.18500,nan) (-0.18000,nan) (-0.17500,nan) (-0.17000,nan) (-0.16500,nan) (-0.16000,nan) (-0.15500,nan) (-0.15000,nan) (-0.14500,nan) (-0.14000,nan) (-0.13500,nan) (-0.13000,nan) (-0.12500,nan) (-0.12000,nan) (-0.11500,nan) (-0.11000,nan) (-0.10500,nan) (-0.10000,nan) (-0.09500,nan) (-0.09000,nan) (-0.08500,nan) (-0.08000,nan) (-0.07500,nan) (-0.07000,nan) (-0.06500,nan) (-0.06000,nan) (-0.05500,nan) (-0.05000,nan) (-0.04500,nan) (-0.04000,nan) (-0.03500,nan) (-0.03000,nan) (-0.02500,nan) (-0.02000,nan) (-0.01500,nan) (-0.01000,nan) (-0.00500,nan) (0.00000,nan) (0.00500,nan) (0.01000,nan) (0.01500,nan) (0.02000,nan) (0.02500,nan) (0.03000,nan) (0.03500,nan) (0.04000,nan) (0.04500,nan) (0.05000,nan) (0.05500,nan) (0.06000,nan) (0.06500,nan) (0.07000,nan) (0.07500,nan) (0.08000,nan) (0.08500,nan) (0.09000,nan) (0.09500,nan) (0.10000,nan) (0.10500,nan) (0.11000,nan) (0.11500,nan) (0.12000,nan) (0.12500,nan) (0.13000,nan) (0.13500,nan) (0.14000,nan) (0.14500,nan) (0.15000,nan) (0.15500,nan) (0.16000,nan) (0.16500,nan) (0.17000,nan) (0.17500,nan) (0.18000,nan) (0.18500,nan) (0.19000,nan) (0.19500,nan) (0.20000,nan) (0.20500,nan) (0.21000,nan) (0.21500,nan) (0.22000,nan) (0.22500,nan) (0.23000,nan) (0.23500,nan) (0.24000,nan) (0.24500,nan) (0.25000,nan) (0.25500,nan) (0.26000,nan) (0.26500,nan) (0.27000,nan) (0.27500,nan) (0.28000,nan) (0.28500,nan) (0.29000,nan) (0.29500,nan) (0.30000,nan) (0.30500,nan) (0.31000,nan) (0.31500,nan) (0.32000,nan) (0.32500,nan) (0.33000,nan) (0.33500,nan) (0.34000,nan) (0.34500,nan) (0.35000,nan) (0.35500,nan) (0.36000,nan) (0.36500,nan) (0.37000,nan) (0.37500,nan) (0.38000,nan) (0.38500,nan) (0.39000,nan) (0.39500,nan) (0.40000,nan) (0.40500,nan) (0.41000,nan) (0.41500,nan) (0.42000,nan) (0.42500,nan) (0.43000,nan) (0.43500,nan) (0.44000,nan) (0.44500,nan) (0.45000,nan) (0.45500,nan) (0.46000,nan) (0.46500,nan) (0.47000,nan) (0.47500,nan) (0.48000,nan) (0.48500,nan) (0.49000,nan) (0.49500,nan) (0.50000,nan) (0.50500,nan) (0.51000,nan) (0.51500,nan) (0.52000,nan) (0.52500,nan) (0.53000,nan) (0.53500,nan) (0.54000,nan) (0.54500,nan) (0.55000,nan) (0.55500,nan) (0.56000,nan) (0.56500,nan) (0.57000,nan) (0.57500,nan) (0.58000,nan) (0.58500,nan) (0.59000,nan) (0.59500,nan) (0.60000,nan) (0.60500,nan) (0.61000,nan) (0.61500,nan) (0.62000,nan) (0.62500,nan) (0.63000,nan) (0.63500,nan) (0.64000,nan) (0.64500,nan) (0.65000,nan) (0.65500,nan) (0.66000,nan) (0.66500,nan) (0.67000,nan) (0.67500,nan) (0.68000,nan) (0.68500,nan) (0.69000,nan) (0.69500,nan) (0.70000,nan) (0.70500,nan) (0.71000,nan) (0.71500,nan) (0.72000,nan) (0.72500,nan) (0.73000,nan) (0.73500,nan) (0.74000,nan) (0.74500,nan) (0.75000,nan) (0.75500,nan) (0.76000,nan) (0.76500,nan) (0.77000,nan) (0.77500,nan) (0.78000,nan) (0.78500,nan) (0.79000,nan) (0.79500,nan) (0.80000,nan) (0.80500,nan) (0.81000,nan) (0.81500,nan) (0.82000,nan) (0.82500,nan) (0.83000,nan) (0.83500,nan) (0.84000,nan) (0.84500,nan) (0.85000,nan) (0.85500,nan) (0.86000,nan) (0.86500,nan) (0.87000,nan) (0.87500,nan) (0.88000,nan) (0.88500,nan) (0.89000,nan) (0.89500,nan) (0.90000,nan) (0.90500,nan) (0.91000,nan) (0.91500,nan) (0.92000,nan) (0.92500,nan) (0.93000,nan) (0.93500,nan) (0.94000,nan) (0.94500,nan) (0.95000,nan) (0.95500,nan) (0.96000,nan) (0.96500,nan) (0.97000,nan) (0.97500,nan) (0.98000,nan) (0.98500,nan) (0.99000,nan) (0.99500,nan) (1.00000,nan) (1.00500,nan) (1.01000,nan) (1.01500,nan) (1.02000,nan) (1.02500,nan) (1.03000,nan) (1.03500,nan) (1.04000,nan) (1.04500,nan) (1.05000,nan) (1.05500,nan) (1.06000,nan) (1.06500,nan) (1.07000,nan) (1.07500,nan) (1.08000,nan) (1.08500,nan) (1.09000,nan) (1.09500,nan) (1.10000,nan) (1.10500,nan) (1.11000,nan) (1.11500,nan) (1.12000,nan) (1.12500,nan) (1.13000,nan) (1.13500,nan) (1.14000,nan) (1.14500,nan) (1.15000,nan) (1.15500,nan) (1.16000,nan) (1.16500,nan) (1.17000,nan) (1.17500,nan) (1.18000,nan) (1.18500,nan) (1.19000,nan) (1.19500,nan) (1.20000,nan) (1.20500,nan) (1.21000,nan) (1.21500,nan) (1.22000,nan) (1.22500,nan) (1.23000,nan) (1.23500,nan) (1.24000,nan) (1.24500,nan) (1.25000,nan) (1.25500,nan) (1.26000,nan) (1.26500,nan) (1.27000,nan) (1.27500,nan) (1.28000,nan) (1.28500,nan) (1.29000,nan) (1.29500,nan) (1.30000,nan) (1.30500,nan) (1.31000,nan) (1.31500,nan) (1.32000,nan) (1.32500,nan) (1.33000,nan) (1.33500,nan) (1.34000,nan) (1.34500,nan) (1.35000,nan) (1.35500,nan) (1.36000,nan) (1.36500,nan) (1.37000,nan) (1.37500,nan) (1.38000,nan) (1.38500,nan) (1.39000,nan) (1.39500,nan) (1.40000,nan) (1.40500,nan) (1.41000,nan) (1.41500,nan) (1.42000,nan) (1.42500,nan) (1.43000,nan) (1.43500,nan) (1.44000,nan) (1.44500,nan) (1.45000,nan) (1.45500,nan) (1.46000,nan) (1.46500,nan) (1.47000,nan) (1.47500,nan) (1.48000,nan) (1.48500,nan) (1.49000,nan) (1.49500,nan) (1.50000,nan) (1.50500,nan) (1.51000,nan) (1.51500,nan) (1.52000,nan) (1.52500,nan) (1.53000,nan) (1.53500,nan) (1.54000,nan) (1.54500,nan) (1.55000,nan) (1.55500,nan) (1.56000,nan) (1.56500,nan) (1.57000,nan) (1.57500,nan) (1.58000,nan) (1.58500,nan) (1.59000,nan) (1.59500,nan) (1.60000,nan) (1.60500,nan) (1.61000,nan) (1.61500,nan) (1.62000,nan) (1.62500,nan) (1.63000,nan) (1.63500,nan) (1.64000,nan) (1.64500,nan) (1.65000,nan) (1.65500,nan) (1.66000,nan) (1.66500,nan) (1.67000,nan) (1.67500,nan) (1.68000,nan) (1.68500,nan) (1.69000,nan) (1.69500,nan) (1.70000,nan) (1.70500,nan) (1.71000,nan) (1.71500,nan) (1.72000,nan) (1.72500,nan) (1.73000,nan) (1.73500,nan) (1.74000,nan) (1.74500,nan) (1.75000,nan) (1.75500,nan) (1.76000,nan) (1.76500,nan) (1.77000,nan) (1.77500,nan) (1.78000,nan) (1.78500,nan) (1.79000,nan) (1.79500,nan) (1.80000,nan) (1.80500,nan) (1.81000,nan) (1.81500,nan) (1.82000,nan) (1.82500,nan) (1.83000,nan) (1.83500,nan) (1.84000,nan) (1.84500,nan) (1.85000,nan) (1.85500,nan) (1.86000,nan) (1.86500,nan) (1.87000,nan) (1.87500,nan) (1.88000,nan) (1.88500,nan) (1.89000,nan) (1.89500,nan) (1.90000,nan) (1.90500,nan) (1.91000,nan) (1.91500,nan) (1.92000,nan) (1.92500,nan) (1.93000,nan) (1.93500,nan) (1.94000,nan) (1.94500,nan) (1.95000,nan) (1.95500,nan) (1.96000,nan) (1.96500,nan) (1.97000,nan) (1.97500,nan) (1.98000,nan) (1.98500,nan) (1.99000,nan) (1.99500,nan) (2.00000,nan) (2.00500,nan) (2.01000,nan) (2.01500,nan) (2.02000,nan) (2.02500,nan) (2.03000,nan) (2.03500,nan) (2.04000,nan) (2.04500,nan) (2.05000,nan) (2.05500,nan) (2.06000,nan) (2.06500,nan) (2.07000,nan) (2.07500,nan) (2.08000,nan) (2.08500,nan) (2.09000,nan) (2.09500,nan) (2.10000,nan) (2.10500,nan) (2.11000,nan) (2.11500,nan) (2.12000,nan) (2.12500,nan) (2.13000,nan) (2.13500,nan) (2.14000,nan) (2.14500,nan) (2.15000,nan) (2.15500,nan) (2.16000,nan) (2.16500,nan) (2.17000,nan) (2.17500,nan) (2.18000,nan) (2.18500,nan) (2.19000,nan) (2.19500,nan) (2.20000,nan) (2.20500,nan) (2.21000,nan) (2.21500,nan) (2.22000,nan) (2.22500,nan) (2.23000,nan) (2.23500,nan) (2.24000,nan) (2.24500,nan) (2.25000,nan) (2.25500,nan) (2.26000,nan) (2.26500,nan) (2.27000,nan) (2.27500,nan) (2.28000,nan) (2.28500,nan) (2.29000,nan) (2.29500,nan) (2.30000,nan) (2.30500,nan) (2.31000,nan) (2.31500,nan) (2.32000,nan) (2.32500,nan) (2.33000,nan) (2.33500,nan) (2.34000,nan) (2.34500,nan) (2.35000,nan) (2.35500,nan) (2.36000,nan) (2.36500,nan) (2.37000,nan) (2.37500,nan) (2.38000,nan) (2.38500,nan) (2.39000,nan) (2.39500,nan) (2.40000,nan) (2.40500,nan) (2.41000,nan) (2.41500,nan) (2.42000,nan) (2.42500,nan) (2.43000,nan) (2.43500,nan) (2.44000,nan) (2.44500,nan) (2.45000,nan) (2.45500,nan) (2.46000,nan) (2.46500,nan) (2.47000,nan) (2.47500,nan) (2.48000,nan) (2.48500,nan) (2.49000,nan) (2.49500,nan) (2.50000,nan) (2.50500,nan) (2.51000,nan) (2.51500,nan) (2.52000,nan) (2.52500,nan) (2.53000,nan) (2.53500,nan) (2.54000,nan) (2.54500,nan) (2.55000,nan) (2.55500,nan) (2.56000,nan) (2.56500,nan) (2.57000,nan) (2.57500,nan) (2.58000,nan) (2.58500,nan) (2.59000,nan) (2.59500,nan) (2.60000,nan) (2.60500,nan) (2.61000,nan) (2.61500,nan) (2.62000,nan) (2.62500,nan) (2.63000,nan) (2.63500,nan) (2.64000,nan) (2.64500,nan) (2.65000,nan) (2.65500,nan) (2.66000,nan) (2.66500,nan) (2.67000,nan) (2.67500,nan) (2.68000,nan) (2.68500,nan) (2.69000,nan) (2.69500,nan) (2.70000,nan) (2.70500,nan) (2.71000,nan) (2.71500,nan) (2.72000,nan) (2.72500,nan) (2.73000,nan) (2.73500,nan) (2.74000,nan) (2.74500,nan) (2.75000,nan) (2.75500,nan) (2.76000,nan) (2.76500,nan) (2.77000,nan) (2.77500,nan) (2.78000,nan) (2.78500,nan) (2.79000,nan) (2.79500,nan) (2.80000,nan) (2.80500,nan) (2.81000,nan) (2.81500,nan) (2.82000,nan) (2.82500,nan) (2.83000,nan) (2.83500,nan) (2.84000,nan) (2.84500,nan) (2.85000,nan) (2.85500,nan) (2.86000,nan) (2.86500,nan) (2.87000,nan) (2.87500,nan) (2.88000,nan) (2.88500,nan) (2.89000,nan) (2.89500,nan) (2.90000,nan) (2.90500,nan) (2.91000,nan) (2.91500,nan) (2.92000,nan) (2.92500,nan) (2.93000,nan) (2.93500,nan) (2.94000,nan) (2.94500,nan) (2.95000,nan) (2.95500,nan) (2.96000,nan) (2.96500,nan) (2.97000,nan) (2.97500,nan) (2.98000,nan) (2.98500,nan) (2.99000,nan) (2.99500,nan) (3.00000,nan) (3.00500,nan) (3.01000,nan) (3.01500,nan) (3.02000,nan) (3.02500,nan) (3.03000,nan) (3.03500,nan) (3.04000,nan) (3.04500,nan) (3.05000,nan) (3.05500,nan) (3.06000,nan) (3.06500,nan) (3.07000,nan) (3.07500,nan) (3.08000,nan) (3.08500,nan) (3.09000,nan) (3.09500,nan) (3.10000,nan) (3.10500,nan) (3.11000,nan) (3.11500,nan) (3.12000,nan) (3.12500,nan) (3.13000,nan) (3.13500,nan) (3.14000,nan) (3.14500,nan) (3.15000,nan) (3.15500,nan) (3.16000,nan) (3.16500,nan) (3.17000,nan) (3.17500,nan) (3.18000,nan) (3.18500,nan) (3.19000,nan) (3.19500,nan) (3.20000,nan) (3.20500,nan) (3.21000,nan) (3.21500,nan) (3.22000,nan) (3.22500,nan) (3.23000,nan) (3.23500,nan) (3.24000,nan) (3.24500,nan) (3.25000,nan) (3.25500,nan) (3.26000,nan) (3.26500,nan) (3.27000,nan) (3.27500,nan) (3.28000,nan) (3.28500,nan) (3.29000,nan) (3.29500,nan) (3.30000,nan) (3.30500,nan) (3.31000,nan) (3.31500,nan) (3.32000,nan) (3.32500,nan) (3.33000,nan) (3.33500,nan) (3.34000,nan) (3.34500,nan) (3.35000,nan) (3.35500,nan) (3.36000,nan) (3.36500,nan) (3.37000,nan) (3.37500,nan) (3.38000,nan) (3.38500,nan) (3.39000,nan) (3.39500,nan) (3.40000,nan) (3.40500,nan) (3.41000,nan) (3.41500,nan) (3.42000,nan) (3.42500,nan) (3.43000,nan) (3.43500,nan) (3.44000,nan) (3.44500,nan) (3.45000,nan) (3.45500,nan) (3.46000,nan) (3.46500,nan) (3.47000,nan) (3.47500,nan) (3.48000,nan) (3.48500,nan) (3.49000,nan) (3.49500,nan) (3.50000,nan) (3.50500,nan) (3.51000,nan) (3.51500,nan) (3.52000,nan) (3.52500,nan) (3.53000,nan) (3.53500,nan) (3.54000,nan) (3.54500,nan) (3.55000,nan) (3.55500,nan) (3.56000,nan) (3.56500,nan) (3.57000,nan) (3.57500,nan) (3.58000,nan) (3.58500,nan) (3.59000,nan) (3.59500,nan) (3.60000,nan) (3.60500,nan) (3.61000,nan) (3.61500,nan) (3.62000,nan) (3.62500,nan) (3.63000,nan) (3.63500,nan) (3.64000,nan) (3.64500,nan) (3.65000,nan) (3.65500,nan) (3.66000,nan) (3.66500,nan) (3.67000,nan) (3.67500,nan) (3.68000,-1.5127938) (3.68500,-1.2776240) (3.69000,-1.1288343) (3.69500,-1.0204467) (3.70000,-0.9355710) (3.70500,-0.8660948) (3.71000,-0.8074906) (3.71500,-0.7569719) (3.72000,-0.7127014) (3.72500,-0.6734035) (3.73000,-0.6381561) (3.73500,-0.6062715) (3.74000,-0.5772224) (3.74500,-0.5505964) (3.75000,-0.5260641) (3.75500,-0.5033588) (3.76000,-0.4822612) (3.76500,-0.4625889) (3.77000,-0.4441888) (3.77500,-0.4269305) (3.78000,-0.4107028) (3.78500,-0.3954096) (3.79000,-0.3809673) (3.79500,-0.3673032) (3.80000,-0.3543530) (3.80500,-0.3420600) (3.81000,-0.3303739) (3.81500,-0.3192497) (3.82000,-0.3086474) (3.82500,-0.2985305) (3.83000,-0.2888666) (3.83500,-0.2796259) (3.84000,-0.2707816) (3.84500,-0.2623093) (3.85000,-0.2541864) (3.85500,-0.2463926) (3.86000,-0.2389090) (3.86500,-0.2317183) (3.87000,-0.2248044) (3.87500,-0.2181526) (3.88000,-0.2117491) (3.88500,-0.2055814) (3.89000,-0.1996375) (3.89500,-0.1939064) (3.90000,-0.1883780) (3.90500,-0.1830425) (3.91000,-0.1778910) (3.91500,-0.1729152) (3.92000,-0.1681071) (3.92500,-0.1634593) (3.93000,-0.1589648) (3.93500,-0.1546172) (3.94000,-0.1504101) (3.94500,-0.1463379) (3.95000,-0.1423948) (3.95500,-0.1385758) (3.96000,-0.1348759) (3.96500,-0.1312904) (3.97000,-0.1278149) (3.97500,-0.1244452) (3.98000,-0.1211773) (3.98500,-0.1180073) (3.99000,-0.1149316) (3.99500,-0.1119469) (4.00000,-0.1090498)};\addlegendentry{O}
\addplot[thick,dotted,blue] coordinates {(-2.00000,nan) (-1.99500,nan) (-1.99000,nan) (-1.98500,nan) (-1.98000,nan) (-1.97500,nan) (-1.97000,nan) (-1.96500,nan) (-1.96000,nan) (-1.95500,nan) (-1.95000,nan) (-1.94500,nan) (-1.94000,nan) (-1.93500,nan) (-1.93000,nan) (-1.92500,nan) (-1.92000,nan) (-1.91500,nan) (-1.91000,nan) (-1.90500,nan) (-1.90000,nan) (-1.89500,nan) (-1.89000,nan) (-1.88500,nan) (-1.88000,nan) (-1.87500,nan) (-1.87000,nan) (-1.86500,nan) (-1.86000,nan) (-1.85500,nan) (-1.85000,nan) (-1.84500,nan) (-1.84000,nan) (-1.83500,nan) (-1.83000,nan) (-1.82500,nan) (-1.82000,nan) (-1.81500,nan) (-1.81000,nan) (-1.80500,nan) (-1.80000,nan) (-1.79500,nan) (-1.79000,nan) (-1.78500,nan) (-1.78000,nan) (-1.77500,nan) (-1.77000,nan) (-1.76500,nan) (-1.76000,nan) (-1.75500,nan) (-1.75000,nan) (-1.74500,nan) (-1.74000,nan) (-1.73500,nan) (-1.73000,nan) (-1.72500,nan) (-1.72000,nan) (-1.71500,nan) (-1.71000,nan) (-1.70500,nan) (-1.70000,nan) (-1.69500,nan) (-1.69000,nan) (-1.68500,nan) (-1.68000,nan) (-1.67500,nan) (-1.67000,nan) (-1.66500,nan) (-1.66000,nan) (-1.65500,nan) (-1.65000,nan) (-1.64500,nan) (-1.64000,nan) (-1.63500,nan) (-1.63000,nan) (-1.62500,nan) (-1.62000,nan) (-1.61500,nan) (-1.61000,nan) (-1.60500,nan) (-1.60000,nan) (-1.59500,nan) (-1.59000,nan) (-1.58500,nan) (-1.58000,nan) (-1.57500,nan) (-1.57000,nan) (-1.56500,nan) (-1.56000,nan) (-1.55500,nan) (-1.55000,nan) (-1.54500,nan) (-1.54000,nan) (-1.53500,nan) (-1.53000,nan) (-1.52500,nan) (-1.52000,nan) (-1.51500,nan) (-1.51000,nan) (-1.50500,nan) (-1.50000,nan) (-1.49500,nan) (-1.49000,nan) (-1.48500,nan) (-1.48000,nan) (-1.47500,nan) (-1.47000,nan) (-1.46500,nan) (-1.46000,nan) (-1.45500,nan) (-1.45000,nan) (-1.44500,nan) (-1.44000,nan) (-1.43500,nan) (-1.43000,nan) (-1.42500,nan) (-1.42000,nan) (-1.41500,nan) (-1.41000,nan) (-1.40500,nan) (-1.40000,nan) (-1.39500,nan) (-1.39000,nan) (-1.38500,nan) (-1.38000,nan) (-1.37500,nan) (-1.37000,nan) (-1.36500,nan) (-1.36000,nan) (-1.35500,nan) (-1.35000,nan) (-1.34500,nan) (-1.34000,nan) (-1.33500,nan) (-1.33000,nan) (-1.32500,nan) (-1.32000,nan) (-1.31500,nan) (-1.31000,nan) (-1.30500,nan) (-1.30000,nan) (-1.29500,nan) (-1.29000,nan) (-1.28500,nan) (-1.28000,nan) (-1.27500,nan) (-1.27000,nan) (-1.26500,nan) (-1.26000,nan) (-1.25500,nan) (-1.25000,nan) (-1.24500,nan) (-1.24000,nan) (-1.23500,nan) (-1.23000,nan) (-1.22500,nan) (-1.22000,nan) (-1.21500,nan) (-1.21000,nan) (-1.20500,nan) (-1.20000,nan) (-1.19500,nan) (-1.19000,nan) (-1.18500,nan) (-1.18000,nan) (-1.17500,nan) (-1.17000,nan) (-1.16500,nan) (-1.16000,nan) (-1.15500,nan) (-1.15000,nan) (-1.14500,nan) (-1.14000,nan) (-1.13500,nan) (-1.13000,nan) (-1.12500,nan) (-1.12000,nan) (-1.11500,nan) (-1.11000,nan) (-1.10500,nan) (-1.10000,nan) (-1.09500,nan) (-1.09000,nan) (-1.08500,nan) (-1.08000,nan) (-1.07500,nan) (-1.07000,nan) (-1.06500,nan) (-1.06000,nan) (-1.05500,nan) (-1.05000,nan) (-1.04500,nan) (-1.04000,nan) (-1.03500,nan) (-1.03000,nan) (-1.02500,nan) (-1.02000,nan) (-1.01500,nan) (-1.01000,nan) (-1.00500,nan) (-1.00000,nan) (-0.99500,nan) (-0.99000,nan) (-0.98500,nan) (-0.98000,nan) (-0.97500,nan) (-0.97000,nan) (-0.96500,nan) (-0.96000,nan) (-0.95500,nan) (-0.95000,nan) (-0.94500,nan) (-0.94000,nan) (-0.93500,nan) (-0.93000,nan) (-0.92500,nan) (-0.92000,nan) (-0.91500,nan) (-0.91000,nan) (-0.90500,nan) (-0.90000,nan) (-0.89500,nan) (-0.89000,nan) (-0.88500,nan) (-0.88000,nan) (-0.87500,nan) (-0.87000,nan) (-0.86500,nan) (-0.86000,nan) (-0.85500,nan) (-0.85000,nan) (-0.84500,nan) (-0.84000,nan) (-0.83500,nan) (-0.83000,nan) (-0.82500,nan) (-0.82000,nan) (-0.81500,nan) (-0.81000,nan) (-0.80500,nan) (-0.80000,nan) (-0.79500,nan) (-0.79000,nan) (-0.78500,nan) (-0.78000,nan) (-0.77500,nan) (-0.77000,nan) (-0.76500,nan) (-0.76000,nan) (-0.75500,nan) (-0.75000,nan) (-0.74500,nan) (-0.74000,nan) (-0.73500,nan) (-0.73000,nan) (-0.72500,nan) (-0.72000,nan) (-0.71500,nan) (-0.71000,nan) (-0.70500,nan) (-0.70000,nan) (-0.69500,nan) (-0.69000,nan) (-0.68500,nan) (-0.68000,nan) (-0.67500,nan) (-0.67000,nan) (-0.66500,nan) (-0.66000,nan) (-0.65500,nan) (-0.65000,nan) (-0.64500,nan) (-0.64000,nan) (-0.63500,nan) (-0.63000,nan) (-0.62500,nan) (-0.62000,nan) (-0.61500,nan) (-0.61000,nan) (-0.60500,nan) (-0.60000,nan) (-0.59500,nan) (-0.59000,nan) (-0.58500,nan) (-0.58000,nan) (-0.57500,nan) (-0.57000,nan) (-0.56500,nan) (-0.56000,nan) (-0.55500,nan) (-0.55000,nan) (-0.54500,nan) (-0.54000,nan) (-0.53500,nan) (-0.53000,nan) (-0.52500,nan) (-0.52000,nan) (-0.51500,nan) (-0.51000,nan) (-0.50500,nan) (-0.50000,nan) (-0.49500,nan) (-0.49000,nan) (-0.48500,nan) (-0.48000,nan) (-0.47500,nan) (-0.47000,nan) (-0.46500,nan) (-0.46000,nan) (-0.45500,nan) (-0.45000,nan) (-0.44500,nan) (-0.44000,nan) (-0.43500,nan) (-0.43000,nan) (-0.42500,nan) (-0.42000,nan) (-0.41500,nan) (-0.41000,nan) (-0.40500,nan) (-0.40000,nan) (-0.39500,nan) (-0.39000,nan) (-0.38500,nan) (-0.38000,nan) (-0.37500,nan) (-0.37000,nan) (-0.36500,nan) (-0.36000,nan) (-0.35500,nan) (-0.35000,nan) (-0.34500,nan) (-0.34000,nan) (-0.33500,nan) (-0.33000,nan) (-0.32500,nan) (-0.32000,nan) (-0.31500,nan) (-0.31000,nan) (-0.30500,nan) (-0.30000,nan) (-0.29500,nan) (-0.29000,nan) (-0.28500,nan) (-0.28000,nan) (-0.27500,nan) (-0.27000,nan) (-0.26500,nan) (-0.26000,nan) (-0.25500,nan) (-0.25000,nan) (-0.24500,nan) (-0.24000,nan) (-0.23500,nan) (-0.23000,nan) (-0.22500,nan) (-0.22000,nan) (-0.21500,nan) (-0.21000,nan) (-0.20500,nan) (-0.20000,nan) (-0.19500,nan) (-0.19000,nan) (-0.18500,nan) (-0.18000,nan) (-0.17500,nan) (-0.17000,nan) (-0.16500,nan) (-0.16000,nan) (-0.15500,nan) (-0.15000,nan) (-0.14500,nan) (-0.14000,nan) (-0.13500,nan) (-0.13000,nan) (-0.12500,nan) (-0.12000,nan) (-0.11500,nan) (-0.11000,nan) (-0.10500,nan) (-0.10000,nan) (-0.09500,nan) (-0.09000,nan) (-0.08500,nan) (-0.08000,nan) (-0.07500,nan) (-0.07000,nan) (-0.06500,nan) (-0.06000,nan) (-0.05500,nan) (-0.05000,nan) (-0.04500,nan) (-0.04000,nan) (-0.03500,nan) (-0.03000,nan) (-0.02500,nan) (-0.02000,nan) (-0.01500,nan) (-0.01000,nan) (-0.00500,nan) (0.00000,nan) (0.00500,nan) (0.01000,nan) (0.01500,nan) (0.02000,nan) (0.02500,nan) (0.03000,nan) (0.03500,nan) (0.04000,nan) (0.04500,nan) (0.05000,nan) (0.05500,nan) (0.06000,nan) (0.06500,nan) (0.07000,nan) (0.07500,nan) (0.08000,nan) (0.08500,nan) (0.09000,nan) (0.09500,nan) (0.10000,nan) (0.10500,nan) (0.11000,nan) (0.11500,nan) (0.12000,nan) (0.12500,nan) (0.13000,nan) (0.13500,nan) (0.14000,nan) (0.14500,nan) (0.15000,nan) (0.15500,nan) (0.16000,nan) (0.16500,nan) (0.17000,nan) (0.17500,nan) (0.18000,nan) (0.18500,nan) (0.19000,nan) (0.19500,nan) (0.20000,nan) (0.20500,nan) (0.21000,nan) (0.21500,nan) (0.22000,nan) (0.22500,nan) (0.23000,nan) (0.23500,nan) (0.24000,nan) (0.24500,nan) (0.25000,nan) (0.25500,nan) (0.26000,nan) (0.26500,nan) (0.27000,nan) (0.27500,nan) (0.28000,nan) (0.28500,nan) (0.29000,nan) (0.29500,nan) (0.30000,nan) (0.30500,nan) (0.31000,nan) (0.31500,nan) (0.32000,nan) (0.32500,nan) (0.33000,nan) (0.33500,nan) (0.34000,nan) (0.34500,nan) (0.35000,nan) (0.35500,nan) (0.36000,nan) (0.36500,nan) (0.37000,nan) (0.37500,nan) (0.38000,nan) (0.38500,nan) (0.39000,nan) (0.39500,nan) (0.40000,nan) (0.40500,nan) (0.41000,nan) (0.41500,nan) (0.42000,nan) (0.42500,nan) (0.43000,nan) (0.43500,nan) (0.44000,nan) (0.44500,nan) (0.45000,nan) (0.45500,nan) (0.46000,nan) (0.46500,nan) (0.47000,nan) (0.47500,nan) (0.48000,nan) (0.48500,nan) (0.49000,nan) (0.49500,nan) (0.50000,nan) (0.50500,nan) (0.51000,nan) (0.51500,nan) (0.52000,nan) (0.52500,nan) (0.53000,nan) (0.53500,nan) (0.54000,nan) (0.54500,nan) (0.55000,nan) (0.55500,nan) (0.56000,nan) (0.56500,nan) (0.57000,nan) (0.57500,nan) (0.58000,nan) (0.58500,nan) (0.59000,nan) (0.59500,nan) (0.60000,nan) (0.60500,nan) (0.61000,nan) (0.61500,nan) (0.62000,nan) (0.62500,nan) (0.63000,nan) (0.63500,nan) (0.64000,nan) (0.64500,nan) (0.65000,nan) (0.65500,nan) (0.66000,nan) (0.66500,nan) (0.67000,nan) (0.67500,nan) (0.68000,nan) (0.68500,nan) (0.69000,nan) (0.69500,nan) (0.70000,nan) (0.70500,nan) (0.71000,nan) (0.71500,nan) (0.72000,nan) (0.72500,nan) (0.73000,nan) (0.73500,nan) (0.74000,nan) (0.74500,nan) (0.75000,nan) (0.75500,nan) (0.76000,nan) (0.76500,nan) (0.77000,nan) (0.77500,nan) (0.78000,nan) (0.78500,nan) (0.79000,nan) (0.79500,nan) (0.80000,nan) (0.80500,nan) (0.81000,nan) (0.81500,nan) (0.82000,nan) (0.82500,nan) (0.83000,nan) (0.83500,nan) (0.84000,nan) (0.84500,nan) (0.85000,nan) (0.85500,nan) (0.86000,nan) (0.86500,nan) (0.87000,nan) (0.87500,nan) (0.88000,nan) (0.88500,nan) (0.89000,nan) (0.89500,nan) (0.90000,nan) (0.90500,nan) (0.91000,nan) (0.91500,nan) (0.92000,nan) (0.92500,nan) (0.93000,nan) (0.93500,nan) (0.94000,nan) (0.94500,nan) (0.95000,nan) (0.95500,nan) (0.96000,nan) (0.96500,nan) (0.97000,nan) (0.97500,nan) (0.98000,nan) (0.98500,nan) (0.99000,nan) (0.99500,nan) (1.00000,nan) (1.00500,nan) (1.01000,nan) (1.01500,nan) (1.02000,nan) (1.02500,nan) (1.03000,nan) (1.03500,nan) (1.04000,nan) (1.04500,nan) (1.05000,nan) (1.05500,nan) (1.06000,nan) (1.06500,nan) (1.07000,nan) (1.07500,nan) (1.08000,nan) (1.08500,nan) (1.09000,nan) (1.09500,nan) (1.10000,nan) (1.10500,nan) (1.11000,nan) (1.11500,nan) (1.12000,nan) (1.12500,nan) (1.13000,nan) (1.13500,nan) (1.14000,nan) (1.14500,nan) (1.15000,nan) (1.15500,nan) (1.16000,nan) (1.16500,nan) (1.17000,nan) (1.17500,nan) (1.18000,nan) (1.18500,nan) (1.19000,nan) (1.19500,nan) (1.20000,nan) (1.20500,nan) (1.21000,nan) (1.21500,nan) (1.22000,nan) (1.22500,nan) (1.23000,nan) (1.23500,nan) (1.24000,nan) (1.24500,nan) (1.25000,nan) (1.25500,nan) (1.26000,nan) (1.26500,nan) (1.27000,nan) (1.27500,nan) (1.28000,nan) (1.28500,nan) (1.29000,nan) (1.29500,nan) (1.30000,nan) (1.30500,nan) (1.31000,nan) (1.31500,nan) (1.32000,nan) (1.32500,nan) (1.33000,nan) (1.33500,nan) (1.34000,nan) (1.34500,nan) (1.35000,nan) (1.35500,nan) (1.36000,nan) (1.36500,nan) (1.37000,nan) (1.37500,nan) (1.38000,nan) (1.38500,nan) (1.39000,nan) (1.39500,nan) (1.40000,nan) (1.40500,nan) (1.41000,nan) (1.41500,nan) (1.42000,nan) (1.42500,nan) (1.43000,nan) (1.43500,nan) (1.44000,nan) (1.44500,nan) (1.45000,nan) (1.45500,nan) (1.46000,nan) (1.46500,nan) (1.47000,nan) (1.47500,nan) (1.48000,nan) (1.48500,nan) (1.49000,nan) (1.49500,nan) (1.50000,nan) (1.50500,nan) (1.51000,nan) (1.51500,nan) (1.52000,nan) (1.52500,nan) (1.53000,nan) (1.53500,nan) (1.54000,nan) (1.54500,nan) (1.55000,nan) (1.55500,nan) (1.56000,nan) (1.56500,nan) (1.57000,nan) (1.57500,nan) (1.58000,nan) (1.58500,nan) (1.59000,nan) (1.59500,nan) (1.60000,nan) (1.60500,nan) (1.61000,nan) (1.61500,nan) (1.62000,nan) (1.62500,nan) (1.63000,nan) (1.63500,nan) (1.64000,nan) (1.64500,nan) (1.65000,nan) (1.65500,nan) (1.66000,nan) (1.66500,nan) (1.67000,nan) (1.67500,nan) (1.68000,nan) (1.68500,nan) (1.69000,nan) (1.69500,nan) (1.70000,nan) (1.70500,nan) (1.71000,nan) (1.71500,nan) (1.72000,nan) (1.72500,nan) (1.73000,nan) (1.73500,nan) (1.74000,nan) (1.74500,nan) (1.75000,nan) (1.75500,nan) (1.76000,nan) (1.76500,nan) (1.77000,nan) (1.77500,nan) (1.78000,nan) (1.78500,nan) (1.79000,nan) (1.79500,nan) (1.80000,nan) (1.80500,nan) (1.81000,nan) (1.81500,nan) (1.82000,nan) (1.82500,nan) (1.83000,nan) (1.83500,nan) (1.84000,nan) (1.84500,nan) (1.85000,nan) (1.85500,nan) (1.86000,nan) (1.86500,nan) (1.87000,nan) (1.87500,nan) (1.88000,nan) (1.88500,nan) (1.89000,nan) (1.89500,nan) (1.90000,nan) (1.90500,nan) (1.91000,nan) (1.91500,nan) (1.92000,nan) (1.92500,nan) (1.93000,nan) (1.93500,nan) (1.94000,nan) (1.94500,nan) (1.95000,nan) (1.95500,nan) (1.96000,nan) (1.96500,nan) (1.97000,nan) (1.97500,nan) (1.98000,nan) (1.98500,nan) (1.99000,nan) (1.99500,nan) (2.00000,nan) (2.00500,nan) (2.01000,nan) (2.01500,nan) (2.02000,nan) (2.02500,nan) (2.03000,nan) (2.03500,nan) (2.04000,nan) (2.04500,nan) (2.05000,nan) (2.05500,nan) (2.06000,nan) (2.06500,nan) (2.07000,nan) (2.07500,nan) (2.08000,nan) (2.08500,nan) (2.09000,nan) (2.09500,nan) (2.10000,nan) (2.10500,nan) (2.11000,nan) (2.11500,nan) (2.12000,nan) (2.12500,nan) (2.13000,nan) (2.13500,nan) (2.14000,nan) (2.14500,nan) (2.15000,nan) (2.15500,nan) (2.16000,nan) (2.16500,nan) (2.17000,nan) (2.17500,nan) (2.18000,nan) (2.18500,nan) (2.19000,nan) (2.19500,nan) (2.20000,nan) (2.20500,nan) (2.21000,nan) (2.21500,nan) (2.22000,nan) (2.22500,nan) (2.23000,nan) (2.23500,nan) (2.24000,nan) (2.24500,nan) (2.25000,nan) (2.25500,nan) (2.26000,nan) (2.26500,nan) (2.27000,nan) (2.27500,nan) (2.28000,nan) (2.28500,nan) (2.29000,nan) (2.29500,nan) (2.30000,nan) (2.30500,nan) (2.31000,nan) (2.31500,nan) (2.32000,nan) (2.32500,nan) (2.33000,nan) (2.33500,nan) (2.34000,nan) (2.34500,nan) (2.35000,nan) (2.35500,nan) (2.36000,nan) (2.36500,nan) (2.37000,nan) (2.37500,nan) (2.38000,nan) (2.38500,nan) (2.39000,nan) (2.39500,nan) (2.40000,nan) (2.40500,nan) (2.41000,nan) (2.41500,nan) (2.42000,nan) (2.42500,nan) (2.43000,nan) (2.43500,nan) (2.44000,nan) (2.44500,nan) (2.45000,nan) (2.45500,nan) (2.46000,nan) (2.46500,nan) (2.47000,nan) (2.47500,nan) (2.48000,nan) (2.48500,nan) (2.49000,nan) (2.49500,nan) (2.50000,nan) (2.50500,nan) (2.51000,nan) (2.51500,nan) (2.52000,nan) (2.52500,nan) (2.53000,nan) (2.53500,nan) (2.54000,nan) (2.54500,nan) (2.55000,nan) (2.55500,nan) (2.56000,nan) (2.56500,nan) (2.57000,nan) (2.57500,nan) (2.58000,nan) (2.58500,nan) (2.59000,nan) (2.59500,nan) (2.60000,nan) (2.60500,nan) (2.61000,nan) (2.61500,nan) (2.62000,nan) (2.62500,nan) (2.63000,nan) (2.63500,nan) (2.64000,nan) (2.64500,nan) (2.65000,nan) (2.65500,nan) (2.66000,nan) (2.66500,nan) (2.67000,nan) (2.67500,nan) (2.68000,nan) (2.68500,nan) (2.69000,nan) (2.69500,nan) (2.70000,nan) (2.70500,nan) (2.71000,nan) (2.71500,nan) (2.72000,nan) (2.72500,nan) (2.73000,nan) (2.73500,nan) (2.74000,nan) (2.74500,nan) (2.75000,nan) (2.75500,nan) (2.76000,nan) (2.76500,nan) (2.77000,nan) (2.77500,nan) (2.78000,nan) (2.78500,nan) (2.79000,nan) (2.79500,nan) (2.80000,nan) (2.80500,nan) (2.81000,nan) (2.81500,nan) (2.82000,nan) (2.82500,nan) (2.83000,nan) (2.83500,nan) (2.84000,nan) (2.84500,nan) (2.85000,nan) (2.85500,nan) (2.86000,nan) (2.86500,nan) (2.87000,nan) (2.87500,nan) (2.88000,nan) (2.88500,nan) (2.89000,nan) (2.89500,nan) (2.90000,nan) (2.90500,nan) (2.91000,nan) (2.91500,nan) (2.92000,nan) (2.92500,nan) (2.93000,nan) (2.93500,nan) (2.94000,nan) (2.94500,nan) (2.95000,nan) (2.95500,nan) (2.96000,nan) (2.96500,nan) (2.97000,nan) (2.97500,nan) (2.98000,nan) (2.98500,nan) (2.99000,nan) (2.99500,nan) (3.00000,nan) (3.00500,nan) (3.01000,nan) (3.01500,nan) (3.02000,nan) (3.02500,nan) (3.03000,nan) (3.03500,nan) (3.04000,nan) (3.04500,nan) (3.05000,nan) (3.05500,nan) (3.06000,nan) (3.06500,nan) (3.07000,nan) (3.07500,nan) (3.08000,nan) (3.08500,nan) (3.09000,nan) (3.09500,nan) (3.10000,nan) (3.10500,nan) (3.11000,nan) (3.11500,nan) (3.12000,nan) (3.12500,nan) (3.13000,nan) (3.13500,nan) (3.14000,nan) (3.14500,nan) (3.15000,nan) (3.15500,nan) (3.16000,nan) (3.16500,nan) (3.17000,nan) (3.17500,nan) (3.18000,nan) (3.18500,nan) (3.19000,nan) (3.19500,nan) (3.20000,nan) (3.20500,nan) (3.21000,nan) (3.21500,nan) (3.22000,nan) (3.22500,nan) (3.23000,nan) (3.23500,nan) (3.24000,nan) (3.24500,nan) (3.25000,nan) (3.25500,nan) (3.26000,nan) (3.26500,nan) (3.27000,nan) (3.27500,nan) (3.28000,nan) (3.28500,nan) (3.29000,nan) (3.29500,nan) (3.30000,nan) (3.30500,nan) (3.31000,nan) (3.31500,nan) (3.32000,nan) (3.32500,nan) (3.33000,nan) (3.33500,nan) (3.34000,nan) (3.34500,nan) (3.35000,nan) (3.35500,nan) (3.36000,nan) (3.36500,nan) (3.37000,nan) (3.37500,nan) (3.38000,nan) (3.38500,nan) (3.39000,nan) (3.39500,nan) (3.40000,nan) (3.40500,nan) (3.41000,nan) (3.41500,nan) (3.42000,nan) (3.42500,nan) (3.43000,nan) (3.43500,nan) (3.44000,nan) (3.44500,nan) (3.45000,nan) (3.45500,nan) (3.46000,nan) (3.46500,nan) (3.47000,nan) (3.47500,nan) (3.48000,nan) (3.48500,nan) (3.49000,nan) (3.49500,nan) (3.50000,nan) (3.50500,nan) (3.51000,-1.8946940) (3.51500,-1.3368392) (3.52000,-1.1034788) (3.52500,-0.9545310) (3.53000,-0.8453268) (3.53500,-0.7593066) (3.54000,-0.6884787) (3.54500,-0.6283680) (3.55000,-0.5762135) (3.55500,-0.5301901) (3.56000,-0.4890276) (3.56500,-0.4518053) (3.57000,-0.4178338) (3.57500,-0.3865829) (3.58000,-0.3576350) (3.58500,-0.3306551) (3.59000,-0.3053690) (3.59500,-0.2815493) (3.60000,-0.2590042) (3.60500,-0.2375697) (3.61000,-0.2171037) (3.61500,-0.1974814) (3.62000,-0.1785917) (3.62500,-0.1603344) (3.63000,-0.1426176) (3.63500,-0.1253560) (3.64000,-0.1084687) (3.64500,-0.0918783) (3.65000,-0.0755090) (3.65500,-0.0592853) (3.66000,-0.0431305) (3.66500,-0.0269655) (3.67000,-0.0107066) (3.67500,0.0057357) (3.68000,0.0224596) (3.68500,0.0395746) (3.69000,0.0572047) (3.69500,0.0754931) (3.70000,0.0946076) (3.70500,0.1147491) (3.71000,0.1361612) (3.71500,0.1591458) (3.72000,0.1840838) (3.72500,0.2114661) (3.73000,0.2419416) (3.73500,0.2763921) (3.74000,0.3160564) (3.74500,0.3627491) (3.75000,0.4192713) (3.75500,0.4902646) (3.76000,0.5842313) (3.76500,0.7193814) (3.77000,0.9476332) (3.77500,1.6131841) (3.78000,nan) (3.78500,nan) (3.79000,nan) (3.79500,nan) (3.80000,nan) (3.80500,nan) (3.81000,nan) (3.81500,nan) (3.82000,nan) (3.82500,nan) (3.83000,nan) (3.83500,nan) (3.84000,-1.6511296) (3.84500,-1.0619753) (3.85000,-0.8445910) (3.85500,-0.7148789) (3.86000,-0.6248914) (3.86500,-0.5573854) (3.87000,-0.5042123) (3.87500,-0.4608929) (3.88000,-0.4247127) (3.88500,-0.3939095) (3.89000,-0.3672797) (3.89500,-0.3439677) (3.90000,-0.3233454) (3.90500,-0.3049400) (3.91000,-0.2883874) (3.91500,-0.2734022) (3.92000,-0.2597569) (3.92500,-0.2472678) (3.93000,-0.2357844) (3.93500,-0.2251825) (3.94000,-0.2153584) (3.94500,-0.2062245) (3.95000,-0.1977066) (3.95500,-0.1897413) (3.96000,-0.1822739) (3.96500,-0.1752570) (3.97000,-0.1686493) (3.97500,-0.1624149) (3.98000,-0.1565218) (3.98500,-0.1509420) (3.99000,-0.1456506) (3.99500,-0.1406253) (4.00000,-0.1358462)};\addlegendentry{X}
\end{axis}\end{tikzpicture}\end{center}

角度未指定使任意斜传播并非唯一曲线；这些明确标注的平行／垂直支覆盖可核查的极限，斜传播须用章\ref{ch:cold}的一般矩阵另给$\theta$。
\end{solution}

\begin{exercise}\label{ex:hist05-9-3}
2005 PS9 原题3。在一维边缘开放磁力线模型中取传导主导、忽略$nu$项、$T_e=T_i=T$；偏滤器通道长$\ell=10\,\mathrm m$，通道入口$T_{\rm in}=50\,\mathrm{eV}$，输入功率密度$q=500\,\mathrm{MW\,m^{-2}}$。求靶板温度并画沿通道距离$x$的$T(x)$。题卷未附所称课堂一维SOL闭合公式，需明确使用的热导系数与边界方向。
\end{exercise}
\begin{solution}
编者独立解答。令$x=0$为入口、$x=\ell$为靶板，若采用常用电子平行Spitzer--H\"arm闭合$\kappa_\parallel(T)=\kappa_0T^{5/2}$（$T$以eV，$\kappa_0$须取与其一致的单位），稳态无体源的$q=-\kappa_0T^{5/2}\dd T/\dd x$给
\begin{equation}
T(x)^{7/2}=T_{\rm in}^{7/2}-\frac{7q}{2\kappa_0}x,\qquad
T_{\rm plate}=\left(T_{\rm in}^{7/2}-\frac{7q\ell}{2\kappa_0}\right)^{2/7}.
\end{equation}
曲线从$50\,$eV单调下降且向下弯曲，只有$T_{\rm in}^{7/2}\ge 7q\ell/(2\kappa_0)$才存在此无体源、纯传导的全段实解。若不满足，给定入口温度和热流不能由该闭合传至靶板，应补充体源、对流、热流限制或改变入口条件。实际题PDF没有提供$\kappa_0$及被提及的课堂公式，故无法从公开题面单独得唯一可靠数值靶板温度；不能捏造一个“官方数值答案”。
\end{solution}
